Универзитет у Београду

Електротехнички факултет

\includegraphics[width=1.54792in,height=1.80903in,alt={proba}]{media/image1.png}

Калкулатор неких специјалних функција

Мастер рад

{\def\LTcaptype{none} % do not increment counter
\begin{longtable}[]{@{}
  >{\raggedright\arraybackslash}p{(\linewidth - 2\tabcolsep) * \real{0.6152}}
  >{\raggedright\arraybackslash}p{(\linewidth - 2\tabcolsep) * \real{0.3848}}@{}}
\toprule\noalign{}
\endhead
\bottomrule\noalign{}
\endlastfoot
Ментор: & Кандидат: \\
проф. др Братислав Иричанин & Сара Лазић, 3238/2021 \\
\end{longtable}
}

Београд, септембар 2025.

\textbf{Садржај}

\begin{enumerate}
\def\labelenumi{\arabic{enumi}.}
\tightlist
\item
\end{enumerate}

\begin{enumerate}
\def\labelenumi{\arabic{enumi}.}
\setcounter{enumi}{1}
\tightlist
\item
\end{enumerate}

\begin{itemize}
\tightlist
\item
\end{itemize}

\begin{itemize}
\tightlist
\item
\end{itemize}

\begin{itemize}
\tightlist
\item
\end{itemize}

\begin{itemize}
\tightlist
\item
\end{itemize}

\begin{itemize}
\tightlist
\item
\end{itemize}

\begin{itemize}
\tightlist
\item
\end{itemize}

\begin{itemize}
\tightlist
\item
\end{itemize}

\begin{itemize}
\tightlist
\item
\end{itemize}

\begin{itemize}
\tightlist
\item
\end{itemize}

\begin{itemize}
\tightlist
\item
\end{itemize}

\begin{itemize}
\tightlist
\item
\end{itemize}

\begin{itemize}
\tightlist
\item
\end{itemize}

\begin{itemize}
\tightlist
\item
\end{itemize}

\begin{itemize}
\tightlist
\item
\end{itemize}

\begin{itemize}
\tightlist
\item
\end{itemize}

\begin{itemize}
\tightlist
\item
\end{itemize}

\begin{itemize}
\tightlist
\item
  \begin{itemize}
  \tightlist
  \item
  \end{itemize}
\end{itemize}

\begin{itemize}
\tightlist
\item
\end{itemize}

\begin{itemize}
\tightlist
\item
  \begin{itemize}
  \tightlist
  \item
  \end{itemize}
\end{itemize}

\begin{itemize}
\tightlist
\item
\end{itemize}

\begin{itemize}
\tightlist
\item
\end{itemize}

\begin{itemize}
\tightlist
\item
\end{itemize}

\begin{itemize}
\tightlist
\item
\end{itemize}

\begin{itemize}
\tightlist
\item
  \begin{itemize}
  \tightlist
  \item
  \end{itemize}
\end{itemize}

\begin{itemize}
\tightlist
\item
\end{itemize}

\begin{itemize}
\tightlist
\item
\end{itemize}

\begin{itemize}
\tightlist
\item
\end{itemize}

\begin{itemize}
\tightlist
\item
\end{itemize}

\begin{itemize}
\tightlist
\item
\end{itemize}

\begin{itemize}
\tightlist
\item
\end{itemize}

\begin{itemize}
\tightlist
\item
\end{itemize}

\begin{itemize}
\tightlist
\item
\end{itemize}

\begin{itemize}
\tightlist
\item
\end{itemize}

\begin{itemize}
\tightlist
\item
\end{itemize}

\begin{itemize}
\tightlist
\item
\end{itemize}

\begin{itemize}
\tightlist
\item
\end{itemize}

\begin{itemize}
\tightlist
\item
\end{itemize}

\begin{itemize}
\tightlist
\item
\end{itemize}

\begin{itemize}
\tightlist
\item
\end{itemize}

\begin{itemize}
\tightlist
\item
\end{itemize}

\begin{itemize}
\tightlist
\item
\end{itemize}

\begin{itemize}
\tightlist
\item
\end{itemize}

\begin{itemize}
\tightlist
\item
\end{itemize}

\begin{itemize}
\tightlist
\item
\end{itemize}

\begin{itemize}
\tightlist
\item
\end{itemize}

\begin{itemize}
\tightlist
\item
\end{itemize}

\begin{itemize}
\tightlist
\item
\end{itemize}

Увод \hyperref[_Toc]{3}Беселове функције \hyperref[_Toc1]{4}Беселова
диференцијална једначина \hyperref[_Toc2]{4}Беселове функције прве врсте
\hyperref[_Toc3]{6}Особине Беселових функција прве врсте
\hyperref[_Toc4]{7}Беселове функције друге врсте
\hyperref[_Toc5]{15}Особине Беселових функција друге врсте
\hyperref[_Toc6]{16}Беселове функције треће врсте
\hyperref[_Toc7]{18}Особине Беселових функција треће врсте
\hyperref[_Toc8]{19}Везе између Беселових функција прве, друге и треће
врсте \hyperref[_Toc9]{20}Гама функција \hyperref[_Toc10]{21}Дефиниције
гама функције \hyperref[_Toc11]{21}Особине гама функције
\hyperref[_Toc12]{24}Апроксимације гама функције
\hyperref[_Toc13]{25}Бета функција \hyperref[_Toc14]{28}Дефиниције бета
функције \hyperref[_Toc15]{28}Особине бета функције
\hyperref[_Toc16]{31}Ортогонални полиноми
\hyperref[_Toc17]{35}Диференцијална једначина
\hyperref[_Toc18]{36}Лежандрови полиноми \hyperref[_Toc19]{39}Лежандрова
диференцијална једначина \hyperref[_Toc20]{39}Лежандров полином
\hyperref[_Toc21]{41}Лагерови полиноми \hyperref[_Toc22]{44}Лагерова
диференцијална једначина \hyperref[_Toc23]{44}Лагеров полином
\hyperref[_Toc24]{45}Ермитови полиноми \hyperref[_Toc25]{48}Ермитовa
диференцијална једначина \hyperref[_Toc26]{48}Ермитов полином
\hyperref[_Toc27]{49}Чебишевљеви полиноми
\hyperref[_Toc28]{53}Чебишевљева диференцијална једначина
\hyperref[_Toc29]{53}Чебишевљеви полиноми прве врсте
\hyperref[_Toc30]{54}Чебишевљеви полиноми друге врсте
\hyperref[_Toc31]{57}Везе између Чебишевљевих полинома прве и друге
врсте \hyperref[_Toc32]{59}Јакобијеви полиноми
\hyperref[_Toc33]{61}Диференцијална једначина
\hyperref[_Toc34]{61}Јакобијеви полиноми \hyperref[_Toc35]{62}Преглед
основних особина специјалних функција \hyperref[_Toc36]{65}Предлози за
рачунарску имплементацију \hyperref[_Toc37]{69}Рачунарска имплементација
Kалкулатора неких специјалних функција \hyperref[_Toc38]{71}Преглед
апликације \hyperref[_Toc39]{71}Компонента за евалуацију израза
\hyperref[_Toc40]{75}Графички приказ функције
\hyperref[_Toc41]{79}Израчунавање вредности специјалне функције
\hyperref[_Toc42]{80}Беселова функција прве врсте
\hyperref[_Toc43]{81}Гама функција \hyperref[_Toc44]{84}Бета функција
\hyperref[_Toc45]{87}Лежандров полином \hyperref[_Toc46]{88}Лагеров
полином \hyperref[_Toc47]{89}Чебишевљеви полиноми
\hyperref[_Toc48]{91}Јакобијеви полиноми \hyperref[_Toc49]{92}Ограничења
апликације \hyperref[_Toc50]{93}Поређење са постојећим програмским
пакетима исте намене \hyperref[_Toc51]{95}Закључак \hyperref[_Toc52]{99}

ЛИТЕРАТУРА \hyperref[_Toc53]{100}

Списак слика \hyperref[_Toc54]{104}

Списак табела \hyperref[_Toc55]{105}

БИОГРАФИЈА \hyperref[_Toc56]{106}

\protect\phantomsection\label{_Toc}{}Увод~

Предмет овог рада је рачунарска имплементација калкулатора неких
специјалних математичких функција.

Специјалне функције су посебна врста математичких функција које своју
примену налазе најчешће приликом решавања одређених проблема у физици и
математици, као и у техничким наукама. Због своје важности и
учесталости, ове функције добиле су посебна имена и устаљене
нотације.~Оно што је интересантно у вези ових функција јесте да су за
већину њих егзактне вредности познате само за неке вредности аргумента,
док се за све остале тачке њихове вредности рачунају нумеричким или
неким другим апроксимацијама {[}44{]}.~Готово све специјалне функције
могу бити записане на један од два начина (што се може и доказати), па
се на основу тога могу разврстати у две групе:~

Функције у облику интеграла

Функције у облику бесконачних конвергентних (најчешће степених) редова.~

У функције у облику интеграла спадају \emph{бета} и \emph{гама}
функције, које су, између осталих својих примена, битан фактор у
анализирању других специјалних функција. \emph{Бета} и \emph{гама}
функције познате су и под именима Ојлеров интергал прве, односно друге
врсте, респективно {[}16{]}.~

Услед њихове веома распрострањене примене у инжењерству, јавља се
потреба за софтвером који ће на ефикасан начин апроксимирати вредности
ових функција - на пример интуитивном веб апликацијом која ће да
обезбеди апроксимацију у минималном броју корака, да дâ графички приказ
функције, као и да пружи неке основне информације о функцији. Ово и
јесте основна идеја овог завршног рада. Истоимени рад ''Калкулатор неких
специјалних функција'' представљен је на XV Српском математичком
конгресу (15. СМАК, 2024) {[}24{]}. Додатно, рад ``Калкулатор неких
специјалних математичких функција'' објављен је у часопису ``Serbian
Journal of Electrical Engineering'' {[}25{]}.

Рад се састоји из четири велике целине. Прву чини теоријски увод у
специјалне функције које су обухваћене у каснијој рачунарској
имплементацији - дефиније, графици, најважније особине, рекурентне везе,
итд. Друга целина представља идеје за рачунарску имплементацију, док је
трећа целина сама имплементација калкулатора за неке специјалне функције
--- коришћене технологије, библиотеке, алгоритми за сваку функцију,
примећени проблеми при имплементацији, ограничења апликације. Коначно,
четврта целина садржи поређење са софтверским алатима који имају
приближне функционалности.

Специјалне функције које су обрађене у овом раду су:

\begin{itemize}
\item
  Беселове функције,
\item
  гама функција,
\item
  бета функција,
\item
  Лежандрови полиноми,
\item
  Лагерови полиноми,
\item
  Ермитови полиноми,
\item
  Чебишевљеви полиноми,
\item
  Јакобијеви полиноми.

  \protect\phantomsection\label{_Toc1}{}Беселове функције
\end{itemize}

О значају Беселових функција можда најбоље говори податак да је о њима
објављен највећи број научних радова од свих специјалних функција. Име
носе по немачком математичару Фридриху Беселу који их је систематизовао,
а ове функције су пре њега проучавали чланови породице Бернули. У
радовима их први спомињу браћа Јохан и Јакоб крајем седамнаестог века, а
потом их дефинише Јакобов син Данијел Бернули {[}44{]}.~

Као и највећи број специјалних функција, прво су се појавиле при
истраживањима из области физике, али оне данас играју битну улогу и у
атомској, нуклеарној и радио физици, електротехници, механици флуида и
хидродинамици,~те у металургији где се појављују нпр. у анализама
преноса топлоте {[}3{]}.

Да бисмо дефинисали Беселове функције, пре свега ћемо се осврнути на
Беселову диференцијалну једначину, а потом дискутовати и једначине прве,
друге и треће врсте - њихове дефиниције и најбитније особине.

\begin{itemize}
\item
  \protect\phantomsection\label{_Toc2}{} Беселова диференцијална
  једначина
\end{itemize}

\emph{\textbf{Дефиниција 1}}: Беселова једначина је хомогена линеарна
диференцијална једначина другог реда чији је општи облик:

\(x^{2}y'' + xy' + \left( x^{2} - \alpha^{2} \right)y = 0\), (1)

где је \(y = f(x)\), \(\alpha \in {\mathbb{C}}\), {[}22{]}.

Дељењем целе једначине са \(x^{2}\) добија се:

\begin{quote}
\(y'' + \frac{1}{x}y' + \left( 1 - \frac{\alpha^{2}}{x^{2}} \right)y = 0\),
\end{quote}

одакле се јасно види да је ово врста Рикатијеве једначине, јер се
редуковањем нелинеарне канонске Рикатијеве диференцијалне једначине
првог реда добија управо једначина облика \(y'' - R(x)y' + S(x)y = 0\).

Дакле, Беселове функције су у тесној вези с поменутом диференцијалном
једначином првог реда, познатом под називом Рикатијева једначина, и то
на начин да се Беселова функција природно дефинише као партикуларно
решење линеарне диференцијалне једначине другог реда (Беселове
једначине), која се добија управо из познате Рикатијеве једначине
применом елементарних трансформација. Наравно, ово није и једини начин
да се дефинишу Беселове функције, о чему ће више речи бити у наставку.~~

Решавањем Беселове једначине добијају се сингуларна решења у 0 и
\(\infty\), као и партикуларна решења која су функције познате као
Беселове функције прве врсте \(J_{\alpha}(x)\) и Беселове функције друге
врсте \(Y_{\alpha}(x)\).

При решавању линеарних диференцијалних једначина другог реда
\(z^{2}u'' + p(z)zu' + q(z)u = 0\) често се користи тзв. Фробениусов
метод: уколико функције \(p(z)\) и \(q(z)\) имају извод у тачки
\(z = 0\) или некој њеној околини, и њихове граничне вредности у тачки
нула постоје (и нису бесконачне), онда постоји решење

\(u(z) = \overset{\infty}{\sum_{k = 0}}A_{k}z^{k + r}\)

које задовољава диференцијалну једначину другог реда {[}9{]}.

Будући да Беселова једначина (1) испуњава наведене услове, на њу може да
се примени Фробениусов метод. Претпоставимо да решење има облик:

\(y(x) =\)\(\overset{\infty}{\sum_{m = 0}}a_{m}x^{m + s}\), (3)

где је \(s\) фиксни параметар, а \(a_{m}\) неодређени коефицијенти. У
том случају су први и други извод функције (3):

\(y'(x)\)
=\(\overset{\infty}{\sum_{m = 0}}a_{m}\overset{}{(m + s)}x^{m + s - 1}\),
(4)

\(y''(x)\)\(= \overset{\infty}{\sum_{m = 0}}a_{m}\overset{}{(m + s)(m + s - 1)}x^{m + s - 2}\).
(5)

Заменом (4) и (5) у Беселову једначину (1), добијамо:

\[\overset{\infty}{\underset{m = 0}{x^{2}\sum}}a_{m}\overset{}{(m + s)(m + s - 1)}x^{m + s - 2} + \overset{\infty}{\underset{m = 0}{x\sum}}a_{m}\overset{}{(m + s)}x^{m + s - 1} + \overset{}{(x^{2} - \alpha^{2})}\overset{\infty}{\sum_{m = 0}}a_{m}x^{m + s}\]

\(=\)\(\overset{\infty}{\sum_{m = 0}}a_{m}\overset{}{(m + s)(m + s - 1)}x^{m + s} + \overset{\infty}{\sum_{m = 0}}a_{m}\overset{}{(m + s)}x^{m + s} + \overset{}{(x^{2} - \alpha^{2})}\overset{\infty}{\sum_{m = 0}}a_{m}x^{m + s}\)

\(=\)\(\overset{\infty}{\sum_{m = 0}}a_{m}\overset{}{((m + s)(m + s - 1) + (m + s) - \alpha^{2})}x^{m + s} + \overset{\infty}{\sum_{m = 0}}a_{m}x^{m + s + 2}\)

\(=\)\(\overset{\infty}{\sum_{m = 0}}a_{m}\overset{}{({(m + s)}^{2} - \alpha^{2})}x^{m + s} + \overset{\infty}{\sum_{m = 2}}a_{m - 2}x^{m + s} = 0\).

Пошто сваки коефицијент испред степена \(x\) мора бити једнак нули,
добијамо следећи систем једначина:

\(m = 0 \Rightarrow a_{0}(s^{2} - \alpha^{2}) = 0\), (6)

\(m = 1 \Rightarrow a_{1}({(s + 1)}^{2} - \alpha^{2}) = 0\), (7)
\(m \geq 2 \Rightarrow a_{m}({(s + m)}^{2} - \alpha^{2}) + a_{m - 2} = 0\).
(8)

Из једначине (6), будући да је \(a_{0} \neq 0\), можемо закључити да је
\(s = \pm \alpha\).

Узмимо прво у разматрање да је \(s = \alpha\). Тада једначине (7) и (8)
добијају облик:

\(a_{1}(2\alpha + 1) = 0\),

\(a_{m}m(m + 2\alpha) + a_{m - 2} = 0\),

редом.

У случају да је \(\alpha\) ненегативан цео број, тада имамо да је

\(a_{1} = 0\),

\(a_{m} = - \frac{a_{m - 2}}{m(2\alpha + m)},m \geq 2\) .

Ово би значило да су сви непарни коефицијенти (када је \(m\) непарно)
једнаки нули, а да се парни одређују по формули:

\(a_{2m} = - \frac{a_{2m - 2}}{4m(\alpha + m)},m \in {\mathbb{N}}\).

Из тога произилази:

\[a_{2} = - \frac{a_{0}}{4(\alpha + 1)} \Rightarrow a_{4} = \frac{a_{0}}{2^{4}2(\alpha + 1)(\alpha + 2)} = \frac{a_{0}}{4^{2}2!(\alpha + 1)(\alpha + 2)}\]

\[\Rightarrow a_{6} = - \frac{a_{4}}{4 \cdot 3(\alpha + 3)} = - \frac{a_{0}}{4^{3}3!(\alpha + 1)(\alpha + 2)(\alpha + 3)}\]

\(\Rightarrow ... \Rightarrow a_{2m} = {( - 1)}^{m}\frac{a_{0}}{4^{m}m!\overset{m}{\prod_{k = 1}}(\alpha + k)},\forall m \in {\mathbb{N}}\).

Будући да \(a_{0}\) можемо сами да бирамо, изабраћемо да буде:

\(a_{0} = \frac{1}{2^{\alpha}\Gamma(\alpha + 1)}\),

при чему је \(\Gamma\) гама функција која се дефинише као
\(\Gamma(x)\)\(= \int_{0}^{\infty}e^{- t}t^{x - 1}dt\) за комплексан
број \(x\) чији је реални део позитиван. Гама функција је генерализација
факторијела на скуп реалних бројева {[}15{]}.

Сада коначно можемо дати израз као решење једначине (1).

\begin{itemize}
\item
  \protect\phantomsection\label{_Toc3}{} Беселове функције прве врсте
\end{itemize}

\emph{\textbf{Дефиниција 2:}} Нека је \(\alpha\) комплексна константа
таква да \(2\alpha\) није негативан цео број. Тада се функција

\(J_{\alpha}(x) =\)\(\overset{\infty}{\sum_{m = 0}}\frac{{( - 1)}^{m}\left( \frac{x}{2} \right)^{\alpha + 2m}}{m!\Gamma(\alpha + m + 1)}\)

назива \emph{Беселова функција прве врсте, реда} \(\alpha\)\emph{, са
аргументом} \(x\)\emph{,} \(x \in {\mathbb{R}}\)\emph{.}

Тако, на пример, за \(\alpha = 0\) и за \(\alpha = 1\) разликујемо

\begin{itemize}
\item
  Беселову функцију нултог реда прве врсте
\end{itemize}

\(J_{0}(x)\)\(= \overset{\infty}{\sum_{m = 0}}\frac{{( - 1)}^{m}\left( \frac{x}{2} \right)^{2m}}{m!\Gamma(m + 1)}\)\(= 1 - \frac{x^{2}}{2^{2}} + \frac{x^{4}}{2^{2}4^{2}} - \frac{x^{6}}{2^{2}4^{2}6^{2}} + ...\)

\begin{itemize}
\item
  Беселову функцију првог реда прве врсте
\end{itemize}

\[J_{1}(x)\]\[= \overset{\infty}{\sum_{m = 0}}\frac{{( - 1)}^{m}\left( \frac{x}{2} \right)^{2m + 1}}{m!\Gamma(m + 2)} = \frac{x}{2} - \frac{x^{3}}{2^{2} \cdot 2} + \frac{x^{5}}{2^{5} \cdot 3!} - \frac{x^{7}}{2^{7}4!} + ...\]

Аналогно је и када се ради са \(- \alpha\), односно када \(2\alpha\)
није позитиван цео број:

\(J_{- \alpha}(x) =\)\(\overset{\infty}{\sum_{m = 0}}\frac{{( - 1)}^{m}\left( \frac{x}{2} \right)^{- \alpha + 2m}}{m!\Gamma( - \alpha + m + 1)}\).

Осим дате дефиниције, постоје и алтернативне дефиниције ових функција
које ће бити дате у наставку.

\emph{\textbf{Дефиниција 3:}} Беселове функције прве врсте дефинишу се
изразом:

\[J_{n}(x) =\]\[\overset{\infty}{\sum_{m = 0}}\frac{{( - 1)}^{m}}{(m + n)!m!}\left( \frac{x}{2} \right)^{2m + n}\]

за свако целобројно \(n\), {[}22{]}\emph{.}

\emph{\textbf{Дефиниција 4:}} Беселова функција за целобројне вредности
реда може да се изрази у облику интеграла {[}49{]}:

\(J_{n}(x) =\)\(\frac{1}{\pi}\int_{0}^{\pi}\cos\)\((n\tau - x\sin\tau)d\tau\),
(9)

односно:

\(J_{n}(x) =\)\(\frac{2}{\pi}\int_{0}^{\frac{\pi}{2}}\)\(\cos(n\tau)\cos(x\sin\tau)d\tau,\)
за свако парно \(n\),

\(J_{n}(x) =\)\(\frac{2}{\pi}\int_{0}^{\frac{\pi}{2}}\)\(\sin(n\tau)\cos(x\sin\tau)d\tau,\)за
свако непарно \(n\).

Одавде такође може проистећи и облик:

\(J_{\alpha}(x) =\)\(\frac{1}{2\pi}\int_{- \pi}^{\pi}e^{- i(n\tau - x\sin\tau)}d\tau\).

\emph{\textbf{Дефиниција 5}:} Беселова функција за реалне вредноси реда
може да се добије уопштавањем израза (9) додавањем новог члана,
{[}36{]}:

\(J_{\alpha}(x) =\)\(\frac{1}{\pi}\int_{0}^{\pi}\cos\)\((\alpha\tau - x\sin\tau)d\tau\)\(- \frac{\sin(\alpha\pi)}{\pi}\int_{0}^{\infty}e^{- xsht - \alpha t}dt\).

\begin{itemize}
\item
  \protect\phantomsection\label{_Toc4}{}Особине Беселових функција прве
  врсте
\end{itemize}

Из дефиниције 2. се види да је у случају парности броја \(\alpha\)
Беселова функција парна, а у случају њене непарности и Беселова функција
је непарна.

На наредној слици 2.3.1. може се видети график Беселове функције прве
врсте за различите вредности реда
\(\alpha\).\includegraphics[width=5.56383in,height=4.32077in,alt={pasted-movie.png}]{media/image2.png}

Слика 2.3.1: График Беселових функција прве врсте за целобројни ред,
\(n = 0(1)4\).

Беселова функција прве врсте дефинисана је за све реалне вредности
независно променљиве.

Графици Беселових функција прве врсте подсећају на синусоиде које
опадају како се променљива приближава бесконачности (било позитивној,
било негативној), и то интензитетом пропорционалним
\(\frac{1}{\sqrt{x}}\), иако њихова решења, у принципу, нису периодична
осим асимптотски за велике вредности променљиве \(x\). Што је ред
функије већи, график мање подсећа на синусоиду.

За позитивне вредности променљиве, Беселова функција целобројног реда
\(J_{n}(x)\) достиже свој први локални максимум када је
\(x \approx \frac{n\pi}{2}\). Ово правило не важи за веће вредности
реда, а када ред узима јако велике вредности максимум узима вредности
приближније \(n\pi\).

Са слике 2.3.1. могу се уочити још неке правилности; једна од њих је да
су нуле функције првог реда управо вредности у којима функција нултог
реда достиже своје локалне минимуме и максимуме. Такође, за вредност
\(x\) за које \(J_{1}(x)\) достиже локални максимум графици \(J_{0}(x)\)
и \(J_{2}(x)\) се секу, а то се понавља и за наредне вредности \(n\).
Ово понашање може да се уопшти: за \(n \geq 1\), променљива за коју
\(J_{n}(x)\) достиже екстремум одговара тачки пресека функција
\(J_{n - 1}(x)\) и \(J_{n + 1}(x)\), {[}35{]}.

Када је \(\alpha\) цео број, тада за функције реда \(\alpha\) и
\(- \alpha\) важи релација:

\(J_{- \alpha}(x) = {( - 1)}^{\alpha}J_{\alpha}(x)\).

Тада је друго решење једначине Беселова функција друге врсте, о чему ће
више речи бити у наредном поглављу.

Када \(\alpha\) није цео број, онда су функције прве врсте реда
\(\alpha\) и \(- \alpha\) независне, и стога представљају два решења
Беселове диференцијалне једначине.

Да би се сагледале још неке важне особине Беселових функција прве врсте,
треба посматрати генеративну функцију\footnote{Генеративна функција за
  низ \(\{ a_{n}\}\) је
  \(g(t) =\)\(\overset{\infty}{\sum_{n = 0}}a_{n}t^{n}\), аналогно томе
  се дефинише и генеративна функција за

  неки низ функцију \(\{ f_{n}(x)\}\):
  \(g(x,t) =\)\(\overset{\infty}{\sum_{n = 0}}f_{n}\)\((x)t^{n}\).
  Понекад се за њу користи израз \emph{функција генератрисе}.} Беселових
једначина прве врсте, целобројног реда:

\(g(x,t) =\)\(\overset{\infty}{\sum_{n = - \infty}}\)\(J_{n}(x)t^{n}\),

\(g(x,t) =\)\(\overset{}{\underset{}{\overset{\infty}{\sum_{n = - \infty}}\overset{\infty}{\sum_{m = 0}}\frac{{( - 1)}^{n}\left( \frac{x}{2} \right)^{n + 2m}}{m!\Gamma(n + m + 1)}}}t^{n}\).

За извођење особина Беселових функција првог реда, као и доказивање
рекурентних веза које важе између њих корисна је чињеница да се ова
функција може записати и као

\(e^{\frac{x}{2}(t - \frac{1}{t})}\). Најпре ће бити изложен доказ ове
тврдње.

\emph{\textbf{Лема 1}}: Функција \(e^{\frac{x}{2}(t - \frac{1}{t})}\) је
генеративна функција Беселових функција првог реда.

\emph{Доказ}: За доказивање ове леме корисна је чињеница да се
експоненцијална функција може представити Маклореновим развојем:

\(e^{x} = \overset{\infty}{\sum_{l = 0}}\frac{x^{l}}{l!}\).

Сада функцију за коју доказујемо тврђење можемо написати на следећи
начин:

\(e^{\frac{x}{2}(t - \frac{1}{t})} = e^{\frac{x}{2}t}e^{- \frac{x}{2}\frac{1}{t}} = \overset{\infty}{\sum_{k = 0}}\frac{\left( \frac{xt}{2} \right)^{k}}{k!}\overset{\infty}{\sum_{l = 0}}\frac{{( - 1)}^{l}\left( \frac{x}{2t} \right)^{l}}{l!}\)\(= \overset{\infty}{\sum_{k = 0}}\frac{\left( \frac{x}{2} \right)^{k}}{k!}t^{k}\overset{\infty}{\sum_{l = 0}}\frac{{( - 1)}^{l}\left( \frac{x}{2} \right)^{l}}{l!}t^{- l}\).

Применом Кошијевог производа\footnote{Ако два бесконачна степена реда
  \(f(x) =\)\(\overset{\infty}{\underset{n = 0}{\sum}}a_{n}x^{n}\) и
  \(g(x) =\)\(\overset{\infty}{\underset{n = 0}{\sum}}b_{n}x^{n}\) имају
  исти полупречник конвергенције \(R > 0\), тада њихов производ може
  такође бити представљен као степени ред у кругу \(|x| < R\):
  \(R(x) \circ g(x) =\)\(\overset{\infty}{\underset{n = 0}{\sum}}c_{n}x^{n}\),
  при чему су коефицијенти
  \(c_{n} = \overset{n}{\underset{k = 0}{\sum}}a_{k}b_{n - k}\).} на два
бесконачна степена реда даље добијамо:

\[\overset{\infty}{\sum_{k = 0}}\frac{\left( \frac{x}{2} \right)^{k}}{k!}t^{k}\overset{\infty}{\sum_{l = 0}}\frac{{( - 1)}^{l}\left( \frac{x}{2} \right)^{l}}{l!}t^{- l} =\]\[\overset{\infty}{\sum_{n = - \infty}}(\overset{\infty}{\sum_{k - l = n,k,l > 0}}\frac{{( - 1)}^{l}\left( \frac{x}{2} \right)^{k + l}}{k!l!})t^{k - l}\]

\(= \overset{\infty}{\sum_{n = - \infty}}(\sum_{k - l = n}\frac{\left( \frac{x}{2} \right)^{k}}{k!}\frac{{( - 1)}^{l}\left( \frac{x}{2} \right)^{l}}{l!})t^{n}\)
\(= \overset{\infty}{\sum_{n = - \infty}}(\overset{\infty}{\sum_{k = 0}}\frac{\left( \frac{x}{2} \right)^{k}}{k!}\frac{{( - 1)}^{k - n}\left( \frac{x}{2} \right)^{k - n}}{(k - n)!})t^{n}\)

\(= \overset{\infty}{\sum_{n = - \infty}}\)\(J_{n}(x)t^{n}\).

\(\square\)

Ова лема ће даље бити коришћена за доказивање других тврђења.

Пре свега, биће изложен доказ за једну од особина већ наведених у овом
поглављу.

\emph{\textbf{Теорема 1:}} За \(\forall n \in {\mathbb{Z}}\) важи
особина:

\(J_{n}( - x) = J_{- n}(x) = {( - 1)}^{n}J_{n}(x)\).

\emph{Доказ:} За доказивање ове теореме неопходна нам је генеративна
функција, у којој ћемо за променљиву узети вредност \(- x\), а уместо
\(t\) користићемо \(t^{- 1}\). Значи, имамо:

\(g( - x,t^{- 1}) =\)\(\overset{\infty}{\sum_{n = - \infty}}\)\(J_{n}( - x)t^{- n}\)
\(= e^{\frac{- x}{2}(\frac{1}{t} - t)}\)
\(= e^{\frac{- x}{2t}}e^{\frac{xt}{2}}\)\emph{\textbf{.}}

Када развијемо експоненцијалне функције у Маклоренов ред, добија се:

\[g( - x,t^{- 1}) =\]\[\overset{\infty}{\sum_{m = 0}}\frac{\left( \frac{- x}{2t} \right)^{m}}{m!}\overset{\infty}{\underset{k = 0}{\sum}}\frac{\left( \frac{xt}{2} \right)^{k}}{k!}\]

\(= \overset{\infty}{\sum_{m = 0}}\frac{\left( \frac{x}{2} \right)^{m}{( - 1)}^{m}}{m!}t^{- m}\overset{\infty}{\underset{k = 0}{\sum}}\frac{\left( \frac{x}{2} \right)^{k}}{k!}t^{k}\)

\(= \overset{\infty}{\sum_{n = - \infty}}(\underset{\underset{m,k \geq 0}{m - k = n,}}{\sum}\frac{{( - 1)}^{m}\left( \frac{x}{2} \right)^{m + k}}{m!k!})t^{- n}\)

\(= \overset{\infty}{\sum_{n = - \infty}}(\overset{\infty}{\sum_{m = 0}}\frac{{( - 1)}^{m}}{m!(m - n)!}\left( \frac{x}{2} \right)^{2m}\left( \frac{x}{2} \right)^{- n})t^{- n}\)

\(= \overset{\infty}{\sum_{n = - \infty}}\)\(J_{- n}(x)\)\(t^{- n}\).

Дакле,

\(\overset{\infty}{\sum_{n = - \infty}}\)\(J_{n}( - x)\)\(t^{- n} =\)\(\overset{\infty}{\sum_{n = - \infty}}\)\(J_{- n}(x)\)\(t^{- n}\).

Поређењем коефицијената, долазимо до:

\(J_{n}( - x) = J_{- n}(x)\).

Овиме је један део теореме доказан.

Да бисмо доказали остатак, користићемо дефиницију Беселове функције
исказану преко гама функције {[}50{]}:

\(J_{\alpha}(x) =\)\(\overset{\infty}{\sum_{m = 0}}\frac{{( - 1)}^{m}\left( \frac{x}{2} \right)^{\alpha + 2m}}{m!\Gamma(\alpha + m + 1)}\).

Посматрамо \(J_{- n}(x)\).

\(J_{- n}(x) =\)\(\overset{\infty}{\sum_{m = 0}}\frac{{( - 1)}^{m}}{m!\Gamma( - n + m + 1)}\)\(\left( \frac{x}{2} \right)^{- n + 2m}\).

Подсећања ради, \(\Gamma(n) = (n - 1)!\), па је

\(\Gamma( - n + m + 1) = (m - n)!\).

Због овога је погодно увести смену \(k = m - n\), и тиме се долази до:

\(J_{- n}(x) =\)\(\overset{\infty}{\sum_{k = - n}}\frac{{( - 1)}^{k + n}}{(k + n)!k!}\left( \frac{x}{2} \right)^{2k + n}\)\(= {( - 1)}^{n}\)\(\overset{\infty}{\sum_{k = - n}}\frac{{( - 1)}^{k}}{(k + n)!k!}\left( \frac{x}{2} \right)^{2k + n}\).

Приметимо сада да овај израз подсећа на израз из дефиниције 3. Разлика
је у бројачу који у овом изразу узима почетну вредност \(- n\). Дакле,
овај израз је једнак:

\(J_{- n}(x) =\)\({( - 1)}^{n}\)\(\overset{- 1}{\sum_{k = - n}}\frac{{( - 1)}^{k}}{(k + n)!k!}\left( \frac{x}{2} \right)^{2k + n}\)\(+ {( - 1)}^{n}\)\(\overset{\infty}{\sum_{k = 0}}\frac{{( - 1)}^{k}}{(k + n)!k!}\left( \frac{x}{2} \right)^{2k + n}\).

Лева сума овог израза је једнака нули, јер за негативан број факторијел
је једнак \(\infty\)\footnote{Факторијел за негативан аргумент није
  дефинисан и формално се третира као бесконачан, чиме члан таквог реда
  не доприноси збиру.}, па је самим тим и \(1/( - n)! = 0\) за сваки
природан број \(n\), па израз постаје:

\(J_{- n}(x) =\)\({( - 1)}^{n}\)\(\overset{\infty}{\sum_{k = 0}}\frac{{( - 1)}^{k}}{(k + n)!k!}\left( \frac{x}{2} \right)^{2k + n}\)\({= ( - 1)}^{n}J_{n}(x)\).
\(\square\)

\emph{\textbf{Теорема 2}} {[}2{]}\textbf{:} Нека је \(J_{n}(x)\)
Беселова функција прве врсте реда \(n\). Тада важе следеће једнакости:

\(\cos x = J_{0}(x) +\)\(2\overset{\infty}{\sum_{n = 1}}{( - 1)}^{n}\)\(J_{2n}(x)\),
(10)

\(\sin x =\)\(2\overset{\infty}{\sum_{n = 0}}{( - 1)}^{n}\)\(J_{2n + 1}(x)\),
(11)

\(1 = J_{0}(x) +\)\(2\overset{\infty}{\sum_{n = 1}}\)\(J_{2n}(x)\). (12)

\emph{Доказ:} За тврђења (10), (11) и (12) погодно је користити Ојлерову
формулу:

\(e^{i\theta} = \cos\theta + i\sin\theta\). (13)

Осим ове формуле важна је и формула:

\(\overset{}{\sin(\arg z)} = \frac{1}{2i}\left( z - \frac{1}{z} \right)\).
(14)

Нека је \(z = e^{i\varphi}\); из (14) је
\(\sin\varphi = \frac{1}{2i}\left( z - \frac{1}{z} \right)\). Када се
замени \(x\sin\varphi\) у формулу (13), добија се:

\(е^{ix\sin\varphi} = \cos\)\((x\sin\varphi)\)\(+ i\sin(x\sin\varphi)\).

Даље важи:

\(\cos(x\sin\varphi) + i\sin(x\sin\varphi) =\)\(e^{\frac{x}{2}(z - \frac{1}{z})}\)
\(= \overset{\infty}{\sum_{n = - \infty}}\)\(J_{n}(x)e^{in\varphi}\).

Уколико се сада поново примени Ојлерова формула на десну страну
претходне једнакости, добија се:

\(\cos(x\sin\phi) + i\sin(x\sin\phi) =\)\(\overset{\infty}{\sum_{n = - \infty}}\)\(J_{n}(x)(\cos n\varphi + i\sin n\varphi)\).
(15)

Раздвајањем реалног и имагинарног дела у једнакости (15) добија се:

\(\cos(x\sin\varphi) =\)\(\overset{\infty}{\sum_{n = - \infty}}\)\(J_{n}(x)\cos n\varphi\),
(16)

\(\sin(x\sin\varphi) =\)\(\overset{\infty}{\sum_{n = - \infty}}\)\(J_{n}(x)\sin n\varphi\).
(17) Сада заменом вредности \(\varphi = \frac{\pi}{2}\) у једначину (16)
имамо:

\(\cos x = \overset{\infty}{\sum_{n = - \infty}}\)\(J_{n}(x)\)\(\cos\frac{n\pi}{2}\).
(18)

Разматрамо следеће случаје:

\textbf{\emph{случај 1}:} Нека је \(n \equiv 1(mod4)\). Тада је
\(\cos(\frac{n\pi}{2}) = \cos(\frac{- n\pi}{2}) = 0\), што би значило:

\(J_{n}(x)\)\(\cos(\frac{n\pi}{2}) +\)\(J_{- n}(x)\)\(\cos(\frac{- n\pi}{2}) = 0\).

То је случај и за било које непарно \(n\), па самим тим и
\(n \equiv 3(mod4)\).

\textbf{\emph{случај 2}:} Нека је \(n \equiv 2(mod4)\). Тада је
\(\cos(\frac{n\pi}{2}) = \cos(\frac{- n\pi}{2}) = - 1\). Већ је наведено
да је једна од особина Беселових функција целобројног реда
\(J_{- \alpha}(x) = {( - 1)}^{\alpha}J_{\alpha}(x)\), па када то
заменимо у претходну једнакост долазимо до:

\(J_{n}(x)\cos\)\((\frac{n\pi}{2}) + J_{- n}\)\((x)\cos\)\((\frac{- n\pi}{2}) =\)\(J_{n}(x)( - 1) + J_{n}(x)( - 1)\)\(= - 2J_{n}(x)\).

\textbf{\emph{случај 3}:} Нека је \(n \equiv 0(mod4)\). Тада је
\(\cos(\frac{n\pi}{2}) = \cos(\frac{- n\pi}{2}) = 1\). Када то заменимо
у претходну једнакост долазимо до:

\(J_{n}(x)\cos\)\((\frac{n\pi}{2}) + J_{- n}\)\((x)\cos\)\((\frac{- n\pi}{2}) =\)\(J_{n}(x)(1) + J_{n}(x)(1)\)\(= 2J_{n}(x)\).

Када се сагледају сва три случаја, може се доћи до закључка да се
једначина (18) може записати на следећи начин:

\(\cos x = J_{0}\)\((x) + 2\)\(\overset{\infty}{\sum_{n = 1}}\)\({( - 1)}^{n}J_{2n}(x)\),

што је једно од тврђења које је требало доказати, (9).

С друге стране, заменом \(\varphi = \frac{\pi}{2}\) у једначину (17)
добија се:

\(\sin x = \overset{\infty}{\sum_{n = - \infty}}\)\(J_{n}(x)\)
\(\sin\frac{n\pi}{2}\).

Аналогно као и за једначину (16), и за једначину (17) разликујемо
случаје на основу остатка при дељењу са четири:

\textbf{\emph{случај 1}:} Нека је \(n \equiv 1(mod4)\). Тада је
\(- n \equiv 3(mod4)\), затим \(\sin(\frac{n\pi}{2}) = 1\), и такође
\(\sin(\frac{- n\pi}{2}) = - 1\). Сада збир сума за ред Беселове
једначине \(n\) и \(- n\) постаје:

\(J_{n}(x)\)\(\sin(\frac{n\pi}{2}) +\)\(J_{- n}(x)\)\(\sin(\frac{- n\pi}{2}) =\)\(J_{n}(x)(1) + J_{- n}(x)( - 1)\).

Већ је наведено да је једна од особина Беселових функција са целобројним
редом \(J_{- \alpha}(x) = {( - 1)}^{\alpha}J_{\alpha}(x)\), па када то
заменимо у претходну једнакост долазимо до:

\(J_{n}(x)\)\(\sin(\frac{n\pi}{2}) +\)\(J_{- n}(x)\)\(\sin(\frac{- n\pi}{2}) =\)\(J_{n}(x) +\)\({( - 1)}^{2}\)\(J_{n}(x)\)\(= 2J_{n}(x)\).

\begin{itemize}
\item
  \textbf{\emph{случај 2}:} Нека је \(n \equiv 2(mod4)\). Тада је
  \(\sin(\frac{n\pi}{2}) = 0\), као и \(\sin(\frac{- n\pi}{2}) = 0\), па
  је збир суме једнак 0. До овог резултата се долази и у случају када је
  \(n \equiv 0(mod4)\).
\item
  \textbf{\emph{случај 3}:} Нека је \(n \equiv 3(mod4)\). Тада је
  \(\sin(\frac{n\pi}{2}) = - 1\) и \(\sin(\frac{- n\pi}{2}) = 1\). Сада
  се као резултат сумирања добија:
\end{itemize}

\(J_{n}(x)\)\(\sin(\frac{n\pi}{2}) +\)\(J_{- n}(x)\)\(\sin(\frac{- n\pi}{2}) =\)\(J_{n}(x)( - 1) + J_{- n}(x)(1)\)\(= J_{n}(x)( - 1) + ( - 1)J_{n}(x)\)\(= - 2J_{n}(x)\).

Дакле, када као резултат овог разматрања долази се до закључка да се
може записати:

\(\sin x =\)\(2\overset{\infty}{\sum_{n = 0}}\)\(J_{2n + 1}(x)\)\({( - 1)}^{n}\),

што и јесте друго тврђење које је требало доказати.

Коначно, треће тврђење доказује се на сличан начин као претходна два,
овај пут заменом \(\varphi = 2\pi\) у једначину (16):

\(\cos(x\sin\varphi) = \cos(x\sin 2\pi)\)\(= \cos(0)\)\(= \overset{\infty}{\sum_{n = - \infty}}\)\(J_{n}(x)\sin(2n\pi)\),

\(1 =\)\(\overset{\infty}{\sum_{n = - \infty}}\)\(J_{n}(x)\).
\(\square\)

\emph{\textbf{Теорема 3}} {[}35{]}\emph{\textbf{:}} За Беселове функције
прве врсте чији је ред цео број важи:

\(\frac{d}{dx}\left\lbrack x^{- \alpha}J_{\alpha}(x) \right\rbrack =\)\(- x^{- \alpha}J_{\alpha + 1}(x)\),
као и

\(\frac{d}{dx}\left\lbrack x^{\alpha}J_{\alpha}(x) \right\rbrack =\)\(x^{\alpha}J_{\alpha - 1}(x)\).

\emph{Доказ:}

\[\frac{d}{dx}\left\lbrack x^{- \alpha}J_{\alpha}(x) \right\rbrack =\]\[\frac{d}{dx}\left\lbrack x^{- \alpha}\overset{\infty}{\sum_{m = 0}}\frac{{( - 1)}^{m}}{m!(\alpha + m)!}\left( \frac{x}{2} \right)^{2m + \alpha} \right\rbrack\]

\[= \frac{d}{dx}\left\lbrack \overset{\infty}{\sum_{m = 0}}\frac{{( - 1)}^{m}}{m!(\alpha + m)!}\left( \frac{x}{2} \right)^{2m}\frac{1}{2^{\alpha}} \right\rbrack\]

Диференцијање реда је дозвољено јер ред апсолутно конвергира, што значи
да се збир реда не мења ако се узму апсолутне вредности чланова, пошто
факторијели у имениоцу расту много брже од степена променљиве.

\[2^{- \alpha}\frac{d}{dx}\left\lbrack \overset{\infty}{\sum_{m = 0}}\frac{{( - 1)}^{m}}{m!(\alpha + m)!}\left( \frac{x}{2} \right)^{2m} \right\rbrack\]\[= 2^{- \alpha}\overset{\infty}{\sum_{m = 0}}\frac{{( - 1)}^{m}}{m!(\alpha + m)!}\frac{1}{2^{2m}}2mx^{2m - 1}\]

\[= 2^{- \alpha}\overset{\infty}{\sum_{m = 1}}\frac{{( - 1)}^{m}}{m!(\alpha + m)!}\frac{2m}{2^{2m}}x^{2m - 1}\]

\(= - 2^{- \alpha}\overset{\infty}{\sum_{m = 1}}\frac{{( - 1)}^{m - 1}}{(m - 1)!(\alpha + 1 + m - 1)!}\left( \frac{x}{2} \right)^{2m - 1}\)\emph{.}

Уводећи смену \(k = m - 1\) долазимо до:

\[\frac{d}{dx}\left\lbrack x^{- \alpha}J_{\alpha}(x) \right\rbrack =\]\[- 2^{- \alpha}\overset{\infty}{\sum_{k = 0}}\frac{{( - 1)}^{k}}{k!(\alpha + 1 + k)!}\left( \frac{x}{2} \right)^{2k + 1}\]

\(= - \overset{\infty}{\sum_{k = 0}}\frac{{( - 1)}^{k}}{k!(\alpha + 1 + k)!}\left( \frac{x}{2} \right)^{2k + \alpha + 1}x^{- \alpha}\)
\(= - x^{- \alpha}\overset{\infty}{\sum_{k = 0}}\frac{{( - 1)}^{k}}{k!(\alpha + 1 + k)!}\left( \frac{x}{2} \right)^{2k + \alpha + 1}\)

\(= - x^{- \alpha}J_{\alpha + 1}(x)\).

Тврђење
\(\frac{d}{dx}\left\lbrack x^{\alpha}J_{\alpha}(x) \right\rbrack =\)\(x^{\alpha}J_{\alpha - 1}(x)\)
се доказује аналогно претходном.

\emph{\textbf{Теорема 4:}} Беселова функција прве врсте има следећу
особину : \(J_{0}(0) = 1\).

\emph{Доказ:} Из

\(J_{0}(x) = 1 - \frac{x^{2}}{2^{2}} + \frac{x^{4}}{2^{2}4^{2}} - \frac{x^{6}}{2^{2}4^{2}6^{2}} + \frac{x^{8}}{2^{2}4^{2}6^{2}8^{2}} - ...\)\emph{,}

када заменимо \(x = 0\), добијамо тражену једнакост. \(\square\)

Рекурзивна формула дужине 3 за Беселову функцију прве врсте је:

\(J_{n + 1}\overset{}{(x)} = \frac{2n}{x}J_{n}\overset{}{(x)} - J_{n - 1}\overset{}{(x)},n \in {\mathbb{N}}\).

\begin{itemize}
\item
  \protect\phantomsection\label{_Toc5}{}Беселове функције друге врсте
\end{itemize}

Беселове функције друге врсте, у нотацији \(Y_{n}\)\((x)\) (у неким
случајевима и \(N_{n}\)\((x)\)) се дефинишу као друго, линеарно
независно решење Беселове једначине.

Ова фамилија функција се такође назива и Нојманове функције или Веберове
функције. Нојманова функција има сингуларитет у нули.

\emph{\textbf{Дефиниција 6:}} Када је ред функције нецелобројан,
Беселова функција друге врсте реда \(\alpha\) се изражава као

\(Y_{\alpha}(x) =\)\(\frac{\cos(\alpha\pi)J_{\alpha}(x) - J_{- \alpha}(x)}{\sin(\alpha\pi)}\),
\(\alpha \notin {\mathbb{Z}}\), \(x \in {\mathbb{R}}\).

\emph{\textbf{Дефиниција 7:}} Нека је \(\alpha\) цео број. Тада се
функција:

\(Y_{\alpha}\overset{}{(x)} =\)\(\lim_{\varphi \rightarrow \alpha}Y_{\varphi}\overset{}{(x)} =\)\(\lim_{\varphi \rightarrow \alpha}\frac{\cos(\varphi\pi)J_{\varphi}(x) - J_{- \varphi}(x)}{\sin(\varphi\pi)}\),
\(\alpha \in {\mathbb{Z}}\),

назива Беселова функција друге врсте реда \(\alpha\), чији је аргумент
\(x\).

Беселове функције друге врсте, односно Нојманове функције, имају и своју
интегралну репрезентацију.

\emph{\textbf{Дефиниција 8:}} Беселова функција друге врсте се може
записати на следећи начин:

\(Y_{0}(x) =\)\(- \frac{2}{\pi}\int_{0}^{\infty}\)\(\cos(xcht)dt =\)\(- \frac{2}{\pi}\int_{1}^{\infty}\)\(\frac{\cos(xt)}{{(t^{2} - 1)}^{0.5}}dt,x > 0\),

\(Y_{n}(x) =\)\(- \frac{1}{\pi}\int_{0}^{\pi}\)\(\sin(xsht - nt)dt -\)\(\frac{1}{\pi}\int_{0}^{\infty}\)\((e^{nt} + {( - 1)}^{n}e^{- nt})e^{- xsht}dt\).

Из дефиниција ових функција је јасно да су решење Беселове једначине,
будући да су представљене линеарном комбинацијом Беселових једначина
прве врсте, које су изведене као решење {[}35{]}.

\begin{itemize}
\item
  \protect\phantomsection\label{_Toc6}{}Особине Беселових функција друге
  врсте
\end{itemize}

Нојманова функција је дефинисана за све вредности променљиве осим за
\(x = 0\). Функција је непарна за све вредности променљиве. На слици
2.5.1. се могу видети графици Нојманових фукнција за различите вредности
реда, тј када је
\(n = 0,1,2,3,4\).\includegraphics[width=4.6653in,height=2.74203in,alt={pasted-movie.png}]{media/image3.png}

Слика 2.5.1: Графици Нојманових функција за различите вредности реда.

Попут Беселових функција прве врсте, и Нојманове функције су
осцилаторне. Како се променљива приближава \(+ \infty\), тако функција
почиње да се понаша као линеарна комбинација синусне и косинусне
функције чији интензитет опада.

У поређењу са Беселовом функцијом прве врсте, интентитет Нојманове
функције опада знатно спорије.

За функције чији је ред произвољан комплексан број важе следеће
рекурентне везе {[}35{]}:

\[Y_{\alpha - 1}(x) - Y_{\alpha + 1}(x) = 2Y_{\alpha}(x)\]

\(Y_{\alpha - 1}(x) + Y_{\alpha + 1}(x) =\)\(\frac{2\alpha}{x}\)\(Y_{\alpha}(x)\).

Нека је ред Беселове функције друге врсте цео број
\(n \in {\mathbb{Z}}\):

\[Y_{n}( - x) = - Y_{- n}(x) = {( - 1)}^{n}Y_{n}(x)\]

\[Y_{n - 1}(x) + Y_{n + 1}(x) =\]\[\frac{2n}{x}\]\[Y_{n}(x)\]

\(Y_{n - 1}(x) - Y_{n + 1}(x) = 2Y_{n}'(x)\)

\(xY_{\alpha}'(x) + \alpha Y_{\alpha}(x) = xY_{\alpha - 1}(x)\)

\(xY_{\alpha}'(x) - \alpha Y_{\alpha}(x) = - xY_{\alpha + 1}(x)\)

\[\frac{d}{dx}\]\[(x^{n}\]\[Y_{n}(x)) =\]\[x^{n}Y_{n - 1}(x)\]

\[\int x^{n + 1}Y_{n}\]\[(x)dx =\]\[x^{n + 1}\]\[Y_{n + 1}(x)\]

\[\int x^{- n + 1}Y_{n}\]\[(x)dx = - x^{- n + 1}\]\[Y_{n + 1}(x)\]

Занимљиве су и особине ових функција када је њихов ред
\(\alpha \in \frac{1}{2} + k,k \in {\mathbb{Z}}\) {[}23{]}:

\[J_{\frac{1}{2}}(x) = Y_{- 1/2}(x) =\]\[\sqrt{\frac{2}{\pi x}}\]\[\sin x\]

\[J_{- 1/2}(x) = - Y_{- 1/2}(x) =\]\[\sqrt{\frac{2}{\pi x}}\]\[\cos x\]

\[Y_{k + 1/2}(x) =\]\[{( - 1)}^{k - 1}\]\[J_{- k - 1/2}(x)\]

\(Y_{- k - 1/2}(x)\)\(= {( - 1)}^{k}\)\(J_{k + 1/2}(x)\).

\begin{itemize}
\item
  \protect\phantomsection\label{_Toc7}{}Беселове функције треће врсте
\end{itemize}

Беселове функције треће врсте дефинишемо помоћу Беселових функција прве
и друге врсте. У ову фамилију функција убрајамо:

\begin{itemize}
\item
  Хенкелове функције прве врсте
\item
  Хенкелове функције друге врсте.
\end{itemize}

Дакле, пре свега ћемо дефинисати наведене функције.

\emph{\textbf{Дефиниција 9:}} Нека је \(J_{\alpha}(x)\) Беселова
функција прве врсте, а \(Y_{\alpha}(x)\) Нојманова функција. Тада се
Хенкелове функције прве врсте \emph{нецелобројног} реда дефинишу
изразом:

\(H_{\alpha}^{(1)}(x) = J_{\alpha}(x) + iY_{\alpha}(x)\),
\(x \in {\mathbb{R}}\).

Из овог израза се лако изводи и други облик записа ове функције:

\(H_{\alpha}^{(1)}(x) =\)\(\frac{J_{- \alpha}(x) - e^{- i\alpha\pi}J_{\alpha}(x)}{i\sin(\alpha\pi)}\).

Слично, дефинише се и Хенкелова функција друге врсте.

\emph{\textbf{Дефиниција 10:}} Нека је \(J_{\alpha}(x)\) Беселова
функција прве врсте, а \(Y_{\alpha}(x)\) Нојманова функција. Тада се
Хенкелове функције прве врсте \emph{нецелобројног} реда дефинишу
изразом:

\(H_{\alpha}^{(2)}(x) = J_{\alpha}(x) - iY_{\alpha}(x)\).

Други запис Хенкелове функције друге врсте гласи {[}16{]}:

\(H_{\alpha}^{(2)}(x) =\)\(\frac{J_{- \alpha}(x) - e^{i\alpha\pi}J_{\alpha}(x)}{- i\sin(\alpha\pi)}\).

Свако решење Беселове једначине се може написати као линеарна
комбинација Хенкелове функције прве врсте и Хенкелове функције друге
врсте. Сада ћемо се позабавити доказивањем овог тврђења.

\emph{\textbf{Теорема 5:}} Нека је дата комплексна константа \(\alpha\)
и Ханкелове функције реда \(\alpha\) су линеарно независне. Тада се
опште решење Беселове једначине може представити на следећи начин:

\(\omega(x) = AH_{\alpha}^{(1)}(x) + BH_{\alpha}^{(2)}(x)\),

при чему су \(A\) и \(B\) комплексне константе {[}15{]}.

\emph{Доказ:}

\(\omega(x) = AH_{\alpha}^{(1)}(x) + BH_{\alpha}^{(2)}(x)\)

\(= A\left( J_{\alpha}(x) + iY_{\alpha}(x) \right) + B\left( J_{\alpha}(x) - iY_{\alpha}(x) \right)\)

\[= (a + ib)\left( J_{\alpha}(x) + iY_{\alpha}(x) \right) + (c + id)\left( J_{\alpha}(x) - iY_{\alpha}(x) \right)\]

\[= \left( a + c + i(b + d) \right)\]\[J_{\alpha}(x) +\]\[(ia - b - ic + d)\]\[Y_{\alpha}(x)\]

\(= \left( a + c + i(b + d) \right)\)\(J_{\alpha}(x) +\)\(\left( d - b + i(a - c) \right)\)\(Y_{\alpha}(x)\)\emph{.}

Ако сада уведемо константе \(C = \left( a + c + i(b + d) \right)\)и
\(D = \left( d - b + i(a - c) \right)\), долазимо до:

\(\omega(x) = CJ_{\alpha}(x) + DY_{\alpha}(x)\).

Такође, Вронскијан Ханкелових функција износи:

\(W(H_{\alpha}^{(1)},H_{\alpha}^{(2)}) =\)

\[= (J_{\alpha}(x) + iY_{\alpha}(x))(J_{\alpha}'(x) - iY_{\alpha}'(x)) - (J_{\alpha}'(x) + iY_{\alpha}'(z))(J_{\alpha}(x) - iY_{\alpha}(x))\]\[= - 2i(J_{\alpha}(x)Y_{\alpha}'(x)) - Y_{\alpha}(x)J_{\alpha}'(x))\]

\[= - \frac{4i}{\pi z} ≢ 0\]

Дакле, функција \(\omega(x)\) је линеарна комбинација решења
\(J_{\alpha}(x)\) и \(Y_{\alpha}(x)\) Беселове једначине, па и сама
представља решење Беселове једначине. \(\square\)

\begin{itemize}
\item
  \protect\phantomsection\label{_Toc8}{}Особине Беселових функција треће
  врсте
\end{itemize}

Пошто Беселове функције прве и друге врсте задовољавају исте рекурентне
везе, то значи да ће и Беселове функције трећег реда, као њихова
линеарна комбинација, задовољавати исте те формуле.

У наставку су дате особине Беслових функција трећег реда {[}2{]}.

\(H_{\alpha - 1}^{(1)}\)\((x) +\)\(H_{\alpha + 1}^{(1)}\)\((x) =\)\(\frac{2\alpha}{x}H_{\alpha}^{(1)}\)\((x)\)
\(H_{\alpha - 1}^{(2)}\)\((x) +\)\(H_{\alpha + 1}^{(2)}\)\((x) =\)\(\frac{2\alpha}{x}H_{\alpha}^{(2)}\)\((x)\)

\(H_{\alpha - 1}^{(1)}\)\((x) -\)\(H_{\alpha + 1}^{(1)}\)\((x) = 2\)\(\frac{dH_{\alpha}^{(1)}(x)}{dx}\)
\(H_{\alpha - 1}^{(2)}\)\((x) -\)\(H_{\alpha + 1}^{(2)}\)\((x) = 2\)\(\frac{dH_{\alpha}^{(2)}(x)}{dx}\)

\(x\frac{dH_{\alpha}^{(1)}(x)}{dx} + \alpha H_{\alpha}^{(1)}\)\((x) = x\)\(H_{\alpha - 1}^{(1)}\)\((x)\)
\(x\frac{dH_{\alpha}^{(2)}(x)}{dx} + \alpha\)\(H_{\alpha}^{(2)}(x) = x\)\(H_{\alpha - 1}^{(2)}\)\((x)\)

\(x\frac{dH_{\alpha}^{(1)}(x)}{dx} - \alpha H_{\alpha}^{(1)}\)\((x) = - x\)\(H_{\alpha + 1}^{(1)}\)\((x)\)
\(x\frac{dH_{\alpha}^{(2)}(x)}{dx} - \alpha H_{\alpha}^{(2)}\)\((x) = - x\)\(H_{\alpha + 1}^{(2)}\)\((x)\)

\(\frac{dH_{0}^{(1)}(x)}{dx} = - H_{1}^{(1)}\)\((x)\)
\(\frac{dH_{0}^{(2)}(x)}{dx} = - H_{1}^{(2)}\)\((x)\)

\(\nabla_{\alpha}H_{\alpha}^{(1)}(x) = 0\)
\(\nabla_{\alpha}H_{\alpha}^{(2)}(x) = 0\)

\[\left( \frac{d}{xdx} \right)^{m}\]\[\{ x^{\alpha}H_{\alpha}^{(1)}\]\[(x)\]\[\} = x^{\alpha - m}H_{\alpha - m}^{(1)}\]\[(x)\]

\[\left( \frac{d}{xdx} \right)^{m}\]\[\{ x^{\alpha}H_{\alpha}^{(2)}\]\[(x)\]\[\} = x^{\alpha - m}H_{\alpha - m}^{(2)}\]\[(x)\]

\[\left( \frac{d}{xdx} \right)^{m}\{\frac{H_{\alpha}^{(1)}(x)}{x^{\alpha}}\} =\]\[{( - 1)}^{m}\]\[\frac{H_{\alpha + 1}^{(1)}(x)}{x^{\alpha + m}}\]

\(\left( \frac{d}{xdx} \right)^{m}\{\frac{H_{\alpha}^{(2)}(x)}{x^{\alpha}}\} =\)\({( - 1)}^{m}\)\(\frac{H_{\alpha + 1}^{(2)}(x)}{x^{\alpha + m}}\).

\begin{itemize}
\item
  \protect\phantomsection\label{_Toc9}{}Везе између Беселових функција
  прве, друге и треће врсте
\end{itemize}

Неке од ових веза су већ изложене у претходним поглављима, будући да се
и Нојманове и Хенкелове функције дефинишу помоћу Беселових функција, а у
наставку ће бити дате још неке интересантне везе {[}49{]}.

\[J_{\alpha}(x) =\]\[\frac{H_{\alpha}^{(1)}(x) + H_{\alpha}^{(2)}(x)}{2} = \frac{Y_{- \alpha}(x) - Y_{\alpha}(x)\cos\alpha\pi}{\sin\alpha\pi}\]

\[J_{- \alpha}(x) =\]\[\frac{e^{\alpha\pi i}H_{\alpha}^{(1)}(x) + e^{- \alpha\pi i}H_{\alpha}^{(2)}(x)}{2} = \frac{Y_{- \alpha}(x)\cos\alpha\pi - Y_{\alpha}(x)}{\sin\alpha\pi}\]

\[Y_{\alpha}(x) =\]\[\frac{J_{\alpha}(x)\cos\alpha\pi - J_{- \alpha}(x)}{\sin\alpha\pi} = \frac{H_{\alpha}^{(1)}(x) - H_{\alpha}^{(2)}(x)}{2i}\]

\[Y_{- \alpha}(x) =\]\[\frac{J_{\alpha}(x) - J_{- \alpha}(x)\cos\alpha\pi}{\sin\alpha\pi} = \frac{e^{\alpha\pi i}H_{\alpha}^{(1)}(x) - e^{- \alpha\pi i}H_{\alpha}^{(2)}(x)}{2i}\]

\[H_{\alpha}^{(1)}(x) =\]\[\frac{J_{- \alpha}(x) - e^{- \alpha\pi i}J_{\alpha}(x)}{i\sin\alpha\pi} = \frac{Y_{- \alpha}(x) - e^{- \alpha\pi i}Y_{\alpha}(x)}{\sin\alpha\pi}\]

\(H_{\alpha}^{(2)}(x) =\)\(\frac{e^{\alpha\pi i}J_{\alpha}(x) - J_{- \alpha}(x)}{i\sin\alpha\pi} = \frac{Y_{- \alpha}(x) - e^{\alpha\pi i}Y_{\alpha}(x)}{\sin\alpha\pi}\).

\protect\phantomsection\label{_Toc10}{}Гама функција~

Гама функција је једна од најчешће коришћених специјалних функција.
Примену налази у теорији бројева, вероватноћи, интегралном рачуну, а
посебно у решавању проблема у физици. Често се спомиње и у теорији
других специјалних функција и као таква има посебан значај. Од других
функција обрађених у овом раду издваја се по једноставности својих
рекурентних веза, што је веома битна особина и, може се рећи, управо то
јој и даје ту улогу коју игра у теорији специјалних фунција {[}35{]}.

Ова функција настаје као одговор на проблем проширења факторијела на
скуп реалних бројева, који се први пут јавља при проучавању
интерполације редова. Дефинисао ју је Леонард Ојлер, па је позната и као
\emph{Ојлеров интеграл друге врсте}. Нотацију која је данас устаљена и
њен назив настали су тек касније и уводи их Адријен-Мари Лежандр
{[}44{]}.

\begin{itemize}
\item
  \protect\phantomsection\label{_Toc11}{}Дефиниције гама функције
\end{itemize}

\emph{\textbf{Дефиниција 11:}} Гама функција, на домену
\(x \in {\mathbb{R}},x > 0\), дефинише се Ојлеровим интегралом:

\(\Gamma(x) =\)\(\int_{0}^{\infty}e^{- t}t^{x - 1}dt\).

Ово је најпознатија интегрална репрезентација, али се често користи и
следећа, дата дефиницијом 12.

\emph{\textbf{Дефиниција 12:}} Гама функција, на домену
\(x \in {\mathbb{R}},x > 0\), може се представити интегралом:

\(\Gamma(x) =\)\(2\int_{0}^{\infty}e^{- t}t^{2x - 1}dt\).

Још неки начини да се гама функција запише у форми одређеног интеграла
су и {[}5{]}, {[}46{]}:

\(\Gamma(x) =\)\(2\int_{0}^{1}\left( \ln\frac{1}{t} \right)^{x - 1}dt\),
(19)

\(\Gamma(x + 1) =\)\(\int_{0}^{\infty}t^{x}\left\lbrack e^{t} - 1 + t - \frac{1}{2!}t^{2} + \frac{1}{3!}t^{3} + ... + {( - 1)}^{n + 1}\frac{1}{n!}t^{n} \right\rbrack dt\),

где је \(n\) цео део од \(x\).

Такође, једноставним сменама у израз (19) могу да се добију и облици:

\(\Gamma(x) =\)\(\int_{0}^{\infty}e^{- a^{\frac{1}{x}}}\frac{1}{x}da\),
када се уведе смена \(t^{x} = a\),

\(\Gamma\left( 1 + \frac{1}{x} \right) = \int_{0}^{\infty}e^{- t^{x}}dt\),
када се \(x\) замени са \(\frac{1}{x}\).

Постоји и начин да се ова фунција дефинише коришћењем лимеса, и он је
дат Гаусовом дефиницијом (дефиниција 13), као и коришћењем бесконачног
производа, дефиниција 14.

\emph{\textbf{Дефиниција 13}} {[}6{]}\emph{\textbf{:}} Гама функција, на
домену \(x \in {\mathbb{R}},x > 0\), дефинише се као:

\[\Gamma(x) =\]\[\lim_{n \rightarrow \infty}\frac{n^{x}}{x(1 + x)(1 + \frac{x}{2})...(1 + \frac{x}{n - 1})(1 + \frac{x}{n})}\]

\[= \lim_{n \rightarrow \infty}\frac{n! \cdot n^{x}}{x(1 + x)(2 + x)...(n - 1 + x)(n + x)}\]

\(= \lim_{n \rightarrow \infty}\frac{1}{x}\overset{n}{\prod_{k = 1}}\left( 1 + \frac{x}{k} \right)^{- 1}n^{x}\).

Ово дефиниција је тачна чак и за негативну променљиву, осим за
целобројне вредности. Ако се у ову дефиницију (13) уведе смена
\(n^{- x} = e^{( - \ln n)x}\), добија се:

\(\frac{1}{\Gamma(x)} = x\lim_{n \rightarrow \infty}e^{( - \ln n)x}\overset{n}{\prod_{k = 1}}\left( 1 + \frac{x}{k} \right)\).

\emph{\textbf{Дефиниција 14}} {[}35{]}\emph{\textbf{:}} Гама функција, у
ознаци \(\Gamma(x)\), \(x \in {\mathbb{R}},x > 0\), дата је изразом:

\(\frac{1}{\Gamma(x)} =\)\(xe^{\gamma x}\overset{\infty}{\prod_{j = 0}}\left( 1 + \frac{x}{j} \right) \cdot e^{- \frac{x}{j}}\),

при чему је \(\gamma\) Ојлерова константа која се дефинише као:

\[\gamma = \lim_{n \rightarrow \infty}\left( 1 + \frac{1}{2} + \frac{1}{3} + ... + \frac{1}{n} - \ln n \right)\]

\(= \lim_{n \rightarrow \infty}\left( \overset{n}{\sum_{j = 1}}\frac{1}{j} - \ln n \right) = 0.577215664901533\).

За неке формуле које дефинишу ову функцију погодније је променљиву
записати као \(x + 1\). На пример:

\(\Gamma(1 + x) =\)\(\int_{0}^{1}\ln^{x}\left( \frac{1}{t} \right)dt\),
или

\(\ln\{\Gamma(1 + x)\} =\)\(\int_{0}^{1}\frac{1}{\ln t}(\frac{t^{x} - 1}{t - 1} - x)dt\).

Још једна од таквих формула је и дефиниција 15, код које се јасно види
да је гама функција проширење факторијела изван скупа природних бројева.

\emph{\textbf{Дефиниција 15}} {[}5{]}\emph{\textbf{:}} Гама функција, у
ознаци \(\Gamma(x)\), \(x \in {\mathbb{R}},x > 0\), дата је изразом:

\[\Gamma(1 + x) =\]\[\lim_{n \rightarrow \infty}\frac{1 \cdot 2 \cdot 3 \cdot .... \cdot n}{(x + 1)(x + 2)(x + 3)...(x + n + 1)} \cdot n^{x + 1}\]\[= \lim_{n \rightarrow \infty}\frac{nx}{x + n + 1} \cdot \frac{1 \cdot 2 \cdot 3... \cdot n}{x(x + 1)(x + 2)...(x + n)}n^{x}\]

\(= x \cdot \Gamma(x)\).

Гама функција се такође може записати и као:

\(\Gamma(x + 1) =\)\(\overset{\infty}{\prod_{j = 1}}\frac{{(j + 1)}^{x}}{(j + x) \cdot j^{x - 1}}\).

Помоћу логаритма се гама функција дефинише изразом:

\(\ln(\Gamma(x + 1)) =\)
\(\int_{0}^{1}\left( \frac{t^{x} - 1}{t - 1} - x \right)\frac{dt}{\ln t}\).

Лапласова трансформација би била {[}35{]}:

\(\Gamma(x) =\)\(s^{x}\int_{0}^{\infty}t^{x - 1}e^{- st}dt,x > 0,s > 0\),

а Фуријеова трансформација:

\(\Gamma(x) =\)\(\omega^{x}\sec\left( \frac{\pi x}{2} \right)\int_{0}^{\infty}t^{x - 1}\)\(\cos\omega tdt,0 < x < 1,\omega > 0\)
.

Као што је познато, за природне бројеве је \(\Gamma(n + 1) = n!\), а
интересантне су и вредности у тачкама \(x = \frac{2k + 1}{2},k \geq 0\);
тада је
\(\Gamma(n + \frac{1}{2}) = \frac{(2n)!}{4^{n} \cdot n!}\sqrt{\pi}\)
{[}10{]}. Неке вредности ове функције дате у табели 3.1.

{\def\LTcaptype{none} % do not increment counter
\begin{longtable}[]{@{}
  >{\centering\arraybackslash}p{(\linewidth - 16\tabcolsep) * \real{0.1111}}
  >{\centering\arraybackslash}p{(\linewidth - 16\tabcolsep) * \real{0.1111}}
  >{\centering\arraybackslash}p{(\linewidth - 16\tabcolsep) * \real{0.1111}}
  >{\centering\arraybackslash}p{(\linewidth - 16\tabcolsep) * \real{0.1111}}
  >{\centering\arraybackslash}p{(\linewidth - 16\tabcolsep) * \real{0.1111}}
  >{\centering\arraybackslash}p{(\linewidth - 16\tabcolsep) * \real{0.1111}}
  >{\centering\arraybackslash}p{(\linewidth - 16\tabcolsep) * \real{0.1111}}
  >{\centering\arraybackslash}p{(\linewidth - 16\tabcolsep) * \real{0.1111}}
  >{\centering\arraybackslash}p{(\linewidth - 16\tabcolsep) * \real{0.1111}}@{}}
\toprule\noalign{}
\begin{minipage}[b]{\linewidth}\centering
\[x\]
\end{minipage} & \begin{minipage}[b]{\linewidth}\centering
\[\frac{1}{2}\]
\end{minipage} & \begin{minipage}[b]{\linewidth}\centering
\[1\]
\end{minipage} & \begin{minipage}[b]{\linewidth}\centering
\[\frac{3}{2}\]
\end{minipage} & \begin{minipage}[b]{\linewidth}\centering
\[2\]
\end{minipage} & \begin{minipage}[b]{\linewidth}\centering
\[\frac{5}{2}\]
\end{minipage} & \begin{minipage}[b]{\linewidth}\centering
\[3\]
\end{minipage} & \begin{minipage}[b]{\linewidth}\centering
\[\frac{7}{2}\]
\end{minipage} & \begin{minipage}[b]{\linewidth}\centering
\[4\]
\end{minipage} \\
\midrule\noalign{}
\endhead
\bottomrule\noalign{}
\endlastfoot
\(\Gamma(x)\) & \(\sqrt{\pi}\) & \(1\) & \(\frac{\sqrt{\pi}}{2}\) &
\(1\) & \(\frac{3\sqrt{\pi}}{4}\) & \(2\) & \(\frac{15\sqrt{\pi}}{8}\) &
\(6\) \\
\end{longtable}
}

Табела 3.1: Вредности гама функције за неке вредности аргумента.

На слици 3.1. приказан је график гама функције за
\(x \in {\mathbb{R}}^{+}\).

\includegraphics[width=6.43784in,height=4.00365in,alt={pasted-movie.png}]{media/image4.png}

Слика 3.1.1: График гама функције.

\begin{itemize}
\item
  \protect\phantomsection\label{_Toc12}{}Особине гама функције
\end{itemize}

Као што је раније поменуто, гама функција има значајне рекурентне везе,
на пример:

\(\Gamma(1 + x) = x \cdot \Gamma(x)\),

\(\Gamma(x - 1) =\)\(\frac{\Gamma(x)}{x - 1}\),

\(\Gamma(x + n) = x(x + 1)(x + 2)...(x + n - 1)\Gamma(x)\),

\(\Gamma(x - n) =\)\(\frac{\Gamma(x)}{(x - 1)(x - 2)...(x - n + 1)(x - n)}\).

Неке од најважнијих особина ове функције су:

\(\Gamma(x) \cdot \Gamma(1 - x) =\)\(\frac{\pi}{\sin\pi x}\),
рефлексиона формула,

\(\Gamma(2x) =\)\(\frac{2^{2x - 1}}{\sqrt{\pi}}\)\(\Gamma(x)\)\(\cdot \Gamma(x + \frac{1}{2})\),
Лежандрова дупликациона формула,

\[\Gamma(nx) =\]\[\frac{n^{nx - \frac{1}{2}}}{{(2\pi)}^{\frac{n - 1}{2}}}\]\[\Gamma(x) \cdot\]\[\Gamma(x + \frac{1}{n}) \cdot \Gamma(x + \frac{2}{n}) \cdot ... \cdot \Gamma(x + \frac{n - 1}{n})\]

\(= \sqrt{\frac{2\pi}{n}}\frac{n^{nx}}{(2\pi)^{\frac{n}{2}}}\overset{n - 1}{\prod_{j = 0}}\left( x + \frac{j}{n} \right)\),
мултипликациона формула. (20)

Из ових особина проистичу и бројне друге, на пример из (20) се изводи
израз {[}41{]}:

\(\Gamma\left( \frac{1}{n} \right) \cdot \Gamma\left( \frac{2}{n} \right) \cdot ... \cdot \Gamma\left( \frac{n - 2}{n} \right) \cdot \Gamma\left( \frac{n - 1}{n} \right) = \frac{{(2\pi)}^{\frac{n - 1}{2}}}{\sqrt{n}}\).

Из интегралних записа могу се извести и изрази {[}6{]}:

\(\frac{\Gamma(x)}{a^{x}} = \int_{0}^{\infty}e^{- at}t^{x - 1}dt,a > 0\),

\(\int_{0}^{\infty}\frac{t^{x - 1}}{{(1 + t)}^{x + y}}dt = \frac{\Gamma(x) \cdot \Gamma(y)}{\Gamma(x + y)}\),

\(\int_{0}^{\frac{\pi}{2}}{(\sin\theta)}^{2x - 1}{(\cos\theta)}^{2y - 1}d\theta = \frac{1}{2}\frac{\Gamma(x) \cdot \Gamma(y)}{\Gamma(x + y)}\),

\(\int_{0}^{1}\frac{t^{m - n}}{\sqrt{1 - t^{n}}} = \frac{\Gamma\left( \frac{m}{n} \right)\sqrt{\pi}}{n \cdot \Gamma\left( \frac{m}{n} + \frac{1}{2} \right)}\).

Такође, интересантно је да двоструки факторијел преко гама функције може
да се дефинише као {[}36{]}:

\(n!! = \left\{ \begin{array}{r}
2^{\frac{n}{2}}\Gamma(\frac{n}{2} + 1),n = 2k,k \in {\mathbb{N}}_{0} \\
\frac{2^{\frac{n + 1}{2}}}{\sqrt{\pi}}\Gamma(\frac{n}{2} + 1),n = 2k + 1,k \in {\mathbb{N}}_{0}
\end{array} \right.\ \).

\begin{itemize}
\item
  \protect\phantomsection\label{_Toc13}{}Апроксимације гама функције
\end{itemize}

Постоји велики број радова чија је тема управо апроксимација гама
функције, те ће у овом одељку бити представљено неколико познатих и
често коришћених метода апроксимације ове функције.

Формула која се често користи због своје једноставности је {[}34{]}:

\(\Gamma(x + 1) \approx\)\(\sqrt{2\pi x}\left( \frac{x}{e} \right)^{x}\).

Ова формула је позната као Стирлингова апроксимација, а постоји пуно
формула насталих побољшањем ове формуле ради боље прецизности.
Стирлинговом формулом се назива и наредна формула и она важи за велике
вредности променљиве {[}31{]}:

\(\Gamma(x + 1) \approx\)\(\sqrt{2\pi}x^{x - \frac{1}{2}}e^{- x}(1 + \frac{1}{12x} + \frac{1}{288x^{2}} - \frac{139}{51840x^{3}} - \frac{571}{2488320x^{4}} + О\left( \frac{1}{x^{5}} \right))\).

Једноставне формуле којe апроксимирају гама фунцкију с већом прецизношћу
од Стирлингове су:

\(\Gamma(x + 1) \approx\)\(\sqrt{2\pi}\left( \frac{x + \frac{1}{2}}{e} \right)^{x + \frac{1}{2}}\)
(Бурнсајд),

\(\Gamma(x + 1) \approx\)\(\sqrt{2\pi\left( x + \frac{1}{6} \right)}\left( \frac{x}{e} \right)^{x}\)
(Госпер),

\(\Gamma(x + 1) \approx\)\(\sqrt{2\pi x}\left( \frac{x}{e} \right)^{x}\left( 1 + \frac{1}{12x^{2} - \frac{1}{10}} \right)^{\frac{x}{2}}\)(Немес),

\(\Gamma(x + 1) =\)\(\sqrt{\pi}\left( \frac{x}{e} \right)^{x}{(8x^{3} + 4x^{2} + x + \frac{1}{30})}^{\frac{1}{6}}\)(Рамануџан).

Мало комплекснија формула која даје исти број тачних децимала као
Немесова формула је {[}11{]}:

\(\Gamma(x + 1) \approx\)\(\sqrt{2\pi x}\left( \frac{x}{e} \right)^{x}\left( x\sin h\frac{1}{x} \right)^{\frac{x}{2}}\)(Виндшитл).

Ова формула даје више од 8 тачних децимала када је променљива већа од 8.

Ако је неопходна већа прецизност, адекватне су и формуле дате помоћу
верижних разломака:

\(\Gamma(x + 1) \approx\)\(\sqrt{2\pi x}\left( \frac{x}{e} \right)^{x}e^{v}\),
где је
\(v = \frac{\frac{1}{12}}{x + \frac{\frac{1}{30}}{x + \frac{\frac{53}{210}}{x + ...}}}\),
(21)

или

\(\Gamma(x + 1) \approx \sqrt{2\pi x}e^{- x}v^{x}\), где је
\(v = x + \frac{1}{12x - \frac{1}{10x + \frac{\frac{- 2369}{252}}{x + \frac{\frac{2117009}{1193976}}{x + \frac{\frac{393032191511}{1324011300744}}{x + ...}}}}}\).

Формула (21) је фактички побољшање Немесове формуле, па се тако у
литератури некад апросимира користећи разломак
\(\frac{1}{12x - \frac{1}{10x - \frac{2369}{252x}}}\)итд. {[}33{]}.

Неке од најновијих апроксимација су дали Махмуд и Алмуаши:

\(\Gamma(x + 1) \approx\)\(\sqrt{2\pi x}\left( \frac{x}{e} \right)^{x}\left( \frac{3(40x^{2} + 3)}{120x - 1} \right),x \rightarrow \infty\),

\(\Gamma(x + 1) \approx\)\(\sqrt{2\pi x}\left( \frac{x}{e} \right)^{x}\left( x\sin\frac{1}{x} \right)^{\frac{x}{2}}e^{\frac{7}{324x^{3}(35x^{2} + 33)}}\),
(22)

\(\Gamma(x + 1) \approx\)\(\sqrt{2\pi x}\left( \frac{x}{e} \right)^{x}\left( \frac{x^{2} + \frac{7}{60}}{x^{2} - \frac{1}{20}} \right)^{\frac{x}{2}}e^{\frac{461}{907200x^{3}(x^{2} + 33)}}\),
(23)

\(\Gamma(x + 1) \approx\)\(\sqrt{2\pi x}\left( \frac{x}{e} \right)^{x}\left( \frac{x^{2} + \frac{7}{60}}{x^{2} - \frac{1}{20}} \right)^{\frac{x}{2}}e^{\frac{461}{907200x^{5}}}\).
(24)

Аутори формула наводе да је су (22) и (23) прецизније од (24), а
најпрецизнијом од наведених се показала формула (23). У сваком случају,
све од ових формула конвергирају када се променљива приближава
бесконачности, а за мале вредности променљиве дају свега неколико тачних
цифара.

Један од најпрецизнијих познатих метода апроксимације је Ланцошев метод
који има много већу комплексност од свих наведених формула. Овај метод
се користи и у неким од најпопуларнијих математичких алата, попут алата
MATHEMATICA и других.

Ланцошева формула за апроксимацију је:

\(\Gamma(x + 1) \approx\)\(\sqrt{2\pi}{(x + r + \frac{1}{2})}^{x + \frac{1}{2}}e^{- (x + r + \frac{1}{2})}\)\(\overset{\infty}{\sum_{k = 0}}\)\(a_{k}(r)H_{k}(x)\),

где је \(H_{0}(x) = 1\),
\(H_{k}(x) =\)\(\frac{x...(x - k + 1)}{(x + 1)(x + 2)...(x + k)}\).

Коефицијенти \(a_{k}(r)\) се рачунају по формули:

\(a_{k}(r) =\)\({( - 1)}^{k}\)\(c_{k}(r)\),

\(c_{k}(r) =\)\(\sqrt{\frac{2}{\pi} \cdot}e^{r} \cdot k \cdot \overset{k}{\sum_{j = 0}}{( - 1)}^{j} \cdot \frac{(k + j - 1)!}{(k - j)! \cdot j!}\left( \frac{e}{j + r + \frac{1}{2}} \right)^{j + \frac{1}{2}}\).

\begin{itemize}
\item
  \protect\phantomsection\label{_Toc14}{}Бета функција~
\end{itemize}

Бета функција позната је и под именом \emph{Ојлеров интеграл прве
врсте}. Своју примену налази у разним областима, а најчешће у
математичкој анализи, теорији вероватноће, статистици, квантној физици
и, наравно, инжењерству.

\begin{itemize}
\item
  \protect\phantomsection\label{_Toc15}{}Дефиниције бета функције
\end{itemize}

\emph{\textbf{Дефиниција 16:}} Бета функција се, преко гама функције у
ознаци \(\Gamma(x),x \in {\mathbb{R}},x > 0\), дефинише изразом:

\(B(p,q) =\)\(\frac{\Gamma(p) \cdot \Gamma(q)}{\Gamma(p + q)}\), где
\(p,q \in {\mathbb{R}}\)и \(p,q > 0\). (25)

\emph{\textbf{Дефиниција 17}} {[}35{]}\emph{\textbf{:}} Интегрални облик
бета функције је дат изразом:

\(B(p + 1,q + 1) =\)\(\int_{0}^{1}t^{p}{(1 - t)}^{q}dt\).

Веза са Гама функцијом, односно израз (25) из дефиниције 16 се може
извести на следећи начин: ако посматрамо производ две гама функције за
произвољне реалне позитивне вредности \(p\) и \(q\), и искористимо
дефиницију 12 {[}5{]}, добијамо:

\(\Gamma(p) \cdot \Gamma(q) =\)\(4 \cdot \int_{0}^{\infty}s^{2p - 1}e^{- s^{2}}ds \cdot \int_{0}^{\infty}t^{2q - 1}e^{- t^{2}}dt\)

\[= 4 \cdot \int_{0}^{\infty}\int_{0}^{\infty}s^{2p - 1} \cdot e^{- s^{2}} \cdot t^{2q - 1} \cdot e^{- t^{2}}dtds\]

\(= 4 \cdot \int_{0}^{\infty}\int_{0}^{\infty}e^{- (s^{2} + t^{2})} \cdot t^{2q - 1} \cdot s^{2p - 1}dtds\).

Квадрати у експоненту су погодни за пребацивање у поларне координате,
тако да уводимо смену
\(t = r \cdot \cos\theta,s = r \cdot \sin\theta,s^{2} + t^{2} = r^{2},ds \cdot dt = r \cdot dr \cdot d\theta\):

\(\Gamma(p) \cdot \Gamma(q) =\)
\(4 \cdot \int_{r = 0}^{\infty}\int_{\theta = 0}^{\frac{\pi}{2}}e^{- r}{(r\cos\theta)}^{2q - 1}{(r\sin\theta)}^{2p - 1}r \cdot dr \cdot d\theta\)

\[= 4 \cdot \int_{r = 0}^{\infty}e^{- r^{2}} \cdot r^{2(p + q) - 1}dr \cdot \int_{\theta = 0}^{\frac{\pi}{2}}{(\sin\theta)}^{2p - 1}{(\cos\theta)}^{2q - 1}d\theta\]

Први интеграл се решава увођењем смене \(r^{2} = a\):

\(\int_{r = 0}^{\infty}{e^{- r}}^{2} \cdot r^{2(p + q) - 1}dr = \frac{1}{2}\int_{0}^{\infty}e^{- a}a^{p + q - 1}da = \frac{1}{2}\)\(\Gamma(p + q)\).

Други интеграл се увођењем смене
\(\sin^{2}\theta = b,d\theta = 2\sin\theta\cos\theta d\theta\), своди
на:

\[\int_{0}^{\frac{\pi}{2}}{(\sin\theta)}^{2p - 1}{(\cos\theta)}^{2q - 1}d\theta = \int_{0}^{\frac{\pi}{2}}{(\sin^{2}\theta)}^{p}{(\cos^{2}\theta)}^{q}\sin\theta\cos\theta d\theta\]

\[= \frac{1}{2}\int_{0}^{1}b^{p}{(1 - b)}^{q}db\]

\(= \frac{1}{2}\)\(B(p,q)\).

Дакле, сада је:

\(\Gamma(p) \cdot \Gamma(q) =\)\(4 \cdot \frac{1}{2}\)\(\Gamma(p,q) \cdot\)\(\frac{1}{2}\)\(B(p,q) =\)
\(\Gamma(p + q)\) \(B(p,q)\),

и коначно
\(B(p,q) =\)\(\frac{\Gamma(p) \cdot \Gamma(q)}{\Gamma(p + q)}\).

У овом извођењу везе са гама функцијом су произишле и две формуле помоћу
којих се бета функција може дефинисати:

\(B(p,q) =\)\(2\int_{0}^{\frac{\pi}{2}}{(\sin t)}^{2p - 1}{(\cos t)}^{2q - 1}dt\),

\(B(p,q) =\)\(2\int_{0}^{1}t^{2p - 1}{(1 - t)}^{q - 1}dt\).

Помоћу одређеног интеграла бета функција се може представити и као:

\(B(p,q) =\)\(\int_{0}^{1}\frac{t^{p - 1} + t^{q - 1}}{{(1 + t)}^{p + q}}dt = \int_{1}^{\infty}\frac{t^{p - 1} + t^{q - 1}}{{(1 + t)}^{p + q}}dt\).

\emph{\textbf{Дефиниција 18}}: За реалне аргументе \(p,q\) бета функција
се може израчунати помоћу бесконачног производа формулом:

\[B(p,q) =\]\[\overset{\infty}{\prod_{j = 0}}\frac{(j + 1)(j + p + q)}{(j + p)(j + q)}\]

\(= \frac{p + q}{pq}\overset{\infty}{\prod_{j = 0}}\frac{1 + \frac{p + q}{j}}{(1 + \frac{p}{n})(1 + \frac{q}{n})}\).

\emph{\textbf{Дефиниција 19}}: За реалне аргументе \(p,q\) бета функција
се може израчунати помоћу бесконачне суме, формулом {[}35{]}:

\(B(p,q) =\)\(\frac{1}{p} + \frac{1 - q}{1 + p} + \frac{(1 - q)(2 - q)}{2(2 + p)} + ... = \overset{\infty}{\sum_{j = 0}}\frac{{(1 - q)}_{j}}{j! \cdot (p + j)}\).

У табели 4.1 су дате вредности бета функције за неке парове променљивих
\(p,q > 0\).

{\def\LTcaptype{none} % do not increment counter
\begin{longtable}[]{@{}
  >{\centering\arraybackslash}p{(\linewidth - 10\tabcolsep) * \real{0.1666}}
  >{\centering\arraybackslash}p{(\linewidth - 10\tabcolsep) * \real{0.1666}}
  >{\centering\arraybackslash}p{(\linewidth - 10\tabcolsep) * \real{0.1666}}
  >{\centering\arraybackslash}p{(\linewidth - 10\tabcolsep) * \real{0.1668}}
  >{\centering\arraybackslash}p{(\linewidth - 10\tabcolsep) * \real{0.1666}}
  >{\centering\arraybackslash}p{(\linewidth - 10\tabcolsep) * \real{0.1666}}@{}}
\toprule\noalign{}
\begin{minipage}[b]{\linewidth}\centering
\[q\]
\end{minipage} & \begin{minipage}[b]{\linewidth}\centering
\[B(\frac{1}{2},q)\]
\end{minipage} & \begin{minipage}[b]{\linewidth}\centering
\[B(1,q)\]
\end{minipage} & \begin{minipage}[b]{\linewidth}\centering
\[B(\frac{3}{2},q)\]
\end{minipage} & \begin{minipage}[b]{\linewidth}\centering
\[B(2,q)\]
\end{minipage} & \begin{minipage}[b]{\linewidth}\centering
\[B(\frac{5}{2},q)\]
\end{minipage} \\
\midrule\noalign{}
\endhead
\bottomrule\noalign{}
\endlastfoot
\(\frac{1}{2}\) & \(\pi\) & \(2\) & \(\frac{3\pi}{4}\) & \(\frac{4}{3}\)
& \(\frac{3\pi}{8}\) \\
\(1\) & \(2\) & \(1\) & \(\frac{3}{8}\) & \(\frac{1}{2}\) &
\(\frac{2}{15}\) \\
\(\frac{3}{2}\) & \(\frac{3\pi}{4}\) & \(\frac{3}{8}\) &
\(\frac{\pi}{8}\) & \(\frac{4}{15}\) & \(\frac{3\pi}{16}\) \\
\(2\) & \(\frac{4}{3}\) & \(\frac{1}{2}\) & \(\frac{4}{15}\) &
\(\frac{1}{6}\) & \(\frac{4}{35}\) \\
\(\frac{5}{2}\) & \(\frac{3\pi}{8}\) & \(\frac{2}{15}\) &
\(\frac{3\pi}{16}\) & \(\frac{4}{35}\) & \(\frac{3\pi}{128}\) \\
\end{longtable}
}

Табела 4.1: Вредности бета функције за неке парове аргумената.

На слици 4.1.1 налази се графички приказ бета функције. На \emph{x} и
\emph{y} осама су представљене вредности параметара \emph{p} и \emph{q},
редом, а вредост бета функције је представљена као висина изнад равни
\((x,y)\), односно вредност \(B(p,q)\) је висина изнад тачке \((p,q)\).
Тамнија боја на графику приказује ниже апсолутне вредности \(B(p,q)\), а
јарке боје више апсолутне
вредности.\includegraphics[width=5.82788in,height=4.38505in,alt={pasted-movie.png}]{media/image5.png}

Слика 4.1.1: Графички приказ бета функције.

\begin{itemize}
\item
  \protect\phantomsection\label{_Toc16}{}Особине бета функције
\end{itemize}

Као и гама функција, и бета функција има пуно корисних особина.

Из саме дефиниције 16 је јасно да важи:

\(B(p,q) = B(q,p)\).

Значајне су рекурзивне везе:

\(B(p + 1,q) =\)\(\frac{p}{p + q}\)\(B(p,q)\), (26)

\(B(p,q + 1) =\)\(\frac{q}{p + q}\)\(B(p,q)\). (27)

Ова тврђења се лако могу извести, посматрањем дефиниције 16, помоћу гама
функције:

\[B(p + 1,q) =\]\[\frac{\Gamma(p + 1) \cdot \Gamma(q)}{\Gamma(p + 1 + q)}\]

\[= \frac{p \cdot \Gamma(p) \cdot \Gamma(q)}{(p + q) \cdot \Gamma(p + q)} =\]\[\frac{p}{p + q} \cdot \frac{\Gamma(p) \cdot \Gamma(q)}{\Gamma(p + q)}\]

\(= \frac{p}{p + q}\)\(B(p,q)\),

што и јесте тврђење (26). Тврђење (27) се изводи аналогно {[}10{]} .

На сличан начин се изводе и особине {[}5{]}:

\(B(p + 1,q) =\)\(\frac{p}{q} \cdot B\)\((p,q + 1)\), (28)

\(B(p,q) = B(p + 1,q) + B(p,q + 1)\), (29)

\(B(p,q) =\)\(\frac{q - 1}{p}\)\(B(p + 1,q - 1)\), (30)

\(B(p,q) \cdot B(p + q,r) = B(q,r) \cdot B(p,q + r)\). (31)

Дакле, израз (28) се добија тако што леву страну израза развијемо по
дедфиницији 16 и помножимо са \(\frac{q}{q}\):

\[B(p + 1,q) =\]\[\frac{\Gamma(p + 1) \cdot \Gamma(q)}{\Gamma(p + 1 + q)} =\]

\[= \frac{p \cdot \Gamma(p) \cdot \Gamma(q)}{\Gamma(p + 1 + q)} \cdot \frac{q}{q} =\]\[\frac{p}{q} \cdot \frac{\Gamma(p) \cdot \Gamma(q + 1)}{\Gamma(p + 1 + q)}\]

\(= \frac{p}{q}B(p,q + 1)\),

док се израз (29) добија полазећи од десне стране и исте дефиниције
(дефиниција 16):

\[B(p + 1,q) + B(p,q + 1) =\]\[\frac{\Gamma(p + 1) \cdot \Gamma(q)}{\Gamma(p + 1 + q)} + \frac{\Gamma(p) \cdot \Gamma(q + 1)}{\Gamma(p + q + 1)} =\]

\[= \frac{p \cdot \Gamma(p) \cdot \Gamma(q) + q \cdot \Gamma(p) \cdot \Gamma(q)}{\Gamma(p + 1 + q)}\]\[= \frac{(p + q) \cdot \Gamma(p) \cdot \Gamma(q)}{(p + q)\Gamma(p + q)} =\]\[\frac{\Gamma(p) \cdot \Gamma(q)}{\Gamma(p + q)}\]

\(= B(p,q)\).

Трећи израз се изводи (30) тако што се лева страна развије по дефиницији
16 и потом помножи са \(\frac{p}{p}\) да би се добило \(\Gamma(p + 1)\):

\(B(p,q) =\)\(\frac{\Gamma(p) \cdot \Gamma(p)}{\Gamma(p + q)}\) \(=\)

\[= \frac{p \cdot \Gamma(p) \cdot \Gamma(q - 1) \cdot (q - 1)}{p \cdot \Gamma(p + q)} =\]\[\frac{q - 1}{p} \cdot \frac{\Gamma(p + 1) \cdot \Gamma(q - 1)}{\Gamma(p + 1 + q - 1)} =\]

\(=\)\(\frac{q - 1}{p}B(p + 1,q - 1)\).

Израз (31) се добија аналогно као претходна три:

\[B(p,q) \cdot B(p + q,r) =\]\[\frac{\Gamma(p) \cdot \Gamma(q)}{\boxed{\Gamma(p + q)}} \cdot \frac{\boxed{\Gamma(p + q)} \cdot \Gamma(r)}{\Gamma(p + q + r)} =\]

\[= \frac{\Gamma(p) \cdot \Gamma(q) \cdot \Gamma(r)}{\Gamma(p + q + r)} \cdot \frac{\Gamma(q + r)}{\Gamma(q + r)}\]\[= \frac{\Gamma(q) \cdot \Gamma(r)}{\Gamma(q + r)} \cdot \frac{\Gamma(p) \cdot \Gamma(q + r)}{\Gamma(p + q + r)} =\]

\(=\) \(B(q,r) \cdot B(p,q + r)\).

Ако је барем један од аргумената природан број, бета функција може да се
израчуна применом формуле {[}35{]}:

\[B(p,n) =\]\[\frac{1 \cdot 2 \cdot 3...(n - 1)}{p(p + 1)...(p + n - 1)} =\]

\(=\)\(\frac{1}{p \cdot \left( \begin{array}{r}
p + n - 1 \\
n - 1
\end{array} \right)}\).

За бета функцију важе следеће везе {[}4{]}{[}5{]}:

\(B(p,q) =\)\(\overset{\infty}{\sum_{j = 0}}B\)\((p + j,q + 1)\),

\(B(p + 1,q + 1) =\)\(2^{- (p + q + 1)}\int_{- 1}^{1}{(1 + x)}^{p}{(1 - x)}^{q}dx\),

\(B(\frac{p + 1}{2},\frac{q - p}{2}) = 2\int_{0}^{\infty}\frac{{sh}^{p}x}{{ch}^{q}x}dx\),

\(B(p,q) =\)\(\frac{(p + q)(p + 1 + 1)}{q(q + 1)}\)\(B(p,q + 2) = ... =\)\(\frac{{(p + q)}_{n}}{q_{n}}\)\(B(p,q + n)\).

Бета функција је због своје природе такође погодна за доказивање многих
теорема, посебно у вези са гама функцијом будући да је једна од њених
дефиниција управо преко производа гама функције. Једна од таквих формула
која се доказује помоћу бета функције је Лежандрова дупликациона формула

\(\Gamma(2x) =\)\(\frac{2^{2x - 1}}{\sqrt{\pi}}\)\(\Gamma(x)\)\(\cdot \Gamma(x + \frac{1}{2})\),

дата у поглављу 4. Нека су оба аргумента бета функције
\(p = q = x + \frac{1}{2}\):

\(B(x + \frac{1}{2},x + \frac{1}{2}) = \int_{0}^{1}t^{x - \frac{1}{2}}{(1 - t)}^{x - \frac{1}{2}}dt\).

Увођењем смене \(s = 1 - 2t,ds = - 2dt\), добија се:

\(B(x + \frac{1}{2},x + \frac{1}{2}) = \int_{0}^{1}t^{x - \frac{1}{2}}{(1 - t)}^{x - \frac{1}{2}}dt\)
\(= \int_{1}^{- 1}\)\(\left( \frac{1 + s}{2}\left( 1 - \frac{1 + s}{2} \right) \right)^{x - \frac{1}{2}}\)\(( - 2ds) =\)

\[= 2\int_{- 1}^{1}{(\frac{1 + s}{2}\frac{1 - s}{2})}^{x - \frac{1}{2}}ds =\]\[2^{- 2x + 1}\int_{0}^{1}{(1 - s^{2})}^{x - \frac{1}{2}}ds\]\[=\]

\(= 2^{- 2x} \cdot B(\frac{1}{2},x + \frac{1}{2})\).

Ако сад обе стране заменимо по дефиницији 16, имамо:

\(\frac{\Gamma(x + \frac{1}{2})\Gamma(x + \frac{1}{2})}{\Gamma(2x + 1)} = 2^{- 2x}\frac{\Gamma(\frac{1}{2})\Gamma(x + \frac{1}{2})}{\Gamma(x + 1)}\),

што је израз који се једноставно може трансформисати у Лежандрову
дупликациону формулу, што је и требало доказати {[}10{]}.

За бета функцију важи и Гаусова мултипликациона формула:

\(B(np,nq) =\)\(\frac{\Gamma(np) \cdot \Gamma(nq)}{\Gamma(n(p + q))}\)
\(=\)

\(= n^{- nq}\frac{B(p,q) \cdot B(p + \frac{1}{n},q)...B(p + \frac{n - 1}{n},q)}{B(q,q) \cdot B(2q,q) \cdot ... \cdot B((n - 1)q,q)}\).

Још неке значајне релације су {[}51{]}:

\(B(p,p)\)\(\cdot B(p + \frac{1}{2},p + \frac{1}{2}) = \frac{\pi}{2^{4p - 1} \cdot p}\),

\[\overset{2n}{\prod_{j = 0}}B(\frac{j}{2n + 1} + p,\frac{j}{2n + 1} + q) =\]

\(\frac{{(2n + 1)}^{\frac{2n + 1}{2}}\pi^{n} \cdot B(n,\frac{1}{2}\lbrack(q + p)(2n + 1) + 1\rbrack) \cdot B(p(2n + 1),q(2n + 1))}{(n - 1)!}\)
за непарно \(n\) ,

\(\overset{2n - 1}{\prod_{j = 0}}B(\frac{j}{2n} + p,\frac{j}{2n} + q) =\)\(\frac{n^{n} \cdot \pi^{n} \cdot B(n,2(p + q)n) \cdot B(2pn,2qn)}{2^{2(p + q)n - n - 1}(n - 1)! \cdot B((p + q)n,(p + q + 1)n)}\).

\begin{itemize}
\item
  \protect\phantomsection\label{_Toc17}{}Ортогонални полиноми~
\end{itemize}

Постоји неколико врста ортогоналних полинома, а неки од њих су споменути
већ у самом уводу рада - Лагерови, Ермитови и Јакобијеви полиноми у које
спадају Лежандрови и Чебишевљеви полиноми (прве и друге врсте). Сви
набројани спадају у тзв. класичне ортогоналне полиноме, који су настали
решавањем диференцијалних једначина другог реда.

Ови полиноми имају широку примену: у стохастичкој анализи, Гаусовим
квадратурним формулама, при минимизацији грешке израчунавања у
нумеричкој математици, теорији вероватноће, комбинаторици, теорији
бројева, затим у физици, техничким наукама, итд.

Поред класичних ортогоналних полинома, постоје и дискретни ортогонални
полиноми, као и полиноми ортогонални на полукругу, али у овом раду ће
фокус бити на класичним {[}47{]}.

Постоји неколико начина да се дефинишу ортогонални полиноми као
партикуларна решења диференцијалних једначина Фухсовог типа, помоћу
тежинске функције, генеративне функције, итд. За потребе овог рада
увешћемо дефиницију помоћу тежинске функције.

\emph{\textbf{Дефиниција 20}} {[}53{]}\emph{:} За низ полинома
\(A_{0}(x),A_{1}(x),A_{2}(x),...\) се каже да је ортогоналан на сегменту
\(\lbrack a,b\rbrack\) ако је њихов скаларни производ у односу на
тежинску функцију \(p(x)\) једнак нули, односно испуњавају једнакост:

\(\int_{a}^{b}\)\(p(x)A_{m}(x)A_{n}(x)dx = 0,\forall m \neq n\)
{[}17{]}.

За низ полинома се каже да је ортономиран на сегменту
\(\lbrack a,b\rbrack\) ако је

\[\int_{a}^{b}\]\[p(x)A_{m}(x)A_{n}(x)dx = \delta_{mn},\forall m,n\]

при чему се \(\delta_{mn}\) назива Кронекеров симбол, односно норма
полинома.

Од сваког низа полинома може се конструисати тачно један ортогонални низ
полинома, коришћењем Грам-Шмитовог поступка ортогонализације.

Ортогонални полиноми имају неколико интересантних особина:

\begin{itemize}
\item
  За три узастопна полинома ортогоналног низа постоји рекурентна
  релација:
\end{itemize}

\(A_{n + 1}(x) - (\alpha_{n}x + \beta_{n})A_{n}(x) + \gamma_{n}A_{n - 1}(x) = 0\),

где су \(\alpha_{n},\beta_{n}\) и \(\gamma_{n}\) константе.

\begin{itemize}
\item
  Све нуле полинома ортогоналног низа су реалне, једноструке и припадају
  интервалу \(\lbrack a,b\rbrack\) {[}29{]}.
\item
  Нека су \(x_{0} < x_{1} < ... < x_{k - 1} < x_{k}\) нуле полинома
  \(A_{k}(x)\) и \(x_{0} = a\) и \(x_{k} = b\). Тада сваки
  интервал\(\lbrack x_{\mu},x_{\mu + 1}\rbrack\)за \(\mu = 0,1,...k\)
  садржи тачно један корен полинома \(A_{k + 1}(x)\). Између два корена
  \(A_{k}(x)\) постоји бар један корен \(A_{m}(x)\) ако је \(m > n\).
\item
  Нека је \(c\) реална константа. Тада полином
  \(p_{k + 1}(x) - cp_{k}(x)\) има тачно \(k + 1\) реалних различитих
  нула. Ако је констата позитивна ове нуле леже унутар интервала
  \(\lbrack a,b\rbrack\) изузев највеће, која лежи у овом интервалу само
  ако је константа мања или једнака количнику полинома \(p_{k + 1}\) и
  \(p_{k}\) у тачки \(b\). Аналогно, ако је константа негативна, све
  нуле датог полинома су унутар интервала осим најмање, која у њему лежи
  под условом да је константа већа или једнака количнику \(p_{k + 1}\) и
  \(p_{k}\) у тачки \(a\).
\end{itemize}

\begin{itemize}
\item
  \protect\phantomsection\label{_Toc18}{}Диференцијална једначина
\end{itemize}

\emph{\textbf{Дефиниција 21}}: Претпоставимо да функција \(\rho(x)\)
задовољава следећу диференцијалну једначину првог реда:

\(\rho'(x) =\)\(\frac{D + Ex}{A + Bx + Cx^{2}}\)\(\rho(x)\),

где су \(A,B,C,D,E\) константе и претпоставимо да
\((A + Bx + Cx^{2})\rho(x)\) тежи нули на крајевима интервала. Овако
дефинисана једначина назива се \emph{Пирсонова диференцијална
једначина}. У случају када је посматрани интервал бесконачан, треба
увести и додатну хипотезу: производ \(\rho(x)\) са произвољним полиномом
се приближава нули када се \(x\) приближава бесконачности.

Тежинске функције које задовољавају Пирсонову диференцијалну једначину
су од велике важности - теоријски и практично. Помоћу ових функција ћемо
увести теорему о диференцијалној једначини за ортогоналне полиноме.

\emph{\textbf{Теорема 6:}} Нека је
\(\{ p_{0}(x),p_{1}(x),p_{2}(x),...,p_{n}(x),...\}\) ортономиран систем
полинома на интервалу \((a,b)\) у односу на тежинску функцију
\(\rho(x)\) која задовољава Пирсонову једначину. Тада сваки полином овог
реда \(p_{n}(x)\) задовољава следећу диференцијалну једначину {[}38{]}:

\((A + Bx + Cx^{2})p_{n}^{''}(x) + \lbrack(B + D) + (2C + E)x\rbrack p_{n}'(x) - \lbrack C_{n}(n + 1) + E_{n}\rbrack p_{n}(x) = 0\).

\emph{Доказ:} Полазимо од диференцијалне једначине:

\(G(x)p_{n}''(x) + H(x)p_{n}'(x) + \gamma_{n}p_{n} = 0\) (32),

где је \(G(x)\) полином другог реда, у овом случају
\(G(x) = A + Bx + Cx^{2}\), \(H(x)\) је полином првог реда независан од
\emph{n}, док је \(\gamma_{n}\) у функцији од \(n\), а не зависи од
\(x\). Доказаћемо да полином \(p_{n}(x)\) задовољава једначину (32).

Нека је \(q_{m}(x)\) произвољан полином степена мањег од \(n\). Када на
интеграл

\(I = \int_{a}^{b}\)\(q_{m}(x)\)\(\frac{d}{dx}\)\(\lbrack G(x)\rho(x)p_{n}'(x)\rbrack dx\)

применимо парцијалну интеграцију:

\(u = q_{m}(x)\) \(\Rightarrow du = q_{m}'(x)dx\),

\(dv =\)\(\frac{d}{dx}\)\(\lbrack G(x)\rho(x)p_{n}'(x)\rbrack dx\)
\(\Rightarrow v = G(x)\rho(x)p_{n}'(x)\),

\(I = uv|_{a}^{b} -\)\(\int_{a}^{b}vdu =\)\(\boxed{q_{m}(x)G(x)\rho(x)p_{n}'(x)|_{a}^{b}}\)\(-\)\(\int_{a}^{b}\)\(G(x)\rho(x)p_{n}'(x)q_{m}'(x)dx\).

Будући да \(uv\) садржи производ \(G(x)\rho(x)\) са полиномом, нестаје
на крајевима интервала и тиме добијамо:
\(I = - \int_{a}^{b}\)\(G(x)\rho(x)p_{n}'(x)q_{m}'(x)dx\).

Поновном применом парцијалне интеграције, \(u = G(x)\rho(x)q_{m}'(x)\) и
\(dv = p_{n}'\)\((x)dx\), слично претходној интеграцији поново нестаје
\(uv\), па је резултат:

\(I = \int_{a}^{b}\)\(p_{n}(x)\)\(\frac{d}{dx}\)\(\lbrack G(x)\rho(x)q_{m}'(x)\rbrack dx\).

Сада даље рачунамо
\(\frac{d}{dx}\)\(\lbrack G(x)\rho(x)q_{m}'(x)\rbrack\) замењујући
вредност \(G(x)\):

\[\frac{d}{dx}\]\[\lbrack G(x)\rho(x)q_{m}'(x)\rbrack =\]\[\frac{d}{dx}\]\[\lbrack(A + Bx + Cx^{2})\rho(x)q_{m}'(x)\rbrack\]

\(= (B + 2Cx)\rho(x)q_{m}'(x) + (A + Bx + Cx^{2})\rho'(x)q_{m}'(x) + (A + Bx + Cx^{2})\rho(x)q_{m}^{''}(x)\).

Будући да је
\(\rho'(x) =\)\(\frac{D + Ex}{A + Bx + Cx^{2}}\)\(\rho(x)\), добија се:

\[\frac{d}{dx}\]\[\lbrack G(x)\rho(x)q_{m}'(x)\rbrack\]\[= (B + 2Cx)\rho(x)q_{m}'(x) + (D + Ex)\rho(x)q_{m}'(x) + (A + Bx + Cx^{2})\rho(x)q_{m}^{''}(x)\]

\(= \rho(x)\lbrack(B + 2Cx)q_{m}'(x) + (D + Ex)q_{m}'(x) + (A + Bx + Cx^{2})q_{m}^{''}(x)\rbrack\).

Можемо цео
израз\(\lbrack(B + 2Cx)q_{m}'(x) + (D + Ex)q_{m}'(x) + (A + Bx + Cx^{2})q_{m}^{''}(x)\rbrack\)
обележити са \(r_{m}(x)\), будући да је \(q_{m}'(x)\) степена \(m - 1\),
као и \(q_{m}^{''}(x)\) степена \(m - 2\), па је цео израз степена
највише \(m\). Уводећи ову смену, следи:

\(I = \int_{a}^{b}\)\(\rho(x)p_{n}(x)r_{m}(x)dx\).

По дефиницији ортогоналних полинома, они испуњавају услов

\(\int_{a}^{b}\)\(p(x)A_{m}(x)A_{n}(x)dx = 0,\forall m \neq n\),

што имплицира да је вредност интеграла \(I\) једнака нули. Дакле, добили
смо:

\(I = \int_{a}^{b}\)\(\rho(x)p_{n}(x)r_{m}(x)dx =\)\(\int_{a}^{b}\)\(q_{m}(x)\)\(\frac{d}{dx}\)\(\lbrack G(x)\rho(x)p_{n}'(x)\rbrack = 0\).

У полазном изразу за овај интеграл имамо:

\[\frac{d}{dx}\]\[\lbrack G(x)\rho(x)p_{n}'(x)\rbrack =\]\[\frac{d}{dx}\]\[\lbrack(A + Bx + Cx^{2})\rho(x)p_{n}'(x)\rbrack\]

\[= (B + 2Cx)\rho(x)p_{n}'(x) + G(x)\rho'(x)p_{n}'(x) + G(x)\rho(x)p_{n}^{''}(x)\]

\(= \rho(x)\lbrack(B + 2Cx)p_{n}'(x) + (D + Ex)p_{n}'(x) + G(x)p_{n}^{''}(x)\rbrack\),

што има облик \(\rho(x)s_{n}(x)\), где \(s_{n}(x)\) представља полином
степена највише \(n\):

\(s_{n}(x) = (B + 2Cx)p_{n}'(x) + (D + Ex)p_{n}'(x) + G(x)p_{n}^{''}(x)\).

Према томе, за произвољан полином\(q_{m}(x)\), \(m < n\), важи:
\(\int_{a}^{b}\)\(\rho(x)s_{n}(x)q_{m}(x)dx = 0\),

односно полином \(s_{n}(x)\) има особину ортогоналности коју за дато
\(n\) поседује само константа помножена с \(p_{n}(x)\). Даље следи да
мора постојати константа \(K_{n}\) таква да је
\(s_{n}(x) = K_{n}p_{n}(x)\).

Овај полином можемо записати и елегантније:

\(s_{n}(x) = (B + Cx)p_{n}'(x) + (D + Ex)p_{n}'(x) + G(x)p_{n}^{''}(x)\)

\(= (A + Bx + Cx^{2})p_{n}^{''}(x) + \lbrack(B + D + (2C + E)x)\rbrack p_{n}'(x)\).

Упоређивањем коефицијената уз \(x^{n}\), где водећи коефицијент у
полиному \(p_{n}(x)\) означавамо са \(a_{nn}\), добијамо:

\(C(n - 1)a_{nn} + (C + E)na_{nn} = K_{n}a_{nn}\),

\(K_{n} = Cn(n + 1) + En\).

Одавде коначно следи:

\((B + 2Cx)p_{n}'(x) + (D + Ex)p_{n}'(x) + G(x)p_{n}^{''}(x) = K_{n}p_{n}(x)\),

а одатле даље:

\((A + Bx + Cx^{2})p_{n}^{''}(x) + \lbrack(B + D) + (2C + E)x\rbrack p_{n}'(x) - \lbrack Cn(n + 1) + En\rbrack p_{n}(x) = 0\),

чиме је теорема доказана. \(\square\)

Ова теорема је применљива на све врсте ортогоналних полинома, нпр., у
случају Лежандорвих полинома узима се \(\rho(x) = 1\) а интервалу
\((a,b)\) одговара интервал \(\lbrack - 1,1\rbrack\); у случају
Eрмитових полинома је \(\rho(x) =\)\(e^{- x^{2}}\), а интервал
\(\lbrack - \infty,\infty\rbrack\), итд {[}38{]}.

\begin{itemize}
\item
  \protect\phantomsection\label{_Toc19}{}Лежандрови полиноми
\end{itemize}

\begin{itemize}
\item
  \protect\phantomsection\label{_Toc20}{}Лежандрова диференцијална
  једначина
\end{itemize}

\emph{\textbf{Дефиниција 22:}} Лежандрова диференцијална једначина је
линеарна диференцијална једначина другог реда чији је облик:

\((1 - x^{2})y^{''}(x) - 2xy' + k(k + 1)y = 0\). (33)

Будући да је ова једначина другог реда, она мора имати два линеарно
независна решења, а њено опште решење када је \(k\) цео ненегативан број
је:

\(y = a_{1}P_{k}(x) + a_{2}Q_{k}(x)\),

где су \(a_{1},a_{2}\) константе, \(P_{k}(x)\) је Лежандров полином
парног експонента, док је \(Q_{k}(x)\)Лежандров полином непарног
експонента. Решење \(P_{k}(x)\) је Лежандрова функција прве врсте, а
\(Q_{k}(x)\) је Лежандрова функција друге врсте, сингуларна у тачкама
\(\pm 1\) {[}2{]}.

Када се цела једначина нормализује, тј. подели са \((1 - x^{2})\) добија
се:

\(y^{''}\)\((x) -\)\(\frac{2x}{1 - x^{2}}\)\(y'\)\((x) +\)\(\frac{k(k - 1)}{1 - x^{2}}y\)\((x) = 0\).

Добијену једначину можемо записати и као:

\(y^{''}(x) - P(x)y'(x) + Q(x)y(x) = 0\),

где је наравно \(P(x) =\)\(\frac{- 2x}{1 - x^{2}}\) и
\(Q(x) =\)\(\frac{k(k + 1)}{1 - x^{2}}\).

Када је \(k\) целобројно, функција прве врсте се редукује у Лежандров
полином {[}7{]}.

Ова једначина има сингуларна решења у тачкама 1, -1 и \(\infty\).

Када се променљива посматра као \(x = \cos\varphi\), једначина добија
облик:

\(\frac{d^{2}y}{d\varphi^{2}} + \frac{\cos\varphi}{\sin\varphi}\frac{dy}{d\varphi} +\)\(k(k + 1)y = 0\).

Метод решавања ове једначине је исти као и у сличају Беселове једначине
- тзв. Фробенијусов метод.

Дакле, вршимо развој са \(k = 0\), и имамо:

\(y = \overset{\infty}{\sum_{n = 0}}a_{n}x^{n}\),

\(y' = \overset{\infty}{\sum_{n = 0}}na_{n}x^{n - 1}\),

\(y^{''} = \overset{\infty}{\sum_{n = 0}}n\)\((n - 1)\)\(a_{n}x^{n - 2}\).

Заменом у једначину се долази до:

\((1 - x^{2})\)\(\overset{\infty}{\sum_{n = 0}}\)\(n(n - 1)a_{n}\)\(x^{n - 2} - 2x\)\(\overset{\infty}{\sum_{n = 0}}\)\(na_{n}x^{n - 1} + k(k + 1)\)\(\overset{\infty}{\sum_{n = 0}}\)\(a_{n}x^{n} = 0\),

\(\overset{\infty}{\sum_{n = 0}}\)\(n(n - 1)a_{n}x^{n - 2} -\)\(\overset{\infty}{\sum_{n = 0}}\)\(n(n - 1)a_{n}x^{n} - 2x\)\(\overset{\infty}{\sum_{n = 0}}\)\(na_{n}x^{n - 1} + k(k + 1)\)\(\overset{\infty}{\sum_{n = 0}}\)\(a_{n}x^{n} = 0\),

\(\overset{\infty}{\sum_{n = 2}}\)\(n(n - 1)a_{n}x^{n - 2} -\)\(\overset{\infty}{\sum_{n = 0}}\)\(n(n - 1)a_{n}x^{n} - 2\)\(\overset{\infty}{\sum_{n = 0}}\)\(na_{n}x^{n} + k(k + 1)\)\(\overset{\infty}{\sum_{n = 0}}\)\(a_{n}x^{n} = 0\),

\(\overset{\infty}{\sum_{n = 0}}\)\((n + 2)(n + 1)a_{n + 2}x^{n} -\)\(\overset{\infty}{\sum_{n = 0}}\)\(n(n - 1)a_{n}x^{n} - 2\)\(\overset{\infty}{\sum_{n = 0}}\)\(na_{n}x^{n} + k(k + 1)\)\(\overset{\infty}{\sum_{n = 0}}\)\(a_{n}x^{n} = 0\),

\(\overset{\infty}{\sum_{n = 0}}\)\(\{(n + 1)(n + 2)a_{n + 2} + \lbrack - n(n - 1) - 2n + k(k + 1)\rbrack a_{n}\} = 0\).

Одавде следи да коефицијент уз \(a_{n}\) мора бити једнак 0:

\((n + 1)(n + 2)a_{n + 2} + \lbrack - n(n - 1) - 2n + k(k + 1)\rbrack a_{n} = 0\),

\(a_{n + 2} = \frac{n(n + 1) - k(k + 1)}{(n + 1)(n + 2)}a_{n}\)\(= - \frac{\lbrack k + (n + 1)\rbrack(k - n)}{(n + 1)(n + 2)}a_{n}\).

Сада из претходне једнакости произилази:

\(a_{2} = - \frac{k(k + 1)}{1 \cdot 2}a_{0}\),

\(a_{4} = - \frac{(k - 2)(k + 3)}{3 \cdot 4}a_{2}\)\(= {( - 1)}^{2}\)\(\frac{\lbrack(k - 2)k\rbrack\lbrack(k + 1)(k + 3)\rbrack}{1 \cdot 2 \cdot 3 \cdot 4}a_{0}\),

\(a_{6} = - \frac{(k - 4)(k + 5)}{5 \cdot 6}a_{4} =\)\({( - 1)}^{3}\)\(\frac{\lbrack(k - 4)(k - 2)k\rbrack\lbrack(k + 1)(k + 3)(k + 5)\rbrack}{1 \cdot 2 \cdot 3 \cdot 4 \cdot 5 \cdot 6}a_{0}\).

Дакле, непарно решење би било изражено у облику:

\(y_{1}(x) = 1 +\)\(\overset{\infty}{\sum_{n = 1}}\)\({( - 1)}^{n}\)\(\frac{\lbrack(k - 2n + 2)...(k - 2)k\rbrack\lbrack(k + 1)(k + 3)...(k + 2n - 1)\rbrack}{(2n)!}x^{2n}\),

док би парно било изражено као:

\(y_{2}(x) = x +\)\(\overset{\infty}{\sum_{n = 1}}\)\({( - 1)}^{n}\)\(\frac{\lbrack(k - 2n + 1)...(k - 3)(k - 1)\rbrack\lbrack(k + 2)(k + 4)...(k + 2n)\rbrack}{(2n + 1)!}x^{2n + 1}\).

Када је \(k\) ненегативан паран број, \(y_{1}(x)\) се своди на полином
степена \(k\) који садржи само парне степене, а \(y_{2}(x)\) дивергира.
Аналогно овоме, у случају када је \(k\) ненегативан непаран број,
\(y_{1}(x)\) дивергира, док \(y_{2}(x)\) бива сведено на полином степена
\(k\) који садржи само непарне степене.

Коначно, опште решење је дато Лежандровим полиномом {[}14{]}:

\(P_{n}(x) = c_{n}\)\(\left\{ \begin{array}{r}
y_{1}(x),k = 2i,i = 0,1,2,... \\
y_{2}(x),k = 2i + 1,i = 0,1,2,...
\end{array} \right.\ \),

при чему се \(c_{n}\) бира тако да добијемо \(P_{n}(1) = 1\).

\begin{itemize}
\item
  \protect\phantomsection\label{_Toc21}{}Лежандров полином
\end{itemize}

\emph{\textbf{Дефиниција 23:}} Лежандрови полиноми у ознаци \(P_{n}(x)\)
су ортогонални полиноми са тежинском функцијом \(w(x) = 1\) на интервалу
\(I = ( - 1,1)\), чији је експлицитни развој:

\(P_{n}(x) =\)\(\frac{1}{2^{n}}\overset{n}{\sum_{k = 0}}\left( \begin{array}{r}
n \\
k
\end{array} \right)^{2}\)\({(x - 1)}^{n - k}{(x + 1)}^{k}\),

односно

\(P_{n}(x) =\)\(2^{n}\overset{n}{\sum_{k = 0}}\)\(x^{k}\)\(\left( \begin{array}{r}
n \\
k
\end{array} \right)\left( \begin{array}{r}
\frac{n + k - 1}{2} \\
n
\end{array} \right)\).

Лежандрови полиноми се у литератури могу наћи и под називом Лежандрове
функције прве врсте, Лежандрови коефицијенти, или зонални хармоници.

Осим на дати начин, Лежандров полином се може дефинисати и Родриговом
формулом која ће бити изложена у наредној дефиницији, као и на још
неколико начина.

\emph{\textbf{Дефиниција 24:}} Лежандров полином се изражава помоћу
следеће формуле:

\(P_{n}(x) =\)\(\frac{1}{2^{n}}\overset{n}{\sum_{k = 0}}\frac{1}{2^{n}n!}\frac{d^{n}}{dx^{n}}\)\(\lbrack{(x^{2} - 1)}^{n}\rbrack\),

познате као Родригова формула.

\emph{\textbf{Дефиниција 25}} {[}35{]}\emph{\textbf{:}} Интегрална
репрезентација Лежандровог полинома, позната и под именом Лапласова
репрезентација, је:

\(P_{n}(x) =\)\(\frac{1}{\pi}\int_{0}^{\pi}\left\lbrack x + \sqrt{x^{2} - 1}\cos t \right\rbrack^{n}dt,x > 1\).

\emph{\textbf{Дефиниција 26}} {[}5{]}\emph{\textbf{:}} Лежандров полином
може се израчунати и из:

\(P_{n}(x) =\)\(\overset{\lbrack n/2\rbrack}{\sum_{k = 0}}{( - 1)}^{k}\frac{(2n - 2k)!}{2^{n}k!(n - k)!(n - 2k)!}x^{n - 2k}\),

при чему је \(\lbrack n/2\rbrack\) највећи природан број мањи од
\(n/2\).

Слично Ермитовом полиному, и Лежандров полином се може дефинисати помоћу
контурног интеграла на следећи начин:
\(P_{n}(z) =\)\(\frac{1}{2\pi i}\oint{(1 - 2zt + t^{2})}^{- \frac{1}{2}}t^{- n - 1}dt\).

Лежандров полином је симетрична функција што је лако уочљиво и са
графика. За непарне вредности је централно симетрична у односу на тачку
(0,0) (непарна функција), а за парне вредности је осно симертрична на
\(y\)-осу (парна функција).

Као што се види са графика, домен функције је цео скуп реалних бројева,
мада је ``најбитнији'' део домена сегмент \(\lbrack - 1,1\rbrack\), што
је и домен на коме је полином ортогоналан. Када се аргумент представља
као \(\cos\theta\) функција се и мора ограничити на тај домен. У том
домену налазе се све нуле Лежандрових полинома и локални екстремуми.
Полином реда \(n\) има \(n\) реалних нула.

На следећој слици 6.2.1. се може видети график првих неколико
Лежандрових
полинома.\includegraphics[width=5.44713in,height=3.3571in,alt={pasted-movie.png}]{media/image6.png}

Слика 6.2.1. Првих пет Лежандрових полинома.

Лежaндров полином има следеће понашање, које дати график и потврђује:

\begin{itemize}
\item
  када је \(x > 1\), функција узима вредности \(P_{n}(x) \geq 1\),
\item
  \(P(1) = 1\), сви графици се секу у тачки (1, 1),
\item
  када \(x \in \lbrack - 1,1\rbrack\), вредности функије су у
  интервалу\(\lbrack - 1,1)\),
\item
  \(P_{n}( - 1) = {( - 1)}^{n}\), видимо да се за непарно \(n\) графици
  секу у тачки \(( - 1, - 1)\), а за парно у \(( - 1,1)\),
\item
  када \(x < - 1\), функција се понаша различито у зависности од
  парности реда, па тако за парно \(n\) функција узима поново вредности
  \(P_{n}(x) \geq 1\), а за непарно \(n\) фунција је
  \(P_{n}(x) \leqslant - 1\).
\end{itemize}

Експлицитни облик првих неколико Лежандрових полинома је:

\(P_{0}(x) = 1\),

\(P_{1}(x) = x\),

\(P_{2}(x) =\)\(\frac{1}{2}(3x^{2} - 1)\),

\(P_{3}(x) =\)\(\frac{1}{2}(5x^{3} - 3x)\),

\(P_{4}(x) =\)\(\frac{1}{8}(35x^{4} - 30x^{2} + 3)\),

\(P_{5}(x) =\)\(\frac{1}{8}(63x^{5} - 70x^{3} + 15x)\).

Рекурзивна формула, за \(n > = 2\) и \(P_{0}(x) = 1\), \(P_{1}(x) = x\),
је {[}35{]}:

\(P_{n}(x) = \frac{2n - 1}{n}xP_{n - 1}(x) - \frac{n - 1}{n}P_{n - 2}(x)\).

Генератриса лежандрових полинома је {[}32{]}:

\(g(x,t) =\)\(\frac{1}{\sqrt{1 - 2xt + t^{2}}}\)\(= \overset{\infty}{\sum_{n = 0}}P_{n}\)\((x)t^{n}\).

Њеним диференцирањем се могу добити бројне рекурентне везе које важе за
ове полиноме :

\(P_{n + 1}'(x) - P_{n - 1}'(x) = (2n + 1)P_{n}(x)\),

\(P_{n + 1}'(x) = (n + 1)P_{n}(x) + xP_{n}'(x)\),

\(P_{n - 1}'(x) = - nP_{n}(x) + xP_{n}'(x)\),

\(P_{n}'(x) = xP_{n + 1}'(x) + nP_{n - 1}(x)\),

\((1 - x^{2})P_{n}'(x) = nP_{n - 1}(x) - nxP_{n}(x)\),

\((1 - x^{2})P_{n}'(x) = (n + 1)xP_{n}(x) - (n + 1)P_{n + 1}(x)\).

Још неке корисне релације између ових полинома и њихових извода су
{[}10{]}:

\(P_{n}'(x) =\)\(\overset{\frac{n - 1}{2}}{\sum_{k = 0}}\)\((2n - 4k - 1)P_{n - 2k - 1}(x)\),

\(xP_{n}(x) =\)\(\frac{n + 1}{2n + 1}\)\(P_{n + 1}(x) +\)\(\frac{n}{2n + 1}\)\(P_{n - 1}(x)\),

\(\overset{n}{\sum_{k = 0}}\)\((2k + 1)P_{k}(x)P_{k}(y) =\)\(\frac{n + 1}{x - y}\)\(\{ P_{n + 1}(x)P_{n}(y) - P_{n}(x)P_{n + 1}(y)\}\).

Такође постоји и велики број одређених интеграла који могу бити корисни
у инжењерским прорачунима, на пример {[}35{]}:

\(\int_{- 1}^{1}\)\(P_{n}(t)dt = 0\), \(n \in {\mathbb{N}}\),

\(\int_{- 1}^{1}\frac{P_{n}(t)}{\sqrt{1 - t^{2}}}dt\)\(= \left\{ \begin{array}{r}
0,парноn \\
\pi\left( \frac{(n - 1)!!}{n!!} \right)^{2},непарноn
\end{array} \right.\ \),

\(\int_{- 1}^{1}\frac{P_{n}(t)}{\sqrt{1 - t}}dt = \frac{\sqrt{8}}{2n + 1},n \in {\mathbb{N}}_{0}\).

\begin{itemize}
\item
  \protect\phantomsection\label{_Toc22}{}Лагерови полиноми
\end{itemize}

\begin{itemize}
\item
  \protect\phantomsection\label{_Toc23}{}Лагерова диференцијална
  једначина
\end{itemize}

\emph{\textbf{Дефиниција 27:}} Лагерова диференцијална једначина је
линеарна диференцијална једначина другог реда:

\(xy^{''} + (1 - x)y' + \lambda y = 0\),

при чему једначина има решења која нису сингуларна ако и само ако је
\(\lambda\) ненегативан цео број.

У светлу придружених Лагерових полинома који су такође дефинисани у
претходном делу поглавља, биће дефинисана и придружена Лагерова
једначина чија су партикуларна решења поменути полиноми:

\(xy^{''} + (\alpha + 1 - x)y' + \lambda y = 0\),

где је такође \(\lambda\) ненегативан цео број.

Будући да је Лагерова једначина случај придружене Лагерове једначине, у
наставку ће се говорити у контексту придружене једначине. Ова једначина
се може решити на више начина. Један од начина решавања би био да се
постави хипотеза (Анзац) да је решење у облику бесконачног реда, тј.
Фробенијусов метод. Дакле,
\(y = \overset{\infty}{\sum_{n = 0}}a_{n}x^{n}\) замењујемо у дату
једначину.

\(x\overset{\infty}{\sum_{n = 0}}n\)\((n - 1)\)\(a_{n}x^{n - 2} + (1 - x)\)\(\overset{\infty}{\sum_{n = 0}}na_{n}x^{n - 1} + \lambda\overset{\infty}{\sum_{n = 0}}a_{n}x^{n} = 0\),

\(\overset{\infty}{\sum_{n = 2}}\)\(n(n - 1)a_{n}x^{n - 1} +\)\(\overset{\infty}{\sum_{n = 1}}na_{n}x^{n - 1} - \overset{\infty}{\sum_{n = 1}}na_{n}x^{n} + \lambda\overset{\infty}{\sum_{n = 0}}a_{n}x^{n} = 0\).

Сада се све суме своде на исти степен да би могли да их запишемо као
једну суму:

\(\overset{\infty}{\sum_{n = 1}}\)\(n(n + 1)a_{n + 1}x^{n} +\)\(\overset{\infty}{\sum_{n = 0}}\)\((n + 1)a_{n + 1}x^{n} -\)\(\overset{\infty}{\sum_{n = 1}}na_{n}x^{n} + \lambda\overset{\infty}{\sum_{n = 0}}a_{n}x^{n} = 0\),

\(a_{1} + \lambda a_{0} + \overset{\infty}{\sum_{n = 1}}\)\((((n + 1)n + (n + 1))a_{n + 1} - na_{n} + \lambda a_{n})x^{n} = 0\),

\(a_{1} + \lambda a_{0} + \overset{\infty}{\sum_{n = 1}}\)\(({(n + 1)}^{2}a_{n + 1} + (\lambda - n)a_{n})x^{n} = 0\).

Одавде је јасно да мора бити \(a_{1} + \lambda a_{0} = 0\) и
\({(n + 1)}^{2}a_{n + 1} + (\lambda - n)a_{n} = 0\), односно:

\(a_{1} = - \lambda a_{0}\),

\(a_{n + 1} = \frac{n - \lambda}{{(n + 1)}^{2}}a_{n}\).

Ово даје решење:

\(y = \overset{\infty}{\sum_{n = 0}}a_{n}x^{n} = a_{0}(1 - \lambda x - \frac{\lambda(1 - \lambda)}{2^{2}}x^{2} - \frac{\lambda(1 - \lambda)(2 - \lambda)}{2^{2} \cdot 3^{2}}x^{3} + ...)\).

Ако је \(\lambda\) ненегативан цеи број, решење је заправо
\(y = a_{0}L_{\lambda}(x)\) {[}1{]}.

Још један од начина решавања ове једначине би био помоћу интеграционог
фактора \(\mu\), при чему је:

\(\mu = e^{\int P(x)dx}\) .

Израчунавањем се добија \(\mu = x^{\alpha + 1}e^{- x}\). Решење
диференцијалне једначине је дато као:

\(y(x) = C_{1}\)\(\int\frac{dx}{\mu} + C_{2}\)\(= C_{1}\)\(\int\frac{e^{x}}{x^{\alpha + 1}}dx + C_{2}\).

па произилази да она има регуларно сингуларно решење у нули, као и
нерегуларно у бесконачости.

Једначина може даље бити решена развојем у ред {[}48{]}.

Ваља напоменути и да, ако бисмо желели да докажемо да је Лагеров полином
решење Лагерове диференцијалне једначине, најпогодније би било на начин
да запишемо Лагеров полином у форму контурног интеграла и заменимо га у
саму једначину.

\begin{itemize}
\item
  \protect\phantomsection\label{_Toc24}{}Лагеров полином
\end{itemize}

Као и Лежандрпви полиноми, и Лагерови полиноми се могу дефинисати
Родриговом формулом, што је и њихова најчешћа дефиниција.

\emph{\textbf{Дефиниција 28}} {[}1{]}\emph{\textbf{:}} Лагерови
полиноми, у ознаци
\(L_{0}^{\alpha},L_{1}^{\alpha},...,L_{n}^{\alpha},...\), су ортогонални
полиноми са тежинском функцијом \(w(x) = e^{- x}\), дефинисани над
доменом реалних бројева, изражени Родриговом формулом као:

\(L_{n}(x) =\)\(\frac{1}{w(x)}\left( \frac{d}{dx} \right)^{n}\)\((w(x)p(x)^{n}) =\)\(\frac{e^{x}}{n!}\frac{d^{n}}{dx^{n}}\)\((x^{n}e^{- x}),\)
односно:

\[L_{n}(x) =\]\[\frac{1}{n!}{(\frac{d}{dx} - 1)}^{n}x^{n}\]

за ред полинома \(n \in {\mathbb{N}}_{0}\).

\emph{\textbf{Дефиниција 29}} {[}32{]}\emph{\textbf{:}} Лагерови
полиноми могу бити дати сумом:

\(L_{n}(x) =\)\(\overset{n}{\sum_{k = 0}}\frac{{( - 1)}^{k}}{k!}\left( \begin{array}{r}
n \\
k
\end{array} \right)x^{k}\).

Полазећи од овакве репрезентације лако можемо доћи до експлицитних
израза за првих неколико полинома:

\(L_{0}(x) = 1\),

\(L_{1}(x) = - x + 1\),

\(L_{2}(x) =\)\(\frac{1}{2}(x^{2} - 4x + 2)\),

\(L_{3}(x) =\)\(\frac{1}{6}( - x^{3} + 9x^{2} - 18x + 6)\),

\(L_{4}(x) =\)\(\frac{1}{24}(x^{4} - 16x^{3} + 72x^{2} - 96x + 24)\),

\(L_{5}(x) =\)\(\frac{1}{120}( - x^{5} + 25x^{4} - 200x^{3} + 600x^{2} - 600x + 120)\).

\emph{\textbf{Дефиниција 30}} {[}1{]}: Лагеров полином може бити
представљен у облику контурног интеграла као:

\(L_{n}(x) =\)\(\frac{2}{2\pi i}\oint\frac{e^{\frac{- xz}{1 - z}}}{(1 - z)z^{n + 1}}dz\).

Када су у питању интегралне репрезентације Лагеровог полинома, постоји и
једна специфична таква репрезентација која изражава Лагеров полином
преко Беселовог полинома прве врсте:

\(L_{n}(x) =\)\(\frac{e^{x}}{n!}\int_{0}^{\infty}t^{n}e^{- t}J_{0}\)\((2\sqrt{xt})dt\)
{[}35{]}.

На наредној слици 7.2.1 може се видети првих неколико
полинома.\includegraphics[width=5.74648in,height=3.68329in,alt={pasted-movie.png}]{media/image7.png}

Слика 7.2.1: Графици првих неколико Лагерових полинома.

Полином је дефинисан за \(\lbrack 0,\infty)\). За позитивну вредност
аргумента, узима вредности у интервалу
\(- e^{\frac{x}{2}} \leq L_{n}\)\((x)\)\(\leq e^{\frac{x}{2}}\). Лагеров
полином је непарна функција, полином реда \(n\) има тачно \(n\) корена
на интервалу \(\lbrack 0,\infty)\).

Генератриса Лагеровог полинома је:

\(g(x,t) =\)\(\overset{\infty}{\sum_{t = 0}}\)\(L_{n}(x)t^{n} =\)\(\frac{e^{- \frac{xt}{1 - t}}}{1 - t} = \frac{e^{x}}{2\pi i}\oint\frac{e^{- z}dz}{z - x - tz}\).

Значајније рекурзивне релације су {[}10{]}:

\((n + 1)L_{n + 1}(x) - (2n + 1 - x)L_{n}(x) + nL_{n - 1}(x) = 0\),

\(x\frac{dL_{n}(x)}{dx} =\)\(nL_{n}(x) - nL_{n - 1}(x)\),

\(L_{n}'(x) =\)\(- \overset{n - 1}{\sum_{k = 0}}L_{k}\)\((x)\).

Дакле, на основу формуле за позната прва два полинома можемо једноставно
одредити вредност произвољног Лагеровог полинома {[}30{]}.

Интересантно је споменути и везу између Лагерових и Ермитових полинома
дату релацијама:

\(H_{2n}(x) =\)\({( - 1)}^{n}2^{2n}n!{L_{n}}^{( - 0,5)}\)\((x^{2})\),

\(H_{2n + 1}(x) =\)\({( - 1)}^{n}2^{2n + 1}n!x{L_{n}}^{0,5}\)\((x^{2})\).

Значајне су и наредне формуле {[}35{]}:

\(L_{n}(x + y) =\)\(\frac{1}{n!}( - \frac{1}{4})\overset{n}{\sum_{j = 0}}\left( \begin{array}{r}
n \\
j
\end{array} \right)\)\(H_{2j}(\sqrt{x})H_{2n - 2j}(\sqrt{y})\),

\(L_{n}(bx) =\)\({(1 - b)}^{n}\)\(\overset{n}{\sum_{k = 0}}\left( \begin{array}{r}
n \\
k
\end{array} \right)\left( \frac{b}{1 - b} \right)^{k}\)\(L_{k}(x)\),

\(\overset{n}{\sum_{k = 0}}\)\(L_{k}(x)L_{k}(y) = (n + 1)\)\(\frac{L_{n}(x)L_{n + 1}(y) - L_{n + 1}(x)L_{n}(y)}{x - y}\).

Постоје и придружени, односно генерализовани Лагерови полиноми, у ознаци
\(L_{n}^{\alpha}(x)\)дефинисани Родриговом формулом као:

\(L_{n}^{\alpha}(x) =\)\(\frac{e^{x}x^{- \alpha}}{n!}\frac{d^{n}}{dx^{n}}\)\((e^{- x}x^{n + \alpha})\),

при чему је \(\alpha > - 1\), на интервалу \(I = \lbrack 0, + \infty)\),
а њихова тежинска функција је \(w(x) = x^{\alpha}e^{- x}\).

Обични Лагерови полиноми су специјални случај придружених када је
\(\alpha = 0\). Веза између обичних и придружених Лагерових полинома је:

\(L_{n}^{\alpha}(x) = {( - 1)}^{\alpha}\)\(\frac{d^{\alpha}}{dx^{\alpha}}\)\(L_{n + \alpha}(x)\).

\begin{itemize}
\item
  \protect\phantomsection\label{_Toc25}{}Ермитови полиноми~
\end{itemize}

Разликујемо две врсте Ермитових полинома - физички и пробабилистички. У
овом поглваљу биће обрађене обе врсте полинома и њихове диференцијалне
једначине.

\begin{itemize}
\item
  \protect\phantomsection\label{_Toc26}{}Ермитовa диференцијална
  једначина
\end{itemize}

\emph{\textbf{Дефиниција 31:}} Линеарна диференцијална једначина другог
реда чије је партикуларно решење Ермитов \textbf{физички} полином
\(H_{n}(x)\), \(x \in {\mathbb{R}}\), \(n \in {\mathbb{N}}_{0}\) је:

\(y'' - xy' + ny = 0\) .

Линеарна диференцијална једначина другог реда чије је партикуларно
решење Ермитов \textbf{пробабилистички} полином \(H_{e_{n}}\)\((x)\) је:

\(y'' - 2xy' + ny = 0\).

Да бисмо извели диференцијалну једначину за Ермитов физички полином,
полазимо од његове тежинске функције. Она задовољава следећу једначину:

\(\phi'(x) + x\phi(x) = 0\).

Диференцирањем ове једначине се добија:

\(\phi^{''}(x) + x\phi'(x) + \phi(x) = 0\),

а \(n\) узастопних диференцирања даје резулат

\(\phi^{(n + 1)}(x) + x\phi^{(n)}(x) + n\phi^{(n - 1)}(x) = 0\).

Будући да је по дефиницији Ермитов полином
\(H_{n}(x) =\)\({( - 1)}^{n}\)\(e^{x^{2}}\phi^{n}\)\((\)\(x\)\()\), када
изразимо \(\phi(x)\) преко Ермитовог полинома и заменимо претходну
једначину добија се:

\(\lbrack\)\({( - 1)}^{n + 1}H_{n + 1}\)\((x) + x{( - 1)}^{n}H_{n}\)\((x) + n\)\({( - 1)}^{n - 1}H_{n - 1}\)\((x)\)\(\rbrack\)\(e^{\frac{- x^{2}}{2}} = 0\).

Даљим сређивањем ове једначине је:

\({( - 1)}^{n - 1}\)\(\lbrack H_{n + 1}(x) - xH_{n}(x) + nH_{n - 1}(x)\rbrack\)\(e^{\frac{- x^{2}}{2}} = 0\),

па се долази до закључка да је:

\(H_{n + 1}(x) - xH_{n}(x) + nH_{n - 1}(x) = 0\). (34)

Ово је рекурентна веза која повезује три узастопна Ермитова полинома.

Из претходних израза, добија се:

\(H_{n}'(x) = nH_{n - 1}(x)\). (35)

Вратимо се сада на једначину (34), њеним диференцирањем добија се да је

\(H_{n + 1}'(x) - xH_{n}'(x) - H_{n}(x) + nH_{n - 1}'(x) = 0\),

уколико у добијену једначину уврстимо (35), добијамо:

\(H_{n}^{''}(x) - xH_{n}'(x) + nH_{n}(x) = 0\).

Дакле, ово је у складу са једначином за физичке полиноме из дефиниције
31. Аналогно се изводи и диференцијална једначина за Ермитов
пробабилистички полином.

Оно што још треба истаћи је да ова једначина одговара и једначини за
ортогоналне полиноме са коефицијентима \(D = B = C = 0,E = - 1,A = 1\),
где је \(\frac{\rho'(x)}{\rho(x)} = - x\).

\begin{itemize}
\item
  \protect\phantomsection\label{_Toc27}{}Ермитов полином
\end{itemize}

\emph{\textbf{Дефиниција 32}} {[}35{]}: Ермитови полиноми су низ
ортогоналних полинома над доменом \(( - \infty,\infty)\) са тежинском
функцијом \(w(x) =\)\(e^{- x^{2}}\) (тада су у ознаци \(H\) и називају
се \textbf{физички}), односно \(w(x) =\)\(e^{- \frac{x^{2}}{2}}\) (у
ознаци \(H_{e}\), називају се \textbf{пробабилистички}).

Конкретно, израз за физички Ермитов полином \(n\)-тог степена гласи:

\(H_{n}(x) =\)\({( - 1)}^{n}\)\(e^{x^{2}}\frac{d^{n}}{dx^{n}}\)\((\)\(e^{- x^{2}}\)\()\)
\(=\)
\(e^{\frac{x^{2}}{2}}\)\({(x - \frac{d}{dx})}^{n}\)\(e^{\frac{- x^{2}}{2}}\).

Формула дата у дефиницији 32 представља Родригову формулу за Ермитов
физички полином. Ред ових полинома је целобројан.

Пратећи дефиницију, за првих неколико Ермитових полинома добија се:

\(H_{0}(x) = 1\),

\(H_{1}(x) = 2x\),

\(H_{2}(x) = 4x^{2} - 2\),

\(H_{3}(x) = 8x^{3} - 12x\),

\(H_{4}(x) = 16x^{4} - 48x^{2} + 12\),

\(H_{5}(x) = 32x^{5} - 160x^{3} + 120x\),

\(H_{6}(x) = 64x^{6} - 480x^{4} + 720x^{2} - 120\).

Када је у питању парност овог полинома, она зависи од реда - паран је за
паран ред, а непарна за непарни ред.

Генератриса Ермитовог физичког полинома је:

\(g(t,x) =\)\(e^{\frac{t}{2}(t - x)}\).

\emph{\textbf{Дефиниција 33:}} Ермитов полином је полином облика:

\[H_{2k}(x) =\]\[{( - 1)}^{k}2^{k}\]\[(2k - 1)!!\]\[\lbrack 1 + \overset{k}{\sum_{j = 1}}\frac{( - 4k)( - 4k + 4)...( - 4k + 4j - 4)}{(2j)!}x^{2j}\rbrack\]

за парно \(n\), односно облика:

\[H_{2k + 1}(x) =\]\[{( - 1)}^{k}2^{k + 1}\]\[(2k + 1)!!\]\[\lbrack x + \overset{k}{\sum_{j = 1}}\frac{( - 4k)( - 4k + 4)...( - 4k + 4j - 4)}{(2j + 1)!}x^{2j + 1}\rbrack\]

за непарно \(n\).

Графички приказано, овај низ полинома изгледа као на слици 8.2.1.

\includegraphics[width=6.59019in,height=3.48568in,alt={pasted-movie.png}]{media/image8.png}

Слика 8.2.1: Ермитови физички полиноми (прва четири).

Једноставнији запис овог полинома, погодан и за рачунарску
имплементацију, је и {[}32{]}:

\(H_{n}(x) =\)\(n!\overset{\lbrack n/2\rbrack}{\sum_{k = 0}}\frac{{( - 1)}^{k}}{k!(n - 2k)!}{(2x)}^{n - 2k}\).

Ермитов физички полином се такође може записати и на следећи начин:

\(H_{n}(z) =\)\(\frac{n!}{2\pi i}\)\(\oint\)\(e^{- t^{2} + 2tz}t^{- n - 1}dt\).

Ермитов полином се може представити следећом детерминантом:

\(H_{n}(x) =\)\(\left| \begin{matrix}
x & n - 1 & 0 & 0 & ... & 0 \\
1 & x & n - 2 & 0 & ... & 0 \\
0 & 1 & x & n - 3 & ... & 0 \\
0 & 1 & 2 & ... & ... & 0 \\
 \vdots & \vdots & \vdots & \vdots & \vdots & \vdots \\
0 & 0 & 0 & ... & 0 & x
\end{matrix} \right|\).

Још једна корисна особина физичког Ермитовог полинома је да се може
одредити помоћу рекурентне формуле:

\(H_{n + 1}(x) = xH_{n}(x) - H_{n}'(x)\),

где је \(n = 0,1,2,...\)

Када се говори о пробабилистичком Ермитовом полиному, првих неколико
полинома гласе:

\(He_{0}(x) = 1\),

\(He_{1}(x) = x\),

\(He_{2}(x) = x^{2} - 1\),

\(He_{3}(x) = x^{3} - 3x\),

\(He_{4}(x) = x^{4} - 6x^{2} + 3\),

\(He_{5}(x) = x^{5} - 10x^{3} + 15x\),

\(He_{6}(x) = x^{6} - 15x^{4} + 45x^{2} - 15\),

\(He_{7}(x) = x^{7} - 21x^{5} + 105x^{3} - 105x\),

\(He_{8}(x) = x^{8} - 28x^{6} + 210x^{4} - 420x^{2} + 105\).

График на коме су приказани ови полиноми налази се на слици 8.2.2.

\includegraphics[width=6.3545in,height=3.51924in,alt={pasted-movie.png}]{media/image9.png}

Слика 8.2.2. Ермитови пробабилистички полиноми (првих четири).

Парност је иста као код Ермитовог физичког полинома.

Генератриса пробабилистичког Ермитовог полинома је:

\(g(t,x) =\)\(e^{t(2x - t)}\).

Између пробабилистичких и физичких Ермитових полинома постоји зависност:

\(H_{n}(x) =\)\(2^{\frac{n}{2}}H_{e_{n}}\)\((\sqrt{2}x)\),

\(H_{e_{n}}\)\((x) =\)\(2^{- \frac{n}{2}}H_{n}\)\((\frac{x}{\sqrt{2}})\).
(36)

За физички Ермитов полином важи:

\(H_{n}( - x) =\)\({( - 1)}^{n}\)\(H_{n}(x)\),

односно парност ове функције зависи од њеног реда. Будући да се
пробабилистички полином може исказати преко физичког на начин дат
формулом (36), иста особина важи и за њега.

Рекурзивна формула којом се физички полином може једноставно израчунати
је

\(H_{n}(x) = 2xH_{n - 1}(x) - 2(n - 1)H_{n - 2}(x),n > 2\),

при чему је \(H_{0}(x) = 1\) и \(H_{1}(x) = 2x\). Ово је погодно за
рачунарску имплементацију, мада може бити проблематично за велике
вредности реда функције, о чему ће бити више речи у делу рада који
конкретно говори о рачунарској имплементацији калкулатора.

Такође, постоје рекурзивне формуле које Ермитови полиноми задовољавају
{[}37{]}:

за пробабилистички полином:

\(H_{e_{n + 1}}\)\((x) = x\)\(H_{e_{n}}\)\((x) -\)\(H_{e_{n}}'\)\((x)\)

\(H_{e_{n}}'\)\((x) = n\)\(H_{e_{n}}\)\((x)\) \footnote{Полиноми који
  задовољавају ову једначину представљају тзв. Апелов низ.}

2) за физички полином:

\(H_{n + 1}\)\((x) = 2xH_{n}\)\((x) -\)\(H_{n}'\)\((x)\)

\(H_{n}'\)\((x) = 2n\)\(H_{n - 1}\)\((x)\) \footnote{Ова једначина је
  такође оно што карактерише Апелов низ.}.

Оно што још може бити интересантна особина физичког Ермитовог полинома
је {[}35{]}:

\(\sum_{k = 0}^{\infty}{( - 1)}^{k}\)\(\frac{H_{2k}}{(2k)!} =\)\(e\cos 2x\),

\(\sum_{k = 0}^{\infty}{( - 1)}^{k}\)\(\frac{H_{2k + 1}}{(2k + 1)!} =\)\(e\sin 2x\).

За велике вредности реда полинома, физички полином се може апроксимирати
изразом:

\(H_{n}(x) \approx\)\(\left\{ \begin{array}{r}
{( - 2)}^{\frac{n}{2}}(n - 1)!!e^{\frac{x^{2}}{2}}\cos(\sqrt{2n + 1}x),nнепарноиn \rightarrow \infty \\
{( - 2)}^{(n - 1)/2}n!!e^{\frac{x^{2}}{2}}\sin(\sqrt{2n + 1}x),nпарноиn \rightarrow \infty
\end{array} \right.\ \).

Ова апроксимација је прецизнија што је апсолутна вредност променљиве
мања, може се искористити искључиво када је \(|x| \leqslant \sqrt{2n}\).
Ова апроксимација имплицира да су нуле полинома

\(r \approx \frac{j\pi}{2\sqrt{2n + 1}}\), где је
\(j = \left\{ \begin{array}{r}
 \pm 1, \pm 2, \pm 5,...непарноn \\
0, \pm 2, \pm 4,...парноn
\end{array} \right.\ \).

\begin{itemize}
\item
  \protect\phantomsection\label{_Toc28}{}Чебишевљеви полиноми
\end{itemize}

Чебишевљеви полиноми представљају низове ортогоналних полинома,
дефинисаних помоћу синусних и косинусних функција. Генерално, синусна и
косинусна функција су такође врсте специјалних функција, настале из
Чебишевљеве диференцијалне једначине, али због њихове огромне примене у
математици спадају у скуп елементарних функција.

Разликујемо Чебишевљеве полиноме прве врсте и Чебишевљеве полиноме друге
врсте. У овој глави ће бити обрађене обе врсте предметних полинома.

\begin{itemize}
\item
  \protect\phantomsection\label{_Toc29}{}Чебишевљева диференцијална
  једначина
\end{itemize}

\emph{\textbf{Дефиниција 34:}} Чебишевљева једначина је хомогена
линеарна диференцијална једначина другог реда дата као:

\((1 - x^{2})y'' - xy' + \alpha^{2}y = 0\),

при чему је \(\alpha\) реална или комплексна константа.

Један од начина да се ова једначина реши је, као код већине претходних
једначина, помоћу степених редова, али будући да је ово решење већ
изложено неколико пута, у овом поглављу ћемо једначину решити на други
начин.

Приметимо да једначина може попримити знатно једноставнији облик уколико
се променљива представи као \(x = \cos\varphi\) \(\Rightarrow\)
\(dx = - \sin\varphi d\varphi\), па је даље
\(\frac{d\varphi}{dx} = - \frac{1}{\sin\varphi}\), и:

\(y' = \frac{dy}{dx} = \frac{dy}{d\varphi}\frac{d\varphi}{dx} = - \frac{1}{\sin\varphi}\frac{dy}{d\varphi}\),

\[\Rightarrow y'' = \frac{d^{2}y}{dx^{2}} = \frac{d}{dx}(\frac{d\varphi}{dx}) = \frac{d}{d\varphi}\frac{d\varphi}{dx}(\frac{dy}{dx}) = - \frac{1}{\sin\varphi}\frac{d}{d\varphi}( - \frac{1}{\sin\varphi}\frac{dy}{d\varphi}) =\]

\(= \frac{1}{\sin\varphi}\lbrack\frac{d}{d\varphi}\left( \frac{1}{\sin\varphi} \right)\frac{dy}{d\varphi} + \frac{1}{\sin\varphi}\frac{d^{2}y}{dt^{2}}\rbrack\)\(= \frac{1}{\sin\varphi}\lbrack( - \frac{\cos\varphi}{\sin^{2}\varphi})\frac{dy}{d\varphi} + \frac{1}{\sin\varphi}\frac{d^{2}y}{d\varphi^{2}}\rbrack\)\(= \frac{1}{\sin^{2}\varphi}\lbrack( - \frac{\cos\varphi}{\sin\varphi})\frac{dy}{d\varphi} + \frac{d^{2}y}{d\varphi^{2}}\rbrack\).

Заменом израза за изводе у деференцијалну једначину добија се:

\((1 - \cos^{2}\varphi)\)\(\frac{1}{\sin^{2}\varphi}\lbrack( - \frac{\cos\varphi}{\sin\varphi})\frac{dy}{d\varphi} + \frac{d^{2}y}{d\varphi^{2}}\rbrack\)\(- \cos\varphi\lbrack - \frac{1}{\sin\varphi}\frac{d\varphi}{dt}\rbrack + n^{2}y = 0\),

\(\frac{\sin^{2}\varphi}{\sin^{2}\varphi}\lbrack - \frac{\cos\varphi}{\sin\varphi}\frac{dy}{d\varphi} + \frac{d^{2}y}{d\varphi^{2}}\rbrack + \frac{\cos\varphi}{\sin\varphi}\frac{dy}{d\varphi} + n^{2}y = 0\),

\(- \boxed{\frac{\cos\varphi}{\sin\varphi}\frac{dy}{d\varphi}} + \frac{d^{2}y}{d\varphi^{2}} + \boxed{\frac{\cos\varphi}{\sin\varphi}\frac{dy}{d\varphi}} + n^{2}y = 0\).

Дакле, сада једначина постаје:
\(\frac{d^{2}y}{d\varphi^{2}} + n^{2}y = 0\). Ово је једначина линеарог
хармонијског осцилатора.

Њено опште решење може се записати као:

\(y(\varphi) = C_{1}\cos n\varphi + C_{2}\sin n\varphi\),

а такође и у облику

\(y(\varphi) = C\cos(n\varphi + \alpha)\),

при чему су \(C,C_{1},C_{2},\alpha\) реални бројеви.

Ради поједностављења може се одабрати \(\alpha = 0\), и тада опште
решење Чебишевљеве једначине постаје функција:

\(y(x) = C\cos(n\arccos x)\).

Константа \(n\) може бити реалан број, али у случају када је целобројна
константа, дата функција се назива Чебишевљев полином прве врсте.

\begin{itemize}
\item
  \protect\phantomsection\label{_Toc30}{}Чебишевљеви полиноми прве врсте
\end{itemize}

\emph{\textbf{Дефиниција 35}} {[}35{]}\emph{\textbf{:}} Чебишевљеви
полиноми \emph{прве врсте,} у ознаци \(T_{n}\), су ортогонални полиноми,
чија је тежинска фукција \(w(x) =\)\({(1 - x^{2})}^{- 1/2}\) над
интервалом \(\lbrack - 1,1\rbrack\), дати изразом:

\(T_{n}(x) =\)\(\cos(n\arccos x),|x| \leq 1,n \in {\mathbb{N}}_{0}\),
односно

\(T_{n}(\cos\varphi) = \cos(n\varphi)\).

\emph{\textbf{Дефиниција 36}} {[}35{]}\emph{\textbf{:}} Чебишевљеви
полиноми прве врсте изражени Родриговом формулом имају облик:

\(T_{n}(x) =\)\(\frac{{( - 2)}^{n}n!}{(2n)!}\sqrt{1 - x^{2}}\frac{d^{n}}{dx^{n}}\)\({(1 - x^{2})}^{n - 1/2}\),
односно

\(T_{n}(x) =\)\(\frac{n}{2}\overset{\lbrack n/2\rbrack}{\sum_{m = 0}}\)\({( - 1)}^{m}\)\(\frac{(n - m - 1)!}{m!(n - 2m)!}\)\({(2x)}^{n - 2m}\).

Осим на дате начине, могу бити изражени и помоћу редова, контурних
интеграла, итд., а интересантан је њихова матрична репрезентација преко
тродијагоналне детерминанте:

\(T_{n}\overset{}{(x)} = \left| \begin{array}{r}
\begin{matrix}
x & 1 & 0 & ... & 0 & 0 \\
1 & 2x & 1 & \ddots & 0 & 0 \\
0 & 1 & 2x & \ddots & 0 & \vdots \\
0 & 0 & 1 & \ddots & 0 & 0 \\
 \vdots & \ddots & \ddots & \ddots & \ddots & 1 \\
0 & 0 & 0 & \cdots & 1 & 2x
\end{matrix}
\end{array} \right|\).

Првих пет Чебишевљевих полинома прве врсте су:

\(T_{0}(x) = 1\)

\(T_{1}(x) = x\)

\(T_{2}(x) = 2x^{2} - 1\)

\(T_{3}(x) = 4x^{3} - 3x\)

\(T_{4}(x) = 8x^{4} - 8x^{2} + 1\).

На слици 9.2.1 може се видети и график Чебишевљевог полинома прве врсте
за првих неколико вредности њиховог
реда.\includegraphics[width=5.88847in,height=3.93019in,alt={pasted-movie.png}]{media/image10.png}

Слика 9.2.1: Графици првих пет Чебишевљевих полинома прве врсте (лево).

За интервал у коме смо дефинсали полином важи да
\(- 1 \leq T_{n}(x) \leq 1\). Као што се са графика примећује, полином
реда \(n\) има \(n\) нула на интервалу \(\lbrack - 1,1\rbrack\) и тачно
\(n + 1\) локални екстремум у истом интервалу. Локални екстремум се
дешава када је \(x = \cos\frac{k\pi}{n},k = 1,2,...n\) и тада функција
узима вредност \(\left| T_{n}(x) \right| = 1\).

Парност Чебишевљевог полинома прве врсте зависи од његовог степена - за
паран степен и фунција је парна, односно за непаран је непарна:
\(T_{n}( - x) = {( - 1)}^{n}T_{n}(x)\).

Генератриса полинома прве врсте је:
\(g(x,t) =\)\(\overset{\infty}{\sum_{n = 0}}\)\(T_{n}(x)t^{n} =\)\(\frac{1 - xt}{1 - 2xt + t^{2}}\).

Чебишевљеви полиноми су специјална врста Јакобијевих полинома
\(P_{n}^{(\alpha,\beta)}\): прве врсте за случај
\(\alpha = \beta = - \frac{1}{2}\), а друге врсте за случај
\(\alpha = \beta = \frac{1}{2}\).

Неке специфичне вредности тих полинома прве врсте:

\(T_{n}(1) = 1\),

\(T_{n}( - 1) = {( - 1)}^{n}\),

\(T_{2n}(0) = {( - 1)}^{n}\),

\(T_{2n + 1}(0) = 0\).

За полином прве врсте, алгебарска репрезентација би била:

\(T_{n}(x) =\)\(\frac{1}{2}\lbrack{(x + \sqrt{x^{2} - 1})}^{n} + {(x - \sqrt{x^{2} - 1})}^{n}\rbrack\).

Такође, постоји и развој попут биномног развоја {[}35{]}:

\(T_{n}(x) =\)\(\left( \begin{array}{r}
n \\
0
\end{array} \right)x^{n} - \left( \begin{array}{r}
n \\
2
\end{array} \right)x^{n - 2}{(1 - x)}^{2} + \left( \begin{array}{r}
n \\
4
\end{array} \right)x^{n - 4}{(1 - x^{2})}^{2} - ...\).

Када је ред ових полинома целобројни ненегативан, производ ових полинома
се може израчунати по формули:
\(T_{m}(x) \cdot T_{n}(x) =\)\(\frac{1}{2}\)\(\lbrack\)\(T_{m + n}\)\((x) +\)\(T_{|m - n|}\)\((x)\)\(\rbrack\).

Такође за \(m,n \geq 0\) важе и формуле {[}39{]}:

\((\)\(T_{m + n}\)\((x) - 1)(\)\(T_{|m - n|}\)\((x) - 1) = (\)\(T_{m}(x) - T_{n}(x)\)\()^{2}\),

\(T_{m}(T_{n}(x)) = T_{mn}(x)\).

Слично претходним специјалним функцијама, и Чебишевљеви полиноми
задовољавају одређене рекурентне релације. У случају Чебишевљевих
полинома прве врсте, то су следеће формуле {[}19{]}:

\(T_{n + 1}(x) =\)\(2xT_{n}(x) - T_{n - 1}(x)\),

\(T_{n + 1}(x) = xT_{n}(x) -\)\(\sqrt{(1 - x^{2})(1 - (T_{n}{(x))}^{2})}\),

\((x - 1)\)\(\lbrack T_{2n + 1}(x) - 1\rbrack =\)\({\lbrack T_{n + 1}(x) - T_{n}(x)\rbrack}^{2},n \geqslant 1\),

\(2(x^{2} - 1)\)\(\lbrack T_{2n}(x) - 1\rbrack\)\(= {\lbrack T_{n + 1}(x) - T_{n - 1}(x)\rbrack}^{2},n \geqslant 1\).

\(T_{2n}(x) = \left( T_{n}(x) \right)^{2} - 1\).

Изводи и интеграли ових полинома такође задовољавају бројне рекурзивне
формуле, неке од њих су {[}10{]}{[}39{]}{[}35{]}:

\(2T_{n}(x) = T_{n + 1}'(x) - 2xT_{n}'(x) + T_{n - 1}'(x)\),

\((1 - x^{2})T_{n}'(x) = - nxT_{n}(x) + nT_{n - 1}(x)\),

\(T_{n}^{''}(x) =\)\(\frac{x}{1 - x^{2}}\)\(T_{n}'(x) -\)\(\frac{n^{2}}{1 - x^{2}}\)\(T_{n}(x)\),

\(\int\)\(T_{n}(x)dx\)\(= \frac{1}{2}\)\((\frac{T_{n + 1}(x)}{n + 1} - \frac{T_{n - 1}(x)}{n - 1}) + C,n \geq 2\),

\(\int_{0}^{x}\)\(T_{n}(t)dt =\)\(\frac{n}{n^{2} - 1}\lbrack\)\(T_{n + 1}(x) +\)\(\sin\frac{n\pi}{2}\rbrack - \frac{x}{n - 1}\)\(T_{n}(x),\)
\(n = 2,3,4,...\).

\begin{itemize}
\item
  \protect\phantomsection\label{_Toc31}{}Чебишевљеви полиноми друге
  врсте
\end{itemize}

\emph{\textbf{Дефиниција 37}} {[}12{]}\emph{\textbf{:}} Чебишевљеви
полиноми \emph{друге врсте,} у ознаци \(U_{n}\), су ортогонални
полиноми, чија је тежинска функција
\(w\overset{}{(x)} = {(1 - x^{2})}^{\frac{1}{2}}\) над интервалом
\(\lbrack - 1,1\rbrack\), дати изразом:

\(U_{n}(x) =\)\(\frac{\sin((n + 1)x)}{\sin x}\),

односно

\(U_{n}(\cos\varphi)\sin\varphi = \sin((n + 1)\varphi)\).

\emph{\textbf{Дефиниција 38:}} Чебишевљеви полиноми прве врсте изражени
Родриговом формулом имају облик:\({(1 - x^{2})}^{\frac{1}{2}}\)

\(U_{n}(x) =\)\(\frac{{( - 1)}^{n}(n + 1)\sqrt{\pi}}{2^{n + 1}(n + \frac{1}{2})!{(1 - x^{2})}^{\frac{1}{2}}}\)
\(\frac{d^{n}}{dx^{n}}\) \(\lbrack{(1 - x^{2})}^{n + 1/2}\rbrack\),
односно

\(U_{n}(x) =\)\(\frac{{( - 1)}^{n}(n + 1)\sqrt{\pi}}{2^{n + 1}\Gamma(n + \frac{3}{2})\sqrt{1 - x^{2}}}\frac{d^{n}}{dx^{n}}\lbrack{(1 - x^{2})}^{n + \frac{1}{2}}\rbrack\),
тј.

\(U_{n}(x) =\)\(\overset{\lbrack n/2\rbrack}{\sum_{m = 0}}\)\({( - 1)}^{m}\)\(\frac{(n - m)!}{m!(n - 2m)!}\)\({(2x)}^{n - 2m}\).

{\def\LTcaptype{none} % do not increment counter
\begin{longtable}[]{@{}
  >{\centering\arraybackslash}p{(\linewidth - 2\tabcolsep) * \real{0.0597}}
  >{\centering\arraybackslash}p{(\linewidth - 2\tabcolsep) * \real{0.6073}}@{}}
\toprule\noalign{}
\begin{minipage}[b]{\linewidth}\centering
\[n\]
\end{minipage} & \begin{minipage}[b]{\linewidth}\centering
\[U_{n}\]
\end{minipage} \\
\midrule\noalign{}
\endhead
\bottomrule\noalign{}
\endlastfoot
\emph{0} & \(1\) \\
\emph{1} & \(2x\) \\
\emph{2} & \(4x^{2} - 1\) \\
\emph{3} & \(8x^{3} - 4x\) \\
\emph{4} & \(16x^{4} - 12x^{2} + 1\) \\
\emph{5} & \(32x^{5} - 32x^{3} + 6x\) \\
\end{longtable}
}

Табела 9.3.1: Првих пет Чебишевљевих полинома друге врсте.

На слици 9.3.1. можемо видети како изгледа првих пет Чебишевљевих
полинома друге врсте. Као и код полинома прве врсте, парност зависи од
степена, \(U_{n}( - x) = {( - 1)}^{n}U_{n}(x)\). Нуле полинома су
\(x_{k} = \cos(\frac{2k + 1}{2n + 2}\pi),k = 0,1,...,n\). Полином
\(n\)-тог степена достиже екстремуме у нулама полинома степена
\(n - 1\).

\includegraphics[width=5.66397in,height=4.02596in,alt={pasted-movie.png}]{media/image11.png}

Слика 9.3.1: Првих пет Чебишевљевих полинома друге врсте.

Генератриса Чебишевљевог полинома друге врсте гласи:

\(g(x,t) =\)\(\overset{\infty}{\sum_{n = 0}}\)\(U_{n}(x)t^{n} =\)\(\frac{1}{1 - 2xt + t^{2}}\).

Овај полином се такође може записати и као детерминанта:

\(U_{n}\overset{}{(x)} = \left| \begin{array}{r}
\begin{matrix}
2x & 1 & 0 & ... & 0 & 0 \\
1 & 2x & 1 & \ddots & 0 & 0 \\
0 & 1 & 2x & \ddots & 0 & \vdots \\
0 & 0 & 1 & \ddots & 0 & 0 \\
 \vdots & \ddots & \ddots & \ddots & \ddots & 1 \\
0 & 0 & 0 & \cdots & 1 & 2x
\end{matrix}
\end{array} \right|\).

Једна од репрезентација Чебишевљевог полинома друге врсте је и {[}39{]}:

\(U_{n}(x) =\)
\(\frac{{(x + \sqrt{x^{2} - 1})}^{n + 1} - {(x - \sqrt{x^{2} - 1})}^{n + 1}}{2\sqrt{x^{2} - 1}}\).

Као и код Чебишевљевог полинома прве врсте, постоји и репрезентација
помоћу биномног развоја {[}35{]}:

\(U_{n - 1}(x) =\)\(\left( \begin{array}{r}
n \\
1
\end{array} \right)x^{n - 1} - \left( \begin{array}{r}
n \\
2
\end{array} \right)x^{n - 3}(1 - x^{2}) + \left( \begin{array}{r}
n \\
3
\end{array} \right)x^{n - 5}{(1 - x^{2})}^{2} - ...\).

Неке специфичне вредности полинома друге врсте су {[}10{]}:

\(U_{n}(1) = n + 1\),

\(U_{n}( - 1) = {( - 1)}^{n}(n + 1)\),

\(U_{2n}(0) = {( - 1)}^{n}\),

\(U_{2n + 1}(0) = 0\).

За Чебишевљеве полиноме друге врсте важe наредне рекурзивне везе:

\(U_{n + 2}(x) = 2xU_{n + 1}(x) - U_{n}(x),n \geq 0\),

\((1 - x^{2})U_{n}'(x) = - nxU_{n}(x) + (n + 1)U_{n - 1}(x)\),

\({(1 - x)}^{2}\)\(U_{n}^{''}(x) - 3xU_{n}'(x) +\)\(n(n + 2)\)\(U_{n}(x) = 0\).

\begin{itemize}
\item
  \protect\phantomsection\label{_Toc32}{}Везе између Чебишевљевих
  полинома прве и друге врсте
\end{itemize}

Постоји велики број релација које повезују Чебишевљеве полиноме прве и
друге врсте, а неке од тих релација ће бити дате у наставку:

\(T_{n}(x) = U_{n}(x) - xU_{n - 1}(x)\), као и

\(T_{n}(x) =\)\(\frac{1}{2}\)\((U_{n}(x) - U_{n - 2}(x))\),

\(T_{n - 1}(x) - T_{n + 1}(x) = 2(1 - x^{2})U_{n - 1}(x),n \geq 1\),

\(U_{n}(x) = 2\)\(\overset{\left\lbrack \frac{n - 1}{2} \right\rbrack}{\sum_{k = 0}}\)\(T_{2k + 1}(x)\),

\(U_{n}(x) = 2\)\(\overset{\left\lbrack \frac{n}{2} \right\rbrack}{\underset{k = 0}{\sum}}\)\(T_{2k}(x) - 1\),

\(T_{n + 2}(x) = 1 - 2(1 - x^{2})\)\(\overset{\left\lbrack \frac{n}{2} \right\rbrack}{\sum_{k = 0}}\)\(U_{2k}(x)\),

\(T_{n + 2}(x) = x - 2(1 - x^{2})\)\(\overset{\left\lbrack \frac{n - 1}{2} \right\rbrack}{\sum_{k = 0}}\)\(U_{2k + 1}(x)\),

\(U_{mn - 1}(x) = U_{m - 1}(T_{n}(x)) \cdot U_{n - 1}(x)\)\(= U_{mn - 1}(x)\),

\(U_{n}(x) =\)\(\overset{n}{\sum_{j = 0}}\)\(x^{j}\)\(T_{n - j}\)\((x)\),
\(T_{n}(x) =\)\(\frac{1}{n + 1}\overset{n}{\sum_{j = 0}}\)\((2T_{j}\)\((x) -\)\(x^{j})T_{n - j}\)\((x)\),

\(2\overset{n}{\sum_{j = 0}}\)\(T_{j}\)\((x)\)
\(T_{n - j}\)\((x)\)\(= U_{n}\)\((x) + (n + 1)\)\(T_{n}\)\((x)\),

\(T_{n}^{2}(x) - (x^{2} - 1)U_{n - 1}^{2}(x) = 1\),

\(\sqrt{(1 - x^{2})}\)\(U_{n - 1}(x) = xT_{n}(x) - T_{n + 1}(x)\),

\(U_{n}(x) =\)\(= 2\overset{\left\lbrack \frac{n}{2} \right\rbrack}{\sum_{j = 0}}T_{n - 2j}\)\((x) -\)\(\frac{{( - 1)}^{n} + 1}{2},n \geq 0\),

\(\frac{U_{n}(x) + U_{n - 1}(x)}{2} = \frac{1}{2} +\)\(T_{1}(x) + .... + T_{n}(x)\),

\(U_{n - 1}(x) =\)\(\frac{1}{n}\)\(T_{n}'(x) =\)\(\frac{\sin n\theta}{\sin\theta},x = \cos\theta\),
{[}39{]}

\(T_{n}'(x)\)=\(\frac{n}{\sqrt{1 - x^{2}}}\)\(U_{n}(x)\),

\(\overset{n}{\sum_{r = 0}}T_{2r}\)\((x) =\)\(\frac{1}{2} + \frac{U_{2n + 1}(x)}{2\sqrt{1 - x^{2}}}\).{[}10{]}

Формуле везане за производ два Чебишевљева полинома су:

\(U_{n}(x)U_{m}(x) =\)\(\frac{T_{n - m}(x) - T_{n + m + 2}(x)}{2(1 - x^{2})}\),
\(n \geq m\),

\(T_{n}(x)U_{m}(x) =\)\(\left\{ \begin{array}{r}
\frac{1}{2}U_{n + m}\overset{}{(x)} + \frac{1}{2}U_{n - m}\overset{}{(x)},n \leq m \\
\frac{1}{2}U_{n + m}\overset{}{(x)},n = m + 1 \\
\frac{1}{2}U_{n + m}\overset{}{(x)} - \frac{1}{2}U_{n - m - 2}\overset{}{(x)},n \geq m + 2
\end{array} \right.\ \),

Значајни интеграли Чебишевљевих полинома су {[}35{]}:

\(U_{n}(x) =\)\(\frac{1}{\pi}\int_{- 1}^{1}\frac{T_{n + 1}\overset{}{(a)}}{(a - x)\sqrt{1 - a^{2}}}da\),
\(T_{n}(x) =\)\(\frac{1}{\pi}\int_{- 1}^{1}\frac{\sqrt{1 - a^{2}}U_{n - 1}\overset{}{(a)}}{x - a}da\),

\(\int_{0}^{x}\)\(U_{n}(t)dt =\)\(\frac{T_{n + 1}\overset{}{(x)} + \sin\frac{n\pi}{2}}{n + 1},\)
\(n = 0,1,2,...\),

\(\int_{1}^{x}\)\(U_{n}(t)dt =\)\(\frac{1 - T_{n + 1}\overset{}{(x)}}{n + 1},\)
\(n = 0,1,2,...\).

\(\frac{1}{\pi}\int_{- 1}^{1}\frac{T_{n}\overset{}{(t)}}{\sqrt{1 - t^{2}}}\frac{dt}{t - y} =\)\(U_{n - 1}(y),\)
\(- 1 < y < 1\),

\(\frac{1}{\pi}\int_{- 1}^{1}\)\(U_{n}(t)\)\(\sqrt{1 - t^{2}}\frac{dt}{t - y} =\)\(T_{n + 1}(y),\)
\(- 1 < y < 1\).

\protect\phantomsection\label{_Toc33}{}Јакобијеви полиноми

\begin{itemize}
\item
  \protect\phantomsection\label{_Toc34}{}Диференцијална једначина
\end{itemize}

\emph{\textbf{Дефиниција 39:}} Диференцијална једначина чијим се
решавањем добија Јакобијев полином као партикуларно решење је хомогена
линеарна диференцијална једначина другог реда:

\((1 - x^{2})y'' + (\beta - \alpha - (\alpha + \beta + 2)x)y' + n(n + \alpha + \beta + 1)y = 0\),

односно

\(\frac{d}{dx}\lbrack{(1 - x)}^{\alpha + 1}{(1 + x)}^{\beta + 1}\gamma'\rbrack +\)\(n(n + \alpha + \beta + 1){(1 - x)}^{\alpha}{(1 + x)}^{\beta}y = 0\).
(37)

Ова једначина је позната као сингуларна Штурм-Луивилова једначина.

Приметимо да када је \(\alpha = \beta = 0\) једначина се своди на
диференцијалну једначину Лежандрових полинома (33).

Заменом релације
\(y = \overset{\infty}{\sum_{\nu = 0}}a_{\nu}{(x - 1)}^{\nu}\) у
претходно наведену једначину, добија се рекурзивна релација

\[\lbrack\gamma - \nu(\nu + \alpha + \beta + 1)\rbrack a_{\nu} - 2(\nu + 1)(\nu + \alpha + 1)a_{\nu + 1} = 0\]

за \(\nu = 0,1,...\), где \(\nu = n(n + \alpha + \beta + 1)\). Решавањем
ове релације добија се Јакобијев полином (38).

Једначина у облику(37) може бити трансформисана у

\(\frac{d^{2}u}{dx^{2}} + \lbrack\frac{1}{4}\frac{1 - \alpha^{2}}{{(1 - x)}^{2}} + \frac{1}{4}\frac{1 - \beta^{2}}{{(1 + x)}^{2}} + \frac{n(n + \alpha + \beta + 1) + \frac{1}{2}(\alpha + 1)(\beta + 1)}{1 - x^{2}}\rbrack u = 0\),

где је
\(u(x) =\)\({(1 - x)}^{(\alpha + 1)/2}{(1 + x)}^{(\beta + 1)/2}P_{n}^{(\alpha,\beta)}\)\((x)\)
и

\[\frac{d^{2}u}{d\theta^{2}} + \lbrack\frac{\frac{1}{4} - \alpha^{2}}{4\sin^{2}(\frac{\theta}{2})} + \frac{\frac{1}{4} - \beta^{2}}{4\cos^{2}(\frac{\theta}{2})} + {(n + \frac{\alpha + \beta + 1}{2})}^{2}\rbrack u = 0\]

а
\(u(\theta) =\)\(\sin^{\alpha + 1/2}(\frac{\theta}{2})\cos^{\beta + 1/2}(\frac{\theta}{2})P_{n}^{(\alpha,\beta)}\)\((\cos\theta)\).

Постоје још два облика једначине чија су решења Јакобијеви полиноми,
назване Јакобијеве једначине:

\(x(1 - x)y^{}'' + \lbrack\gamma - (\alpha + 1)x\rbrack y' + n(\alpha + n)y = 0\),
и

\((a_{1} + b_{1}x + c_{1}y)(xy' - y) - (a_{2} + b_{2}x + c_{2}y)y' + (a_{3} + b_{3}x + c_{3}y) = 0\).

\begin{itemize}
\item
  \protect\phantomsection\label{_Toc35}{}Јакобијеви полиноми
\end{itemize}

\emph{\textbf{Дефиниција 40:}} Јакобијеви полиноми, са параметрима
\(\alpha > - 1,\beta > - 1\), су низови полинома
\(P_{n}^{\alpha,\beta}\)\((x),n \in {\mathbb{N}}_{0}\),
\(x \in {\mathbb{R}}\), које имају следећи облик:

\({P_{n}}^{(\alpha,\beta)}\)\((x) =\)\(\frac{{( - 1)}^{n}}{2^{n}n!}{(1 - x)}^{- \alpha}{(1 + x)}^{- \beta}\frac{d^{n}}{dx^{n}}\left\lbrack {(1 - x)}^{\alpha + n}{(1 + x)}^{\beta + n} \right\rbrack\).
(38)

Домен ортогоналности ове класе полинома је коначан интервал, за који се
без губитка општости може узети (-1, 1).

Тежинска функција у овом случају би била {[}38{]}:

\(\rho(x) =\)\({(1 - x)}^{\alpha}{(1 + x)}^{\beta}\).

Јакобијеви полиноми су низ полинома који задовољава следећу релацију
ортогоналности:

\(\int_{- 1}^{1}{(1 - x)}^{\alpha}{(1 + x)}^{\beta}P_{n}^{(\alpha,\beta)}\)\((x)\)\(P_{m}^{(\alpha,\beta)}\)\((x)dx =\)\(\left\{ \begin{array}{r}
0,m \neq n \\
h_{n},m = n
\end{array} \right.\ \),

при чему је

\(h_{n} = \frac{2^{\alpha + \beta + 1}\Gamma(\alpha + n + 1)\Gamma(\beta + n + 1)}{(2n + \alpha + \beta + 1)n!(n + \alpha + \beta + 1)}\).

Слично једначинама обрађеним у претходним поглављима, и Јакобијеве
једначине могу бити добијене из Родригове формуле као:

\({(1 - x)}^{\alpha}{(1 + x)}^{\beta}P_{n}^{(\alpha,\beta)}\)\((x) =\)\(\frac{{( - 1)}^{n}d^{n}}{2^{n}n!dx^{n}}\lbrack{(1 - x)}^{\alpha + n}{(1 + x)}^{\beta + n}\rbrack\).

Даље произилази:

\(\frac{d}{dx}(P_{n}^{(\alpha,\beta)}\)\((x)\)\() = \frac{1}{2}(n + \alpha + \beta + 1)P_{n - 1}^{(\alpha + 1,\beta + 1)}\)\((x)\).

Применом Лајблицовог правила\footnote{Лајбницово правило гласи: ако
  постоје n-ти изводи функција \(f\) и \(g\), тада је њихов производ
  функција чији \emph{n}-ти извод постоји и износи:
  \({(f \cdot g)}^{(n)} = \overset{n}{\sum_{k = 0}}\left( \begin{array}{r}
  n \\
  k
  \end{array} \right)f^{(k)}g^{(n - k)}\), при чему је \(f^{(k)}\)
  \emph{k}-ти извод функције \(f\)} на Родригову формулу може се добити
експлицитни израз:

\[{( - 2)}^{n} \cdot n! \cdot {(1 - x)}^{\alpha}{(1 - x)}^{\beta} \cdot P_{n}^{(\alpha,\beta)}\]\[(x) =\]\[\overset{n}{\sum_{k = 0}}\left( \begin{array}{r}
n \\
k
\end{array} \right)\frac{d^{n - k}}{dx^{n - k}}{(1 - x)}^{n + \alpha}\frac{d^{k}}{dx^{k}}{(1 + x)}^{n + \beta} =\]

\({( - 1)}^{n}{(1 - x)}^{\alpha}{(1 + x)}^{\beta}n!\overset{n}{\sum_{k = 0}}\left( \begin{array}{r}
n + \alpha \\
n - k
\end{array} \right)\left( \begin{array}{r}
n + \beta \\
k
\end{array} \right){(x - 1)}^{k}{(x + 1)}^{n - k}\),

што даље даје:

\(P_{n}^{(\alpha,\beta)}\)\((x) =\)\(\frac{1}{2^{n}}\overset{n}{\sum_{k = 0}}\left( \begin{array}{r}
n + \alpha \\
k
\end{array} \right)\left( \begin{array}{r}
n + \beta \\
n - k
\end{array} \right){(x - 1)}^{n - k}{(1 + x)}^{k}\),
\(n \in {\mathbb{N}}_{0}\).

Дакле, Јакобијеви полиноми су парни или непарни у зависности од реда
полинома.

Вредности првих неколико Јакобијевих полинома су {[}20{]}:

\(P_{0}^{(\alpha,\beta)}\)\((x) = 1\),

\(P_{1}^{(\alpha,\beta)}\)\((x) =\)\(\frac{1}{2}\)\((\alpha - \beta + (\alpha + \beta + 2)x)\),

\(P_{2}^{(\alpha,\beta)}\overset{}{(x)} =\)\(\frac{1}{8}\)\(\lbrack 4(\alpha + 1)(\alpha + 2) + 4(\alpha + \beta + 3)(\alpha + 2)(x - 1) + (\alpha + \beta + 3)(\alpha + \beta + 4)\)\({(x - 1)}^{2}\)\(\rbrack\).

Генератриса Јакобијевих полинома је:

\(g(t,x) =\)\({(1 - 2xt + t^{2})}^{- \alpha} \cdot {(1 + 2xt + t^{2})}^{- \beta}\).

Јакобијеви полиноми су нормализовани тако да је:
\(P_{n}^{(\alpha,\beta)}\)\((1) =\)\(\left( \begin{array}{r}
n + \alpha \\
n
\end{array} \right) = \frac{\Gamma(n + \alpha + 1)}{n!\Gamma(\alpha + 1)}\).

На слици 10.2.1 може се видети график Јакобијевих полинома за различите
вредности параметара и реда. Ови полиноми имају нуле на интервалу
\(( - 1,1)\).\includegraphics[width=6.69298in,height=2.7795in,alt={pasted-movie.png}]{media/image12.png}

Слика 10.2.1: Првих неколико Јакобијевих полинома за различите вредности
параметара.

Као значајна последица симетричности тежинске функције Јакобијевих
полинома и ортогоналности јавља се симетрична релација {[}52{]}:

\(P_{n}^{(\alpha,\beta)}\)\(( - x) =\)\({( - 1)}^{n}P_{n}^{(\beta,\alpha)}\)\((x)\).

Из претходне две релације следи и:
\(P_{n}^{(\alpha,\beta)}\)\(( - 1) =\)\({( - 1)}^{n}\)\(\left( \begin{array}{r}
n + \beta \\
n
\end{array} \right)\).

Јакобијеви полиноми задовољавају рекурзивне релацијe:

\[2n(n + \alpha + \beta)(2n + \alpha + \beta - 2)\]\[P_{n}^{(\alpha,\beta)}\]\[(x) =\]

\((2n + \alpha + \beta - 1)\)\(\{(2n + \alpha + \beta)(2n + \alpha + \beta - 2)x + \alpha^{2} - \beta^{2}\}\)\(P_{n - 1}^{(\alpha,\beta)}\)\((x) -\)\(2(n + \alpha - 1)(n + \beta - 1)(2n + \alpha + \beta)\)\(P_{n - 2}^{(\alpha,\beta)}\)\((x)\),
\(n \geq 2\),

\((\alpha + \beta + 2n)\)\(P_{n}^{(\alpha,\beta - 1)}\)\((x) = (\alpha + \beta + n)\)\(P_{n}^{(\alpha,\beta)}\)\((x) + (\alpha + n)\)\(P_{n - 1}^{(\alpha,\beta)}\)\((x)\),

\((\alpha + \beta + 2n)\)\(P_{n}^{(\alpha - 1,\beta)}\)\((x) = (\alpha + \beta + n)\)\(P_{n}^{(\alpha,\beta)}\)\((x) - (\beta + n)\)\(P_{n - 1}^{(\alpha,\beta)}\)\((x)\).

Изводи Јакобијевих полинома такође задовољавају рекурзивне формуле:

\(P_{n}^{(\alpha,\beta)'}\)\((x) =\)\(\frac{1}{2}\)\((1 + \alpha + \beta + n)P_{n - 1}^{(\alpha + 1,\beta + 1)}\)\((x)\),

\((x + 1)\)\(P_{n}^{(\alpha,\beta)'}\)\((x) = n\)\(P_{n}^{(\alpha,\beta)}\)\((x) + (\beta + n)\)\(P_{n - 1}^{(\alpha,\beta + 1)}\)\((x)\),

\((x + 1)\)\(P_{n}^{(\alpha,\beta)'}\)\((x) = n\)\(P_{n}^{(\alpha,\beta)}\)\((x) - (\alpha + n)\)\(P_{n - 1}^{(\alpha,\beta + 1)}\)\((x)\),

\(P_{n}^{(\alpha,\beta)'}\)\((x) =\)\(\frac{1}{2}\)\((\overset{}{(\beta + n)}P_{n}^{(\alpha + 1,\beta)}\overset{}{(x)} + \overset{}{(\alpha + n)}P_{n - 1}^{(\alpha,\beta + 1)}\overset{}{(x)})\).

\begin{itemize}
\item
  \protect\phantomsection\label{_Toc36}{}Преглед основних особина
  специјалних функција
\end{itemize}

У овој глави дат је табеларни преглед најважнијих особина специјалних
фунција обрађених у овом раду (табеле 11.1-11.8): диференцијална
једначина чије је партикуларно решење дата специјална функција, њена
(најчешће коришћена) формула, тежинска функција и домен ортогоналности,
генератриса, Родригова формула, рекурзивне формуле, као и формула
записана помоћу контурног интеграла.

{\def\LTcaptype{none} % do not increment counter
\begin{longtable}[]{@{}
  >{\raggedright\arraybackslash}p{(\linewidth - 2\tabcolsep) * \real{0.3082}}
  >{\centering\arraybackslash}p{(\linewidth - 2\tabcolsep) * \real{0.6918}}@{}}
\toprule\noalign{}
\begin{minipage}[b]{\linewidth}\raggedright
Особина
\end{minipage} & \begin{minipage}[b]{\linewidth}\centering
\end{minipage} \\
\midrule\noalign{}
\endhead
\bottomrule\noalign{}
\endlastfoot
Диференцијална једначина &
\(x^{2}y'' + xy' + \left( x^{2} - \alpha^{2} \right)y = 0\) \\
Формула &
\(J_{n}(x) =\)\(\overset{\infty}{\sum_{m = 0}}\frac{{( - 1)}^{m}}{(m + n)!m!}\left( \frac{x}{2} \right)^{2m + n}\) \\
Домен & \(x \in {\mathbb{R}},n \in {\mathbb{N}}_{0}\) \\
Генератриса & \(g(x,t) = e^{\frac{x}{2}(t - \frac{1}{t})}\) \\
Значајне рекурзивне формуле &
\(J_{n + 1}\overset{}{(x)} = \frac{2n}{x}J_{n}\overset{}{(x)} - J_{n - 1}\overset{}{(x)},n \in {\mathbb{N}}\) \\
Контурна формула &
\(J_{n}\overset{}{(z)} = \frac{1}{2\pi i}\oint e^{\frac{z}{2}\left( t - \frac{1}{t} \right)}t^{- n - 1}dt.\) \\
\end{longtable}
}

Табела 11.1: Преглед најважнијих особина Беселових функције прве врсте.

{\def\LTcaptype{none} % do not increment counter
\begin{longtable}[]{@{}
  >{\raggedright\arraybackslash}p{(\linewidth - 2\tabcolsep) * \real{0.3375}}
  >{\centering\arraybackslash}p{(\linewidth - 2\tabcolsep) * \real{0.6625}}@{}}
\toprule\noalign{}
\begin{minipage}[b]{\linewidth}\raggedright
Особина
\end{minipage} & \begin{minipage}[b]{\linewidth}\centering
\end{minipage} \\
\midrule\noalign{}
\endhead
\bottomrule\noalign{}
\endlastfoot
Диференцијална једначина &
\((1 - x^{2})y^{''}(x) - 2xy' + k(k + 1)y = 0\) \\
Формула &
\(P_{n}(x) =\)\(\frac{1}{2^{n}}\overset{n}{\sum_{k = 0}}\left( \begin{array}{r}
n \\
k
\end{array} \right)^{2}\)\({(x - 1)}^{n - k}{(x + 1)}^{k}\) \\
Ортогоналност & Домен: \(I = ( - 1,1)\) , тежинска функција:
\(w(x) = 1\) \\
Генератриса &
\(g(x,t) =\)\(\frac{1}{\sqrt{1 - 2xt + t^{2}}}\)\(= \overset{\infty}{\sum_{n = 0}}P_{n}\)\((x)t^{n}\) \\
Родригова формула &
\(P_{n}(x) =\)\(\frac{1}{2^{n}}\overset{n}{\sum_{k = 0}}\frac{1}{2^{n}n!}\frac{d^{n}}{dx^{n}}\)\(\lbrack{(x^{2} - 1)}^{n}\rbrack\) \\
Значајне рекурзивне формуле &
\(P_{n}(x) = \frac{2n - 1}{n}xP_{n - 1}(x) - \frac{n - 1}{n}P_{n - 2}(x)\),
\(n > = 2\)

\(P_{n}'(x) = xP_{n + 1}'(x) + nP_{n - 1}(x)\) \\
Контурна формула &
\(P_{n}(z) =\)\(\frac{1}{2\pi i}\oint{(1 - 2zt + t^{2})}^{- \frac{1}{2}}t^{- n - 1}dt\) \\
\end{longtable}
}

Табела 11.2: Преглед најважнијих особина Лежандрових полинома.

{\def\LTcaptype{none} % do not increment counter
\begin{longtable}[]{@{}
  >{\raggedright\arraybackslash}p{(\linewidth - 2\tabcolsep) * \real{0.3082}}
  >{\centering\arraybackslash}p{(\linewidth - 2\tabcolsep) * \real{0.6918}}@{}}
\toprule\noalign{}
\begin{minipage}[b]{\linewidth}\raggedright
Особина
\end{minipage} & \begin{minipage}[b]{\linewidth}\centering
\end{minipage} \\
\midrule\noalign{}
\endhead
\bottomrule\noalign{}
\endlastfoot
Диференцијална једначина & \(xy^{''} + (1 - x)y' + \lambda y = 0\) \\
Формула & \(L_{n}(x) =\)\(\frac{1}{n!}{(\frac{d}{dx} - 1)}^{n}x^{n}\) \\
Ортогоналност & домен: \(x \in {\mathbb{R}}\), тежинска функција:
\(w(x) = e^{- x}\) \\
Генератриса &
\(g(x,t) =\)\(\overset{\infty}{\sum_{t = 0}}\)\(L_{n}(x)t^{n} =\)\(\frac{e^{- \frac{xt}{1 - t}}}{1 - t} = \frac{e^{x}}{2\pi i}\oint\frac{e^{- z}dz}{z - x - tz}\) \\
Родригова формула &
\(L_{n}^{\alpha}(x) =\)\(\frac{e^{x}x^{- \alpha}}{n!}\frac{d^{n}}{dx^{n}}\)\((e^{- x}x^{n + \alpha})\) \\
Значајне рекурзивне формуле &
\((n + 1)L_{n + 1}(x) - (2n + 1 - x)L_{n}(x) + nL_{n - 1}(x) = 0\)\emph{\textbf{,}}

\(x\frac{dL_{n}(x)}{dx} =\)\(nL_{n}(x) - nL_{n - 1}(x)\) \\
Контурна формула &
\(L_{n}(x) =\)\(\frac{2}{2\pi i}\oint\frac{e^{\frac{- xz}{1 - z}}}{(1 - z)z^{n + 1}}dz\) \\
\end{longtable}
}

Табела 11.3: Преглед најважнијих особина Лагерових полинома.

{\def\LTcaptype{none} % do not increment counter
\begin{longtable}[]{@{}
  >{\raggedright\arraybackslash}p{(\linewidth - 2\tabcolsep) * \real{0.3082}}
  >{\centering\arraybackslash}p{(\linewidth - 2\tabcolsep) * \real{0.6918}}@{}}
\toprule\noalign{}
\begin{minipage}[b]{\linewidth}\raggedright
Особина
\end{minipage} & \begin{minipage}[b]{\linewidth}\centering
\end{minipage} \\
\midrule\noalign{}
\endhead
\bottomrule\noalign{}
\endlastfoot
Диференцијална једначина & \(y'' - xy' + ny = 0\) \\
Формула &
\(H_{n}(x) =\)\(n!\overset{\lbrack n/2\rbrack}{\sum_{k = 0}}\frac{{( - 1)}^{k}}{k!(n - 2k)!}{(2x)}^{n - 2k}\) \\
Ортогоналност & домен: \(x \in {\mathbb{R}}\), тежинска функција:
\(w(x) =\)\(e^{- x^{2}}\) \\
Генератриса & \(g(t,x) = e^{t(2x - t)}\) \\
Родригова формула &
\(H_{n}(x) =\)\(e^{\frac{x^{2}}{2}}\)\({(x - \frac{d}{dx})}^{n}\)\(e^{\frac{- x^{2}}{2}}\) \\
Значајне рекурзивне формуле &
\(H_{n}(x) = 2xH_{n - 1}(x) - (2n - 2)H_{n - 2}(x),n = 2,3,4...\) \\
Контурна формула &
\(H_{n}(z) =\)\(\frac{n!}{2\pi i}\)\(\oint\)\(e^{- t^{2} + 2tz}t^{- n - 1}dt\). \\
\end{longtable}
}

Табела 11.4: Преглед најважнијих особина Ермитовог физичког полинома.

{\def\LTcaptype{none} % do not increment counter
\begin{longtable}[]{@{}
  >{\raggedright\arraybackslash}p{(\linewidth - 2\tabcolsep) * \real{0.3082}}
  >{\centering\arraybackslash}p{(\linewidth - 2\tabcolsep) * \real{0.6918}}@{}}
\toprule\noalign{}
\begin{minipage}[b]{\linewidth}\raggedright
Особина
\end{minipage} & \begin{minipage}[b]{\linewidth}\centering
\end{minipage} \\
\midrule\noalign{}
\endhead
\bottomrule\noalign{}
\endlastfoot
Диференцијална једначина & \(y'' - 2xy' + ny = 0\). \\
Формула &
\(H_{e_{n}}\)\((x) =\)\(2^{- \frac{n}{2}}H_{n}\)\((\frac{x}{\sqrt{2}})\) \\
Ортогоналност & домен: \(x \in {\mathbb{R}}\), тежинска функција:
\(w\overset{}{(x)} = e^{- \frac{x^{2}}{2}}\) \\
Генератриса & \(g(t,x) =\)\(e^{\frac{t}{2}(t - x)}\) \\
Родригова формула & \(H_{e_{n}} = {(x - \frac{d}{dx})}^{n}\) \\
Значајне рекурзивне формуле &
\(H_{e_{n + 1}}\overset{}{(x)} = xH_{e_{n}}\overset{}{(x)} - {H_{e_{n}}}'\overset{}{(x)}\) \\
Контурна формула &
\(H_{e_{n}}\overset{}{\overset{}{(x)}} = \frac{n!}{2\pi i}\oint\frac{e^{tx} - \frac{t^{2}}{2}}{t^{n + 1}}dt\) \\
\end{longtable}
}

Табела 11.5: Преглед најважнијих особина Ермитовог пробабилистичког
полинома.

{\def\LTcaptype{none} % do not increment counter
\begin{longtable}[]{@{}
  >{\raggedright\arraybackslash}p{(\linewidth - 2\tabcolsep) * \real{0.3200}}
  >{\centering\arraybackslash}p{(\linewidth - 2\tabcolsep) * \real{0.6800}}@{}}
\toprule\noalign{}
\begin{minipage}[b]{\linewidth}\raggedright
Особина
\end{minipage} & \begin{minipage}[b]{\linewidth}\centering
\end{minipage} \\
\midrule\noalign{}
\endhead
\bottomrule\noalign{}
\endlastfoot
Диференцијална једначина & \((1 - x^{2})y'' - xy' + n^{2}y = 0\) \\
Формула & \(T_{n}(x) =\)\(\cos(n\arccos x)\) \\
Ортогоналност & Домен \(\lbrack - 1,1\rbrack\), тежинска функција
\(w(x) =\)\({(1 - x^{2})}^{- 1/2}\) \\
Генератриса &
\(T_{n}(x) =\)\(\frac{n}{2}\overset{\lbrack n/2\rbrack}{\sum_{m = 0}}\)\({( - 1)}^{m}\)\(\frac{(n - m - 1)!}{m!(n - 2m)!}\)\({(2x)}^{n - 2m}\) \\
Родригова формула &
\(g(x,t) =\)\(\overset{\infty}{\sum_{n = 0}}\)\(T_{n}(x)t^{n} =\)\(\frac{1 - xt}{1 - 2xt + t^{2}}\). \\
Значајне рекурзивне формуле &
\(T_{n + 1}(x) =\)\(2xT_{n}(x) - T_{n - 1}(x)\),

\((1 - x^{2})T_{n}'(x) = - nxT_{n}(x) + nT_{n - 1}(x)\) \\
\end{longtable}
}

Табела 11.6: Преглед најважнијих особина Чебишевљевих полинома прве
врсте.

{\def\LTcaptype{none} % do not increment counter
\begin{longtable}[]{@{}
  >{\raggedright\arraybackslash}p{(\linewidth - 2\tabcolsep) * \real{0.3350}}
  >{\centering\arraybackslash}p{(\linewidth - 2\tabcolsep) * \real{0.6650}}@{}}
\toprule\noalign{}
\begin{minipage}[b]{\linewidth}\raggedright
Особина
\end{minipage} & \begin{minipage}[b]{\linewidth}\centering
\end{minipage} \\
\midrule\noalign{}
\endhead
\bottomrule\noalign{}
\endlastfoot
Диференцијална једначина & \((1 - x^{2})y'' - 3xy' + n(n + 2)y = 0\) \\
Формула & \(U_{n}(x) =\)\(\frac{\sin((n + 1)x)}{\sin x}\) \\
Ортогоналност & Домен \(\lbrack - 1,1\rbrack\), тежинска функција
\(w\overset{}{(x)} = {(1 - x^{2})}^{\frac{1}{2}}\) \\
Генератриса &
\(U_{n}(x) =\)\(\overset{\lbrack n/2\rbrack}{\sum_{m = 0}}\)\({( - 1)}^{m}\)\(\frac{(n - m)!}{m!(n - 2m)!}\)\({(2x)}^{n - 2m}\) \\
Родригова формула &
\(g(x,t) =\)\(\overset{\infty}{\sum_{n = 0}}\)\(U_{n}(x)t^{n} =\)\(\frac{1}{1 - 2xt + t^{2}}\) \\
Значајне рекурзивне формуле &
\(U_{n + 2}(x) = 2xU_{n + 1}(x) - U_{n}(x),n \geq 0\)

\((1 - x^{2})U_{n}'(x) = - nxU_{n}(x) + (n + 1)U_{n - 1}(x)\) \\
\end{longtable}
}

Табела 11.7: Преглед најважнијих особина Чебишевљевих полинома друге
врсте.

{\def\LTcaptype{none} % do not increment counter
\begin{longtable}[]{@{}
  >{\raggedright\arraybackslash}p{(\linewidth - 2\tabcolsep) * \real{0.3082}}
  >{\centering\arraybackslash}p{(\linewidth - 2\tabcolsep) * \real{0.6918}}@{}}
\toprule\noalign{}
\begin{minipage}[b]{\linewidth}\raggedright
Особина
\end{minipage} & \begin{minipage}[b]{\linewidth}\centering
\end{minipage} \\
\midrule\noalign{}
\endhead
\bottomrule\noalign{}
\endlastfoot
Диференцијална једначина &
\((1 - x^{2})y'' + (\beta - \alpha - (\alpha + \beta + 2)x)y' + n(n + \alpha + \beta + 1)y = 0\) \\
Формула &
\({P_{n}}^{(\alpha,\beta)}\)\((x) =\)\(\frac{{( - 1)}^{n}}{2^{n}n!}{(1 - x)}^{- \alpha}{(1 + x)}^{- \beta}\frac{d^{n}}{dx^{n}}\left\lbrack {(1 - x)}^{\alpha + n}{(1 + x)}^{\beta + n} \right\rbrack\) \\
Ортогоналност & Домен: , тежинска функција:
\(\rho(x) =\)\({(1 - x)}^{\alpha}{(1 + x)}^{\beta}\). \\
Генератриса &
\(g(t,x) =\)\({(1 - 2xt + t^{2})}^{- \alpha} \cdot {(1 + 2xt + t^{2})}^{- \beta}\) \\
Родригова формула &
\({(1 - x)}^{\alpha}{(1 + x)}^{\beta}P_{n}^{(\alpha,\beta)}\)\((x) =\)\(\frac{{( - 1)}^{n}d^{n}}{2^{n}n!dx^{n}}\lbrack{(1 - x)}^{\alpha + n}{(1 + x)}^{\beta + n}\rbrack\) \\
Значајне рекурзивне формуле &
\(2n(n + \alpha + \beta)(2n + \alpha + \beta - 2)\)\(P_{n}^{(\alpha,\beta)}\)\((x) =\)

\((2n + \alpha + \beta - 1)\)\(\{(2n + \alpha + \beta)(2n + \alpha + \beta - 2)x + \alpha^{2} - \beta^{2}\}\)\(P_{n - 1}^{(\alpha,\beta)}\)\((x) -\)\(2(n + \alpha - 1)(n + \beta - 1)(2n + \alpha + \beta)\)\(P_{n - 2}^{(\alpha,\beta)}\)\((x)\),
\(n \geq 2\),

\(P_{n}^{(\alpha,\beta)'}\)\((x) =\)\(\frac{1}{2}\)\((1 + \alpha + \beta + n)P_{n - 1}^{(\alpha + 1,\beta + 1)}\)\((x)\) \\
\end{longtable}
}

Табела 11.8: Преглед најважнијих особина Јакобијевих полинома.

\begin{itemize}
\item
  \protect\phantomsection\label{_Toc37}{}Предлози за рачунарску
  имплементацију
\end{itemize}

Избор алгоритама који ће се користити у имплементацији \emph{Калкулатора
неких специјалних функција} зависи пре свега од избора програмских
језика, библиотека и, генерално, врсте апликације - веб апликација,
мобилна, десктоп апликација, итд. Будући да је основна идеја направити
нешто што је лако доступно, не мора да се инсталира, захтева што мање
корисничких акција да се дође до резултата - избор врсте апликације пада
на веб апликацију.

Када су у питању нумеричке апроксимације и, генерално, разне врсте
прорачуна, често избор пада на програмски језик \emph{Python}. Разлози
за то су бројни: он има велику подршку управо за математичке прорачуне,
у виду библиотека попут \emph{SciPy} (која чак поседује и модул који има
већ имплементиране методе за рачунање вредности специјалних функција,
али идеја овог рада је свакако имплементација калкулатора, а не
коришћење готових решења) и других, као и библиотеке које омогућавају
визуелизацију. Осим тога, \emph{Python} има велику популарност због
своје синтаксе и читљивости кода, објектно-оријентисан је, лако се
интегрише са другим програмским језицима, преносив је, итд. У поређењу с
другим често коришћеним програмским језицима попут \emph{C/C++, Java,
JavaScript} и других, \emph{Python} је спорији језик, захтеван је по
питању ресурса и није добар избор за развој мобилних апликација
{[}27{]}.

\emph{C} је процедурални језик који је добио на популарности због своје
моћи да извуче максимум из архитектуре рачунара без коришћења
асемблерског кода, као и због минималистичке синтаксе, високе читљивости
кода, одличних перформанси. Контрола над меморијом је доста прецизнија
јер се меморија мануелно алоцира и деалоцира, али ово може и успорити
развој (код других језика високог нивоа програмер не мора да води рачуна
о алокацији меморије) и повећава могућност грешака везаних за меморију.
Библиотека \emph{\textless math.h\textgreater{}} поседује бројне
функције и макрое који су потребни за развој апликације попут
\emph{Калкулатора специјалних функција} {[}8{]}. \emph{GNU Scientific
Library (GSL)} је нумеричка библиотека написана управо у језику \emph{C}
и намењена извођењу научних прорачуна, изузетно добрих перформанси и са
преко 1000 корисних функција. \emph{GSL} je намењена програмирању у
\emph{C} и \emph{C++,} али је могуће интегрисати је и са другим
програмским језицима писањем омотача {[}55{]}.

\emph{C++} би такође био сјајан избор за програмирање овог калкулатора,
будући да је заснован на \emph{C} програмском језику даје све набројане
предности које даје и \emph{C} , али за разлику од њега нуди и
објектно-оријентисан приступ и адресира неке проблеме и питања која
постоје у језику \emph{C} {[}42{]}. Осим поменутих математичких
библиотека, за научне прорачуне постоји и подршка у виду
\emph{Armadillo} библиотеке, као и бројних других {[}13{]}.

Joш један од адекватних избора програмског језика би била \emph{Java -}
објектно-оријентисани језик, високих перформанси, преносив. Један од
највећих разлога за његову популарност језика је агилност - овај
програмски језик се годинама прилагођава променама у програмерском
окружењу и променама у самом начину рада програмера {[}40{]}. У поређењу
са до сада набројаним језицима у овом поглављу: смањује могућности за
грешке са меморијом јер за разлику од \emph{C/C++} поседује
\emph{Garbage collector}, може бити мало спорији него \emph{C}++,
једноставнији је за коришћење јер се сви типови преносе по вредности,
итд. Поседује веома обимну стандардну библиотеку, а и математичке
библиотеке попут \emph{JScience, Colt} и бројне друге. Додатно,
\emph{Java} има добро развијене радне оквире (\emph{frameworks}) за
развој веб апликација.

\emph{JavaScript} je брз, објектно-оријентисан jeзик високог нивоа,
једноставне синтаксе, коришћен за веб програмирање. Може се користити и
за серверску и за клијентску страну, иако је првобитно направљен за
клијентску страну и може се извршавати само у прегледачима. Да би се
користио на серверској страни, потребно је извршно окружење
(\emph{runtime environment}), попут \emph{Node.js}-a. Може бити мало
спорији од \emph{C++} и \emph{Java-}е\emph{,} али то је изражено тек код
огромних апликација. Он је слабо типизиран и ради динамичко везивање,
што даје већи степен слободе на неки начин, али даје и доста простора за
грешке при извршавању. Тај проблем делимично превазилази
\emph{TypeScript -} он се назива надскупом програмског језика
\emph{JavaScript,} а заправо се преводи у \emph{JavaScript:} он најбоље
ради када има све типове за време компајлирања, али није неопходно да их
зна. \emph{Math.js} je моћна библиотека за ~\emph{Node.js} са мноштвом
корисних метода, константи, па и типова попут \emph{BigNumber} за рад са
бројевима у већем опсегу од уобичајеног.

Разлог због чега je \emph{TypeScript} одабран за развој
\emph{Kалкулатора неких специјалних функција}, упркос томе што не
доминира над осталим наведеним језицима по брзини је што је у питању
израда веб апликације - ради брзине самог развоја апликације и
једноставности, користиће се исти језик на клијентској и серверској
страни. Архитектура саме апликације ће бити монолитна јер у овом
тренутку нема потребе за развојем микросервисне архитектуре, али уколико
би се додавале функционалности у њу, радиле оптимизације итд, појавила
би се и потреба за издвајањем микросервиса. Дакле, за потребе развоја
серверске (бекенд) стране апликације користиће се \emph{Node.js}
(верзија v21.7.0)\emph{,} док ће се за развој кликентске стране
(фронтенд) користити \emph{Angular} (верзија 16.2.12).

\begin{itemize}
\item
  \protect\phantomsection\label{_Toc38}{}Рачунарска имплементација
  \emph{Kалкулатора неких специјалних функција}
\end{itemize}

За сваку од специјалних функција која је представљена у првом делу овог
рада дато је више различитих репрезентација. За Чебишевљеве полиноме
прве и друге врсте дефиниције су такве да се њихова вредност може
једноставно израчунати. Вредност Ермитовог пробабилистичког полинома се
такође може израчунати брзо и једноставно, користећи везу са Ермитовим
физичким полиномом, под претпоставком да је рачунање физичког полинома
имплементирано. Остале специјалне функције имају већ комплексније
алгоритме и у овом делу рада ће бити објашњени принципи по којима су они
имплементирани.

Такође, будући да је један од основних мотива овог рада било пружити
функционалну апликацију за што једноставније израчунавање вредности
неких често коришћених специјалних функција, апликација не би била
савремена без подршке за унос у виду израза. Другим речима, имплемениран
је и класични калкулатор помоћу ког се евалуира вредност неког израза
која ће представљати улазну вредност - променљиву, ред специјалне
функције, прецизност (у неким случајевима).

Још једна функционалност без које апликација не би била потпуна је и
графички приказ вредности функције. Ово је реализовано уз помоћ
библиотеке \emph{plotly.js -} при сваком рачунању за унете вредности
исцртава се график за унети ред функције.

У наредним поглављима биће дат преглед апликације, имплементација
евалуатора израза (калкулатора на пољима за унос), као и имплементација
за израчунавање вредности самих специјалних функција.

\begin{itemize}
\item
  \protect\phantomsection\label{_Toc39}{}Преглед апликације
\end{itemize}

Жеља нам је била да апликација поседује веома интуитиван графички
кориснички интерфејс. На почетној страници се налази мени за избор саме
функције, као и заглавље и подножје (енг. \emph{header} и
\emph{footer}), које се налазe на свакој страници апликације и садрже
пречице до неких страница попут информативне странице, странице с
контактом, дугме за промену језика, повратак на почетну
страну.\includegraphics[width=6.17212in,height=2.91635in,alt={pasted-movie.png}]{media/image13.png}

Слика 13.1.1: Почетна страна апликације.

По избору се учитава страница за ту одређену функцију која садржи форму
за унос података. Може се унети вредност у децималном или
експоненцијалном запису или се може евалуирати вредност израза. Кликом
на дугме на одређеном пољу отвара се калкулатор и вредност која се
израчуна може се употребити као улазна вредност за то поље. На наредној
слици 13.1.2 може се видети форма за унос улазних вредности за Беселову
функцију прве
врсте.\includegraphics[width=4.85785in,height=2.71546in,alt={pasted-movie.png}]{media/image14.png}

Слика 13.1.2: Изглед форме за унос података.

Сама форма зависи од одабране функције, односно од тога који су улазни
подаци неопходни, па тако у случају Беселове функције поља су ред,
прецизност и променљива, док нпр. за Јакобијеве полиноме форма укључује
ред, променљиву и два параметра. На сваком пољу се налази лабела која
даје информацију о типу тог поља. Постављене су и иницијалне вредности
за свако поље. Апликација ради и валидацију унетих поља, тако да се у
случају невалидног уноса јавља порука о грешци и дугме за израчунавање
резултата је онемогућено, као на слици
13.1.3.\includegraphics[width=3.94072in,height=2.2028in,alt={pasted-movie.png}]{media/image15.png}

Слика 13.1.3: Пример некоректног уноса.

На пољу је уочљиво и дугме на чији се клик отвара компонента која
представља калкулатор који може евалуирати неки израз како би он био
употребљен као улазна вредност за то поље. Компонента калкулатор се
појављује испод форме за унос, а кликом на исто дугме она нестаје.

\includegraphics[width=5.04549in,height=4.97454in,alt={pasted-movie.png}]{media/image16.png}

Слика 13.1.4: Унос поља помоћу калкулатора.

Израчуната вредност се може употребити на том пољу на ком је иницијално
кликнуто на дугме, тако да уколико би корисник хтео да употреби
израчунату вредност на другом пољу, морао би да кликне на поље на коме
жели да га употреби, без обзира што је компонента већ отворена. Ово
можда не звучи као најинтуитивније решење, али поље на коме је
калкулатор отворен је истакнуто другом бојом и ефектима тако да то
кориснику јасно сугерише да вредност може употребити само на том пољу.

Када је израз евалуиран, дугме ''Употреби ову вредност'' ће постати
омогућено, тако да кликом на то дугме корисник уноси вредност у форму.
На слици 13.4.1. се може видети и како описана акција клика на дугме за
калкулатор изгледа. О самој компоненти и њеној имплементацији ће бити
више речи у наредном делу рада.

На претходној слици 13.1.4. такође може уочити да поред наслова функције
постоји и дугме за више информација. Преласком курсора преко овог
дугмета појављује се и објашњење (енг. \emph{explainer}) о намени
дугмета. Кликом на дугме отвара се нови прозор на коме се налазе опште
информације о одабраној специјалној функцији: њена дефиниција, домен,
дефинисаност, понашање, итд. Ово дугме наменски води корисника на нови
прозор из два разлога: да не би изгубио фокус на главну функционалност у
апликацији, али и да не би изгубио прогрес уколико је израчунао већ неку
вредност, како не би морао да поново уноси податке.

\includegraphics[width=6.80482in,height=3.2153in,alt={pasted-movie.png}]{media/image17.png}

Слика 13.1.5: Опште информације о функцији.

Наслов апликације на заглављу је такође и веза ка почетној страници.
Преласком курсора преко слова (енг. \emph{on hover}) она мењају боју што
сугерише да је реч о дугмету/линку. У заглављу се такође налази и дугме
за избор језика које отвара падајући мени. Функционалност промене језика
је имплементирана користећи посебну класу названу \emph{LanguageService}
која стринг који представља изабрани језик чува у \emph{localStorage.}
При промени избора језика, сачува се нови избор и емитује се догађај
који ослушкује свака компонента која поседује текстуална поља.

Компоненте увозе речнике сачуване у \emph{json} фајловима и сав текст
приказан у \emph{html} страницама је динамички
рендерован.\includegraphics[width=1.18in,height=1.52881in,alt={pasted-movie.png}]{media/image18.png}

Слика 13.1.6: Избор језика у апликацији.

У подножју сваке странице се налазе линкови до две странице: страницу са
основним информацијама о апликацији, као и страницу са формом за
контактирање (''корисничка подршка'').

Форма за контактирање је замипљена тако да се аутоматски пошаље и-мејл с
подацима из форме на адресу направљену за потребе апликације.

На слици 13.1.7. се може видети страница за контактирање корисничке
подршке и страница са основним информацијама о одабраној функцији.

Слика 13.1.7: Остале странице у
апликацији.\includegraphics[width=6.69298in,height=3.15537in,alt={pasted-movie.png}]{media/image19.png}\includegraphics[width=6.69298in,height=3.15537in,alt={pasted-movie.png}]{media/image20.png}

\begin{itemize}
\item
  \protect\phantomsection\label{_Toc40}{}Компонента за евалуацију израза
\end{itemize}

На слици 13.1.4. се може видети компонента калкулатора за евалуацију
израза. Осим основних рачунских операција, подржава и тригонометријске
функције, експоненцијалну, логаритамску, хиперболичке функције, као и
више нивоа заграда.

За имплементацију ове компоненте коришћен је стек, а алгоритми у
псеудокоду који су коришћени, уз проширења и одређене модификације, су
алгоритми за обраду аритметичких израза преузети из књигa
{[}45{]}{[}26{]}.

Сама \emph{angular} компонента \emph{CalculatorComponent} користи класу
\emph{CalculatorService} у којој се практично налази сва логика за
евалуацију израза. \emph{CalculatorComponent} дефинише омотачке методе
како би из \emph{html} фајла ове компоненте могле да се позивају
одговарајуће методе сервиса.

private infixToPostfix() \{

const stack = new Stack\textless string\textgreater();

let rank = 0, i = 0, x: string, postfix: string{[}t{]} = {[}t{]}

let next: string \textbar{} undefined = this.infix{[}tI++{]};

while (next) \{

if (!isOperator(next) \&\& next != \textquotesingle(\textquotesingle{}
\&\& next != \textquotesingle)\textquotesingle) \{

// operand

postfix.push(next);

rank = rank + 1;

\} else \{

if (isUnaryOperator(next)) \{

stack.push(next);

rank = rank + getPriority(next).R;

\} else \{

while (

!stack.isEmpty() \&\& getPriority(next).ipr \textless=
getPriority(stack.top()).spr) \{

x = stack.pop() ?? \textquotesingle\textquotesingle; // pop return type
can be undefined

postfix.push(x);

rank = rank + getPriority(x).R;

if (rank \textless{} 1) \{ return INVALID\_EXPRESSION; \}

\} // end\_while

if (stack.isEmpty() \textbar\textbar{} next !=
\textquotesingle)\textquotesingle) \{ stack.push(next); \}

else \{ x = stack.pop() ?? `\textquotesingle; \}

\}// end\_if for case binary operator

\}

next = i \textless{} this.infix.length ? this.infix{[}ti++{]} :
undefined;

\} // finished iterations trough infix expression

while (!stack.isEmpty()) \{

x = stack.pop() ?? \textquotesingle\textquotesingle;

postfix.push(x);

rank = rank + getPriority(x).R;

\}

if (rank != 1) \{

return INVALID\_EXPRESSION;

\}

return postfix;

\}

Разлог за коришћење метода датог изнад је тај што је постфиксна нотација
знатно погоднија за израчунавање вредности израза него инфиксна. Стек се
у овом алгоритму користи из разлога што сама природа израза сугерише
потребу за структуром података са LIFO понашањем - подизраз који је у
најугњежденијем пару заграда мора бити израчунат први.

Пошто морамо дефинисати редослед операција, уводи се IPR-\emph{input
priority} и SPR-\emph{stack priority}, oдносно тзв. улазни приоритет на
стек и приоритет када је елемент већ на стеку. Унарни оператори имају
већи приоритет од бинарних, операције сабирања и одузимања имају исте
приоритете, али су им приоритети најмањи од свих бинарних оператора,
итд.

Да би алгоритам истовремено могао и да валидира и исправност унетог
инфиксног израза, уведен је и ранг сваком оператору, а он се дефинише
као \emph{R=1-n,} где је \emph{n} број операнада за тај оператор, па је
тако вредност ранга за све бинарне операторе -1, а за све унарне 0.

У наредној табели ће бити дате вредности приоритета и ранга за сваки
операнд (и заграде).

{\def\LTcaptype{none} % do not increment counter
\begin{longtable}[]{@{}
  >{\raggedright\arraybackslash}p{(\linewidth - 6\tabcolsep) * \real{0.3134}}
  >{\raggedleft\arraybackslash}p{(\linewidth - 6\tabcolsep) * \real{0.1204}}
  >{\raggedleft\arraybackslash}p{(\linewidth - 6\tabcolsep) * \real{0.1347}}
  >{\raggedleft\arraybackslash}p{(\linewidth - 6\tabcolsep) * \real{0.1228}}@{}}
\toprule\noalign{}
\begin{minipage}[b]{\linewidth}\raggedright
\end{minipage} & \begin{minipage}[b]{\linewidth}\raggedleft
IPR
\end{minipage} & \begin{minipage}[b]{\linewidth}\raggedleft
SPR
\end{minipage} & \begin{minipage}[b]{\linewidth}\raggedleft
Ранг
\end{minipage} \\
\midrule\noalign{}
\endhead
\bottomrule\noalign{}
\endlastfoot
+ , - & 2 & 2 & -1 \\
* , / & 3 & 3 & -1 \\
\^{} & 4 & 5 & -1 \\
Остали бинарни оператори & 6 & 6 & -1 \\
Унарни оператори & 7 & 7 & 0 \\
( & 8 & 0 & - \\
) & 1 & 0 & - \\
\end{longtable}
}

Табела 13.2.1: Приоритет оператора коришћен у алгоритму за конверзију
инфиксног израза у постфиксни.

Као што је већ напоменуто, псеудокод за овај алгоритам налази се у
{[}45{]}, али он обухвата само простије бинарне операторе: сабирање,
множење, одузимање, дељење и степеновање. За потребе имплементације у
овом раду је алгоритам проширен са унарним операторима: унарни минус,
природни логаритам, логаритам за основу 10, квадратни и кубни корен,
тригонометријске функције, итд., као и бинарним операторима попут
логаритма за задату основу, \emph{n-}тог корена, итд.

Инфиксни израз је у овој имплементацији низ стрингова - сваки операнд,
оператор и заграда су један елемент низа. Алгоритам итерира кроз овај
низ, односно инфиксни облик унетог израза и када се наиђе на унарни
оператор или операнд, он се ставља на стек и вредност ранга се увећава у
зависности од тренутног елемента. Када се наиђе на бинарни оператор, док
год се не испразни стек или улазни приоритет тренутног елемента не
постане већи од приоритета на стеку оног елемента који је на врху стека,
узима се последњи елемент са стека и ставља се на крај низа
\emph{postfix}, који представља излаз ове функције - постфиксни облик
задатог израза. Увећава се вредност ранга и у свакој итерацији се врши
провера вредности ове променљиве - ако је она мања од 1 то значи да је
израз некоректан и одмах се излази из функције. Када се изађе из ове
петље, ако је стек празан или тренутни елемент није затворена заграда,
елемент се ставља на стек. У супротном, узима се последња вредност са
стека.

Када је алгоритам прошао кроз цео улазни низ, ако стек није празан,
узима се вредност с врха стека, ставља на крај излазног низа и ажурира
вредност ранга.

На крају, врши се последња провера исправности унетог израза у инфиксној
нотацији - ако је вредност ранга тачно 1, израз је коректан, у супротном
треба исписати поруку о грешци.

Дакле, када је позван метод \emph{evaluate()}, он позива приватни метод
i\emph{nfixToPostfix(),} који може вратити или низ стрингова који
представља постфиксни облик израза, или поруку о грешци. Уколико се
врати порука о грешци, функција прекида са радом и исписује грешку на
дисплеју компоненте. Ако је враћен коректни постфиксни запис, зове се
приватни метод \emph{evaluatePostfixExpression}, која за задати
постфиксни израз евалуира његову вредност. У наставку ће бити дата
имплементација ове функције.

private evaluatePostfixExpression(postfix: string{[}t{]}): string \{

const stack = new Stack\textless string\textgreater();

let i = 0, x: string, rez;

while (i \textless{} postfix.length) \{

x = postfix{[}ti++{]};

if (!isOperator(x)) \{

// operand

stack.push(x);

\} else \{

if (isUnaryOperator(x)) \{

if (stack.isEmpty()) return INVALID\_EXPRESSION;

const operand = this.math.bignumber(stack.pop());

rez = this.calculateUnaryOperation(x, operand);

stack.push(rez.toString());

\} else \{

// if it\textquotesingle s not operand nor unary operator,
it\textquotesingle s binary operator

if (stack.size() \textless{} 2) return INVALID\_EXPRESSION;

const operand2 = this.math.bignumber(stack.pop());

const operand1 = this.math.bignumber(stack.pop());

rez = this.calculateBinaryOperation(operand1, operand2, x);

stack.push(rez.toString());

\}

\}

\} // end while

rez = stack.pop() ?? INVALID\_EXPRESSION;

if (stack.isEmpty()) return rez;

return INVALID\_EXPRESSION;

\}

Уколико је тренутни елемент бинарни оператор, поново треба проверити
стање стека. Уколико је све како треба, скинути две вредности са врха
стека, над њима применити бинарни оператор и резултат вратити на стек.

Када је алгоритам прошао кроз цео израз, под условом да је унесени израз
у постфиксној нотацији коректан, на стеку би морала да буде тачно једна
вредност - резултат евалуације унетог израза.

На крају, сам метод \emph{evaluate}, који је и позвао дате две функције,
поставља вредности одређених променљивих које се користе за сам унос и
потом враћа резултат. Постављање поменутих променљивих се врши како би
добијени резултат могао да се користи као операнд за даље рачунање.

Када су у питању методи \emph{calculateBinaryOperation} и
\emph{calculateUnaryOperation}, који се користе при евалуацији израза,
су прилично једноставне функције које се састоје од по једне наредбе
вишеструког гранања \emph{switch}, која утврђује која од подржаних
операција треба да се изврши и потом је извршава коришћењем метода
библиотеке \emph{mathjs} (верзија коришћена при изради је
11.12.0)\emph{.} Сва израчунавања врше се с прецизношћу од 64 децимале.

Када је израчуната вредност унетог израза, корисник може наставити даље
израчунавање (користећи резултат као операнд), ресетовати калкулатор и
рачунати од почетка, или употребити добијену вредност као унос за поље с
ког је компонента отворена. Ако се деси да је израчуната вредност коју
корисник жели да употреби ван опсега за дато поље, добиће се порука о
грешци.

Дугме ''Употреби ову вредност'' је онемогућено док корисник не израчуна
вредност, тек када је вредност у калкулатору израчуната (и није дошло до
грешке), дугме ће моћи да се употреби јер калкулаторска компонента, при
израчунавању, шаље догађај родитељској компоненти у којој се налази то
дугме.

Када је у питању унос у калкулатор, на дисплеју се приказује последњи
унети операнд, а изнад калкулатора ће се приказати последњи унети
оператор или заграда. Последњи унети операнд се у класи
\emph{CalculatorService} чува у пољу \emph{currentOperand,} он се додаје
на крај низа \emph{infix} у моменту када се унесе наредни оператор,
затворена заграда или једнакост. То је имплементирано на тај начин због
саме природе уноса операнада - треба потенцијално притиснути више
дугмића за један операнд, тако да је једини начин да знамо да је унос
операнда завршен када се унесе оператор, заграда или једнакост.

\begin{itemize}
\item
  \protect\phantomsection\label{_Toc41}{}Графички приказ функције
\end{itemize}

Као што је раније поменуто, при сваком израчунавању функције за задате
параметре, графички се приказује и график функције помоћу библиотеке
\emph{plotly.js,} верзија 2.29.1\emph{.}

У зависности од одабране функције генерише се 200 тачака за које се
израчунавају вредности функције и представљају се на графику.

За Лежандрове, Чебишевљеве прве и друге врсте, као и за Јакобијеве
полиноме, се узимају се вредности \(x\) за променљиву од -0.999999 до
0.999999, јер су ове функције и дефинисане на домену (-1,1). За бета и
гама функције се узима интервал \((x - 2,x + 2)\), односно \((0,x + 2)\)
када је \(x < 2\). За све остале специјалне функције се узима интервал
\((x - 3,x + 3)\), где је \(x\) унета променљива. Овај интервал се бира
да би се на графику видела тачка у \(x\) као централна у некој својој
симетричној околини.

const step: number = (endValue - startValue) / (numParameters - 1);

const xArr: number{[}t{]} = Array.from(

\{ length: numParameters \},

(\_, index) =\textgreater{} startValue + index * step

);

const yArr = xArr.map((x) =\textgreater{} spef?.calculate(\{ ...data, x:
x \}) ?? 0);

За унету вредност променљиве тачка се представља црвеном бојом да би се
јасно видело где се та вредност на графику налази.

Компонента која која приказује график је
\emph{SpecialFunctionComponent.} Она је родитељска компонента за
компоненту за унос података, која, када корисник унесе податке, шаље
догађај родитељској компоненти. \emph{SpecialFunctionComponent,} када
прими тај догађај, зове метод за рачунање вредности одабране специјалне
функције за унете податке, потом функцију која генерише координате
тачака које ће бити приказане на графику (део овог метода је и код
изнад) и на крају с тим координатама и координатама тачке која одговара
унетим вредностима (њу шаљемо посебно јер ће се исцртати другачијом
бојом да би била истакнута) зове функцију која ће на задатом HTML
елементу коришћењем \emph{Plotly} библиотеке исцртати тај график. Ова
функција је дата у наредном исечку кода.

export function drawGraph(

element: HTMLElement,

xCoordinates: number{[}t{]},

yCoordinates: number{[}t{]},

x: number,

y: number

) \{

const functionGraph: Partial\textless Plotly.ScatterData\textgreater{} =
\{

x: xCoordinates,

y: yCoordinates,

type: \textquotesingle scatter\textquotesingle,

mode: \textquotesingle lines+markers\textquotesingle,

marker: \{ size: 0.5 \},

\};

const calculatedValue: Partial\textless Plotly.ScatterData\textgreater{}
= \{

x: {[}tx{]},

y: {[}ty{]},

mode: \textquotesingle markers\textquotesingle,

type: \textquotesingle scatter\textquotesingle,

marker: \{ size: 5, color: \textquotesingle red\textquotesingle{} \}, //
Color for the specific spot

\};

const layout: Partial\textless Plotly.Layout\textgreater{} = \{

title: \textquotesingle Graph\textquotesingle,

xaxis: \{ title: `x\textquotesingle{} \},

yaxis: \{ title: `f(x)\textquotesingle{} \},

showlegend: false,

margin: \{ r: 50, b: 50, t: 50 \},

hovermode: `closest\textquotesingle,

autosize: false,

\};

const data = {[}tfunctionGraph, calculatedValue{]};

Plotly.newPlot(element, data, layout);

\}

Последња два параметра метода су координате тачке која треба да се
истакне другачијом бојом. Оне су прослеђене засебно јер је вредност већ
израчуната, да се не би функција звала више пута за исти скуп вредности.

\begin{itemize}
\item
  \protect\phantomsection\label{_Toc42}{}Израчунавање вредности
  специјалне функције
\end{itemize}

За сврху израчунавања вредности специјалне функције креирана је
апстрактна класа \emph{SpecialFunction} која декларише метод
\emph{calculate(params:~FunctionParamsForCalculation)} и
\emph{calculate64(params: FunctionParamsForCalculationWithBigNumbers).}
Оба метода имплементирају исти алгоритам, само што други апроксимира
вредност функције с прецизношћу од 64 значајне цифре, а прва са
стандардном прецизношћу. Разлози за постојање оба метода су: графички
приказ - није неопходно оптерећивати апликацију прецизним израчунавањима
у сврху исцртавања на графику, као и тестирање и проналажење грешака -
упоређивањем добијених резутата ових метода много лакше се уочавају
грешке.

За сваку функцију постоји посебна класа и све те класе наслеђују
поменуту апстрактну класу, па се у функцији calculate налази рачунарска
имплементација евалуације вредности одабране функције. У овом поглављу
ће бити дата имплементација за сваку од функција.

\begin{itemize}
\item
  \protect\phantomsection\label{_Toc43}{}Беселова функција прве врсте
\end{itemize}

За рачунарску имплементацију евалуације Беселове функције прве врсте
погодна је

дефиниција 2 према којој је:
\(J_{\alpha}(x) =\)\(\overset{\infty}{\sum_{m = 0}}\frac{{( - 1)}^{m}\left( \frac{x}{2} \right)^{\alpha + 2m}}{m!\Gamma(\alpha + m + 1)}\).

Будући да је у питању бесконачна сума, неопходно је да корисник унесе и
прецизност с којом ће сума бити израчуната, односно, у оном моменту када
је однос тренутног сабирка и суме мањи од прецизности петља ће бити
прекинута. Прецизност може бити до 64 децимале, што се такође проверава
при валидацији.

Када би имплементација била само за целобројне вредности реда, погодна
би била и дефиниција 3, која функцију представља у облику
\(J_{n}(x) =\)\(\overset{\infty}{\sum_{m = 0}}\frac{{( - 1)}^{m}}{(m + n)!m!}\left( \frac{x}{2} \right)^{2m + n}\),
што знатно поједностављује имплементацију. Међутим, за реалне
нецелобројне вредности реда, неопходно је имплементирати и апроксимацију
гама функције.

Нека је \(m\)-ти члан суме дате дефиницијом 3:

\(t_{m} = \frac{{( - 1)}^{m}\left( \frac{x}{2} \right)^{\alpha + 2m}}{m!\Gamma(\alpha + m + 1)}\).

Однос било која два узастопна члана је тада:

\(R_{m + 1} =\)\(\frac{t_{m + 1}}{t_{m}}\)\(=\)
\(\frac{\frac{{( - 1)}^{m + 1}\left( \frac{x}{2} \right)^{\alpha + 2m + 2}}{(m + 1)!\Gamma(\alpha + m + 2)}}{\frac{{( - 1)}^{m}\left( \frac{x}{2} \right)^{\alpha + 2m}}{m!\Gamma(\alpha + m + 1)}}\)
\(= - \frac{\left( \frac{x}{2} \right)^{2}}{m + 1}\frac{\Gamma(\alpha + m + 1)}{\Gamma(\alpha + m + 2)}\),
\(m \geq 0\).

Ово може да се искористи да се напише ефикасан алгоритам - уместо да се
сваки сабирак суме рачуна по формули, он ће бити израчунат једноставно
као умножак претходног и \(R_{m}\).

Сада је још потребно извести и први сабирак:

\(t_{0} =\)\(\frac{{( - 1)}^{0}\left( \frac{x}{2} \right)^{\alpha + 0}}{0!\Gamma(\alpha + 1)} =\)\(\frac{\left( \frac{x}{2} \right)^{\alpha}}{\Gamma(\alpha + 1)}\).

Алгоритам се тада своди на:

\begin{itemize}
\item
  Израчунаваење \(t_{0}\) по горе изведеној формули, иницијална сума
  постаје \(t_{0}\).
\item
  Петљу која полази од \(m = 1\) и извршава се док год
  је\(\left| \frac{t_{m}}{S} \right| \geq \varepsilon\), при чему је
  \(S\) сума (у датој итерацији са \(m\) чланова), а \(\varepsilon\)
  унета прецизност. У петљи сума добија вредност
  \(S = S + t_{m - 1}R_{m}\).
\end{itemize}

Да би се алгоритам додатно оптимизовао, минимализоваће се број
операција, тако да ће у свакој итерацији за израчунавање \(R_{m + 1}\)
да се користе већ израчунате вредности из претходне итерације и само се
рачунати вредност гама функције за вредност тренутног аргумента.

calculate64(params: FunctionParamsForCalculationWithBigNumbers): string
\{

const \{ alpha, x, eps \} = this.stringToBigNumber(params);

let gammaArg = this.math.add(alpha, BIG\_NUMBER\_CONSTANTS.ONE);

let gammaCurrent = this.calculateGamma64(gammaArg);

const xHalf = this.math.divide(x, BIG\_NUMBER\_CONSTANTS.TWO);

const xHalfSqr = this.math.pow(xHalf, BIG\_NUMBER\_CONSTANTS.TWO);

const xHalfPowered = this.math.pow(xHalf, alpha);

let t = this.math.divide(xHalfPowered, gammaCurrent);

let sum = t, gammaPrevious, R;

for (

let m = 1;

Number(this.math.compare(this.math.abs(this.math.divide(t, sum)), eps))
\textgreater{} 0;

m++

) \{

gammaPrevious = gammaCurrent;

gammaArg = this.math.add(gammaArg, BIG\_NUMBER\_CONSTANTS.ONE);

gammaCurrent = this.calculateGamma64(gammaArg);

R = this.math.divide(gammaPrevious, gammaCurrent);// R=G(n+m+1)/G(n+m+2)

R = this.math.multiply(R, xHalfSqr); // R = R / sqrt(x/2)

R = this.math.divide(R, m); // R = R/m

R = this.math.unaryMinus(R); // R = - R

t = this.math.multiply(R, t); // t = t * R

sum = this.math.add(t, sum); // sum = sum + t

\}

return sum.toString();

\}

Дакле, пре свега се израчунају аргумент за гама функцију и вредност гама
функције за нулти елемент суме, као и константе \(\sqrt{\frac{x}{2}}\) и
\(\left( \frac{x}{2} \right)^{\alpha}\), а потом се израчунава и нулти
елемент као
\(\frac{\left( \frac{x}{2} \right)^{\alpha}}{\Gamma(\alpha + 1)}\).
Будући да је сума бесконачна, у петљи се елементи сумирају док год однос
наредног елемента и до тада израчунате суме не постане мањи од унете
прецизности. У петљи се рачуна вредност количника између тренутног и
претходног елемента, потом тренутни елемент и на крају се тренутни
сабирак додаје суми.

За рачунање ове функције користи се и гама функција, тако да се у
конструктору саме класе креира објекат класе која имплементира рачунање
гама функције, а приватан метод \emph{calculateGamma64} је само омотачка
функција за метод класе \emph{GammaFunction,} креирана због неопходних
конверзија\emph{.} У следећем потпоглављу ће бити детаљљно описана
имплементација апроксимације гама функције Ланцошевим методом.

Резултат рада метода \emph{calculate64} приказан је приказан на дисплеју
с десне стране странице, док се испод дисплеја налази график функције,
као на слици 13.4.1.

С леве стране остаје форма за унос и корисник може наставити да уноси
нове вредности уколико му је то потребно.

\includegraphics[width=6.69298in,height=2.84396in,alt={pasted-movie.png}]{media/image21.png}

Слика 13.4.1. Приказ резултата рада класе која апроксимирa вредност
Беселове функције.

calculate(params: FunctionParamsForCalculation): number \{

const \{ alpha, x \} = params;

const eps: number = params.eps ?? 10 ** -15;

const xHalf = x / 2.0;

const xHalfSqr = xHalf ** 2;

let gammaArg = alpha + 1.0;

let gammaCurrent: number = this.math.gamma(gammaArg);

let t = xHalf ** alpha / gammaCurrent;

let sum = t;

let gammaPrevious: number;

let R: number;

for (let m = 1; math.abs(t / sum) \textgreater{} eps; m++) \{

gammaPrevious = gammaCurrent;

gammaArg += 1;

gammaCurrent = this.math.gamma(gammaArg);

R = -(xHalfSqr * gammaPrevious) / (m * gammaCurrent);

t *= R;

sum += t;

\}

return sum;

\}

Као што се може видети, иако имплементирају исти алгоритам, метод
\emph{calculate} је доста читљивији и интуитивнији од метода
\emph{calculate64. З}бог тога и јесте доста једноставније прво
имплементирати метод \emph{calculate,} па потом тај код ''пребацити'' у
\emph{calculate64.}

\begin{itemize}
\item
  \protect\phantomsection\label{_Toc44}{}Гама функција
\end{itemize}

При имплементацији апроксимације вредности Беселове функције, као и бета
функције, коришћени су методи класе \emph{GammaFunction.} Стога је
неопходно искористити што прецизнији алгоритам да се не би губило на
прецизности у алгоритмима поменутих функција\emph{.}

У поглављу 3 у коме је дефинисана гама функција изложене су и бројне
често коришћене апроксимације, а као најпрецизнија наведена је тзв.
\emph{Ланцошева апроксимација}. При имплементацији Гама функције су пре
саме Ланцошеве имплементације тестиране неке апроксимације изложене у
трећем поглављу: варијанте Стирлингове апроксимације, побољшане верзије
Рамануџанове формуле итд. Међутим ни једна од њих није дала
задовољавајуће резултате, тако да је као најпрецизнија позната
имплементирана Ланцошева апроксимација, упркос томе што је доста
комплекснија од наведених формула.

Проблематичан део ове апроксимације јесте израчунавање њених
коефицијената који су дати формулом

\(a_{k}(r) =\)\({( - 1)}^{k}\)\(\sqrt{\frac{2}{\pi} \cdot}e^{r} \cdot k \cdot \overset{k}{\sum_{j = 0}}{( - 1)}^{j} \cdot \frac{(k + j - 1)!}{(k - j)! \cdot j!}\left( \frac{e}{j + r + \frac{1}{2}} \right)^{j + \frac{1}{2}}\).

Пошто знак сваког члана суме зависи од парности индекса, ова сума јако
споро конвергира и једноставан алгоритам написан по овој формули (слично
попут алгоритма за Беселову фунцију) не даје резултате - програм уђе у
практично бесконачну петљу и извршавање се прекине после неког времена.

Због тога је имплементиран алгоритам који доста ефикасно израчунава
Ланцошеве коефицијенте, а који предлаже Пол Годфри {[}18{]}. Ланцош
заправо апроксимира вредност природног логаритма функције гама формулом:

\(\ln\Gamma(x + 1) = \ln(Z \cdot P) + (x + 0.5) \cdot \ln(x + g + 0.5) - (x + g + 0.5)\),

при чему је \(Z\) вектор добијен као:

\(Z = \left\lbrack 1,\frac{1}{x + 1},\frac{1}{x + 2},...,\frac{1}{x + n - 1} \right\rbrack\),

\(P\) вектор добијен множењем матрица \(P = D \cdot B \cdot C \cdot F\)
(дефинисане у наставку), а \(n\) природан и \(g\)реалан параметар. Избор
ових параметара значајно утиче на добијену прецизност.

Матрице \(B,C,D,F\) које дефинишу вектор \(P\) се добијају по следећим
формулама:

\(B_{ij} = \left\{ \begin{array}{r}
1,i = 0 \\
{( - 1)}^{i - j}\left( \begin{array}{r}
i + j - 1 \\
j - i
\end{array} \right),i > 0,j \geq i \\
0,иначе
\end{array} \right.\ \),

\(C_{ij} = \left\{ \begin{array}{r}
\frac{1}{2},i = j = 0 \\
0,j > i \\
{( - 1)}^{i - j}\overset{i}{\sum_{k = 0}}\left( \begin{array}{r}
2i \\
2k
\end{array} \right)\left( \begin{array}{r}
k \\
k + j - i
\end{array} \right),иначе
\end{array} \right.\ \),

\(D_{ij} = \left\{ \begin{array}{r}
0,i \neq j \\
1,i = j = 0 \\
 - 1,i = j = 1 \\
\frac{2(2i - 1) \cdot D_{i - 1,j - 1}}{i - 1},иначе
\end{array} \right.\ \),

\(F_{i} = \frac{(2i)! \cdot e^{i + g + 0.5}}{i! \cdot 2^{2i - 1} \cdot {(i + g + 0.5)}^{i + 0.5}}\).

\emph{C} представља матрицу коефицијената за Чебишевљеве полиноме и она
се може рачунати и рекурзивно.

Будући да матрице све зависе само од параметара \(n\) и \(g\), ако
одаберемо ове параметре тако да буду фиксни, матрице се могу израчунати
на самом покретању програма, при иницијализацији, и тако можемо смањити
време рада самог алгоритма јер нема потребе да се израчунавају сваки
пут. Самим тим и матрица \emph{P} која је производ ових матрица је
фиксна и она се такође може израчунати у иницијализацији.

Неке имплементације Ланцошеве апроксимације које се могу наћи због тога
имају сачуван вектор P као константу и уопште не израчунавају његову
вредност већ користе фиксне вредености што имплементацију своди само на
израчунавање једноставног израза.

Као што је већ споменуто, избор параметара је јако битан јер он одређује
прецизност апроксимације. Најсигурнији начин за избор адекватних
параметара је написати програм који ће да изврши претрагу за које
комбинације параметара је максимална грешка најмања. Годфри {[}18{]}
грешку даје формулом:

\(|\varepsilon| = \left| \frac{\pi}{2\sqrt{2e}}\left( e^{g}\sqrt{\pi} - \overset{n - 1}{\sum_{n = 0}}{( - 1)}^{i}E_{i} \right) \right|\),

при чему се вектор Е рачуна по формули: \(E = C \cdot F\). Будући да је
\(g\) реалан параметар и да треба да се испитају комбинације пара
\((n,g)\), ово испитвање је доста захтевно и у овом раду није рађено,
већ је испробано неколико познатих парова \((n,g)\) који су дати у
табели 13.4.1. и на крају је изабран онај пар који даје највише
прецизних цифара.

{\def\LTcaptype{none} % do not increment counter
\begin{longtable}[]{@{}
  >{\centering\arraybackslash}p{(\linewidth - 4\tabcolsep) * \real{0.1780}}
  >{\centering\arraybackslash}p{(\linewidth - 4\tabcolsep) * \real{0.4110}}
  >{\centering\arraybackslash}p{(\linewidth - 4\tabcolsep) * \real{0.4109}}@{}}
\toprule\noalign{}
\endhead
\bottomrule\noalign{}
\endlastfoot
\(n\) & \(g\) & \emph{број тачних децимала иза зареза} \\
13 & 6.024680040776729583740234375 & 49 \\
13 & 13.144565 & 50 \\
17 & 12.2252227365970611572265625 & 13 \\
21 & 20.1200911879863739013679875 & 9 \\
\end{longtable}
}

Табела 13.4.1. Парови параметара који дају најбоље апроксимације.

Од наведених парова из табеле 12.3.1. најбољи резултати добијени су за
\(n = 13\) и \(g = 13.144565\), а за само једну децималу лошији резултат
је за \(n = 13\) и\(g = 6.024680040776729583740234375\). За одабране
вредности апроксимација даје, нпр.

\(\Gamma(1) = 1.0000000000000000000000000000000000000000000000000015983089672\),

\(\Gamma(2) = 1.00000000000000000000000000000000000000000000000000786292157837\),

\(\Gamma(3) = 2.000000000000000000000000000000000000000000000000044728092714001\),

\(\Gamma(4) = 6.000000000000000000000000000000000000000000000000289665138806081\),

\(\Gamma(5) = 24.00000000000000000000000000000000000000000000000210901915273219\),

што даје чак 50 тачних децимала иза зареза.

calculate64(params: FunctionParamsForCalculationWithBigNumbers): string
\{

let \{ x \} = this.stringToBigNumber(params);

console.log(x.toString());

x = this.math.subtract(x, BIG\_NUMBER\_CONSTANTS.ONE);

// P = D * B * C * F

let P = this.math.multiply(this.D, this.B);

P = this.math.multiply(P, this.C);

P = this.math.multiply(P, this.F);

const vectorZ: BigNumber{[}{]} = this.populateVectorZ(x);

let R = this.math.multiply(vectorZ, P); // Z * P

let lnZP = this.math.log(R{[}0{]}); // ln(Z * P)

// x+0.5

let xPlusHalf = this.math.add(x, this.math.bignumber(0.5));

// x+0.5+g

let t = this.math.add(xPlusHalf, this.math.bignumber(this.g));

// (x+0.5)* ln(x+0.5+g)

let t1 = this.math.multiply(xPlusHalf, this.math.log(t));

// ln(G) = ln(Z * P) + (x + 0.5)*ln(x + g + 0.5)-(x + g + 0.5)

let logGx = this.math.add(lnZP, t1);

logGx = this.math.subtract(logGx, t);

const Gx = this.math.exp(logGx as BigNumber);

return Gx.toString();

\}

Матрице \emph{D, B, C, F} су иницијализоване у конструктору класе.
Вектор \emph{Z} зависи од унете променљиве, тако да он не може да буде
израчунат када и ове матрице, већ се метод за његово израчунавање позива
у самом методу \emph{calculate64}.

Кao што је претходно описано, рачуна се \emph{P} множењем матрица, потом
се добијени производ множи вектором \emph{Z} и израчунава се природни
логаритам добијене вредности. Потом се рачунају изрази
\((x + 0.5) \cdot \ln(x + g + 0.5)\) и \((x + g + 0.5)\) и логаритам се
сабира с првим изразом и од збира се одузима други израз. Добијено
решење представља вредност природног логаритма \(\Gamma(x + 1)\), тако
да је потребно израчунати експоненцијалну функцију добијене вредности да
би се добила вредност \(\Gamma(x + 1)\).

\begin{itemize}
\item
  \protect\phantomsection\label{_Toc45}{}Бета функција
\end{itemize}

Бета функција се рачуна помоћу гама функције као

\(B(p,q) =\)\(\frac{\Gamma(p) \cdot \Gamma(q)}{\Gamma(p + q)}\),

па се самим тим и имплементација своди на рачунање вредности гама
функције за три вредности и израчунавања израза. Због тога класа
\emph{BetaFunction} у конструктору иницијализује члан типа
\emph{GammaFunction.}

calculate64(params: FunctionParamsForCalculationWithBigNumbers): string
\{

const \{ x, y \} = this.stringToBigNumber(params);

// gamma(x)

const Gx = this.math.bignumber(

this.gamma.calculate64(\{

...initializeParams64(),

xBig: x.toString(),

\})

);

// gamma(y)

const Gy = this.math.bignumber(

this.gamma.calculate64(\{

...initializeParams64(),

xBig: y.toString(),

\})

);

// gamma(x+y)

const Gxy = this.math.bignumber(

this.gamma.calculate64(\{

...initializeParams64(),

xBig: this.math.add(x, y).toString(),

\})

);

// G(x) * G(y)

const mul = this.math.multiply(Gx, Gy);

// G(x) * G(y)/G(x+y)

return this.math.divide(mul, Gxy).toString();

\}

Навешћемо један пример. Резултат рада функције за \emph{p=6, 1=4} je:

,

\(B(5,4) = 0.00357142857142857142857142857142857142857142857142765043458645275\)

док за исте вредности \emph{WolframAlpha} даје број

0.00357142857142857142857142857142857142857142857142765043458645275,

што је 49 истих децимала иза децималне тачке.

\begin{itemize}
\item
  \protect\phantomsection\label{_Toc46}{}Лежандров полином
\end{itemize}

Дефиницијом 23 дат је израз
\(P_{n}(x) =\)\(\overset{\lbrack n/2\rbrack}{\sum_{k = 0}}{( - 1)}^{k}\frac{(2n - 2k)!}{2^{n}k!(n - k)!(n - 2k)!}x^{n - 2k}\)
који се може употребити за имплементацију евалуације вредности
Лежандровог полинома за задату вредност реда и аргумента. Као и код
Лагерових полинома и Беселове функције, изводи се однос између
узастопних сабирака ради оптимизације алгоритма, па се добија:

\(t_{k} = {( - 1)}^{k}\frac{(2n - 2k)!}{2^{n}k!(n - k)!(n - 2k)!}x^{n - 2k} \Rightarrow t_{0} = {(x - 1)}^{n}\),

\[R_{k + 1} = \frac{\left( \begin{array}{r}
n \\
k + 1
\end{array} \right)^{2}{(x - 1)}^{n - k - 1}{(x + 1)}^{k + 1}}{\left( \begin{array}{r}
n \\
k
\end{array} \right)^{2}{(x - 1)}^{n - k}{(x + 1)}^{k}} =\]

\(=\)
\(\frac{\frac{n^{2}}{{(n - k - 1)}^{2}{(k + 1)}^{2}} \cdot {(x - 1)}^{n - k - 1}{(x + 1)}^{k + 1}}{\frac{n^{2}}{{(n - k)}^{2}k^{2}} \cdot {(x - 1)}^{n - k}{(x + 1)}^{k}}\)\(= \left( \frac{n - k}{k + 1} \right)^{2}{(x - 1)}^{- 1}\)\((x + 1)\)\(=\)\(= \left( \frac{n - k}{k + 1} \right)^{2}\frac{x + 1}{x - 1}\).

calculate64(params: FunctionParamsForCalculationWithBigNumbers): string
\{

const \{ alpha, x \} = this.stringToBigNumber(params);

const xMinus1 = this.math.subtract(x, BIG\_NUMBER\_CONSTANTS.ONE);

const xPlus1 = this.math.add(x, BIG\_NUMBER\_CONSTANTS.ONE);

let t = this.math.pow(xMinus1, alpha), sum = t, R: MathType;

const Rx = this.math.divide(xPlus1, xMinus1); // (x+1)/(x-1)

for (

let k = 1;

Number(this.math.compare(this.math.bignumber(k), alpha)) \textless= 0;

k++

) \{

const kBig = this.math.bignumber(k);

R = this.math.subtract(alpha, kBig); // n-k

R = this.math.add(R, BIG\_NUMBER\_CONSTANTS.ONE); // k+1

R = this.math.divide(R, kBig); // (n-k)/(k+1)

R = this.math.pow(R, BIG\_NUMBER\_CONSTANTS.TWO); // ((n-k)/(k+1))\^{}2

R = this.math.multiply(R, Rx);

t = this.math.multiply(t, R); // t = t*R

sum = this.math.add(sum, t); // sum = sum+t

\}

let res = this.math.pow(

BIG\_NUMBER\_CONSTANTS.TWO,

this.math.unaryMinus(alpha)

);

res = this.math.multiply(res, sum);

return res.toString();

\}

\begin{itemize}
\item
  \protect\phantomsection\label{_Toc47}{}Лагеров полином
\end{itemize}

Израз из дефиниције 28 ће бити искоришћен за имплементацију Лагерових
полинома:

\(L_{n}(x) =\)\(\overset{n}{\sum_{k = 0}}\frac{{( - 1)}^{k}}{k!}\left( \begin{array}{r}
n \\
k
\end{array} \right)x^{k}\).

Ако пробамо да применимо оптимизацију као код Беселових функција, добија
се:

\(t_{k} = \frac{{( - 1)}^{k}}{k!}\left( \begin{array}{r}
n \\
k
\end{array} \right)x^{k}\) \(\Rightarrow t_{0} = 1\),

\[R_{k + 1} = \frac{\frac{{( - 1)}^{k + 1}}{(k + 1)!} \cdot \frac{n!}{(k + 1)! \cdot (n - k - 1)!} \cdot x^{k + 1}}{\frac{{( - 1)}^{k}}{k!} \cdot \frac{n!}{k! \cdot (n - k)!} \cdot x^{k}}\]

\(= \frac{- 1}{k + 1} \cdot \frac{k! \cdot (n - k) \cdot (n - k - 1)!}{(k + 1) \cdot k! \cdot (n - k - 1)!}\)\(= - \frac{n - k}{{(k + 1)}^{2}}\).

Будући да је именилац у сваком случају дефинисан, ова оптимизација јесте
применљива за Лагеров полином. Још један могућ приступ је рекурзивни,
будући да важи једнакост
\((n + 1)L_{n + 1}(x) - (2n + 1 - x)L_{n}(x) + nL_{n - 1}(x) = 0\), па
би рекурзивна функција подразумевала постављање иницијалних вредности:
\(L_{0}(x) = 1\), \(L_{1}(x) = - x + 1\) и рекурзиван

позив јер је
\(L_{n}(x) =\)\(\frac{(2n - 1 - x)L_{n - 1}(x) + (n - 1)L_{n - 2}(x)}{n}\).

calculate64(params: FunctionParamsForCalculationWithBigNumbers): string
\{

const \{ alpha, x \} = this.stringToBigNumber(params);

const alphaPlus1 = this.math.add(alpha, BIG\_NUMBER\_CONSTANTS.ONE);

let t: MathType = this.math.bignumber(BIG\_NUMBER\_CONSTANTS.ONE);

let sum: MathType = t, R: MathType, kBig: MathType;

for (

let k = 1;

Number(this.math.compare(this.math.bignumber(k), alpha)) \textless= 0;
k++

) \{

kBig = this.math.bignumber(k); // needs explicit conversion

R = this.math.subtract(alphaPlus1, k); //(n+1-k)

R = this.math.divide(R, this.math.pow(kBig,
BIG\_NUMBER\_CONSTANTS.TWO));

R = this.math.multiply(R, x); // (n+1-k)*x/(k\^{}2)

R = this.math.unaryMinus(R); // -(n+1-k)*x/(k\^{}2)

t = this.math.multiply(t, R); t=t*R

sum = this.math.add(sum, t); // s=s+t

\}

return sum.toString();

\}

\begin{itemize}
\item
  \textbf{Ермитови полиноми}
\end{itemize}

Ермитов пробабилистички полином се дефинише преко релације с физичким
полиномом, тако да посао имплементације калкулатора за ова два полинома
фактички спада само на физички полином.

Запис Ермитовог физичког полинома употребљен за имплементацију је
изведен из дефиниције 32:

\(H_{n}(x) =\)\(n!\overset{\lbrack n/2\rbrack}{\sum_{k = 0}}\frac{{( - 1)}^{k}}{k!(n - 2k)!}{(2x)}^{n - 2k}\).

Дакле,
\(t_{k} = \frac{{( - 1)}^{k}}{(n - 2k)! \cdot k!}{(2x)}^{n - 2k}\),
одакле се долази до \(t_{0} = \frac{1}{n!}{(2x)}^{n}\) и, коначно,

\(R_{k + 1} = - {(2x)}^{- 2}\frac{(n - 2k - 1)(n - 2k)}{k + 1}\).

Проблем са овим јесте што, иако је Ермитов полином дефинисан у свакој
тачки, \(R_{k}\) није дефинисано у нули због \({(2x)}^{- 2}\). Будући да
је сума код Ермитовог полинома коначна за разлику од Беселове функције и
она је сразмерна реду полинома, перформансе ће бити прихватљиве и ако
једноставно у сваком кораку израчунамо \(t_{k}\), тако да нема потребе
користити два различита алгоритма у зависности од вредности променљиве.

calculate64(params: FunctionParamsForCalculationWithBigNumbers): string
\{

const \{ alpha, x \} = this.stringToBigNumber(params);

const alphaFactorial = this.math.factorial(alpha);

const x2 = this.math.multiply(BIG\_NUMBER\_CONSTANTS.TWO, x);

let sum: math.MathType = 0,k = 0, t, n2k, m, den;

while (Number(this.math.compare(this.math.multiply(2, k), alpha))
\textless= 0) \{

n2k = this.math.subtract(

alpha,

this.math.multiply(BIG\_NUMBER\_CONSTANTS.TWO, k)

);

// m = (-1)\^{}k

m = this.math.pow(-1, k);

// k! * (n-2k)!

den = this.math.multiply(

this.math.factorial(k),

this.math.factorial(n2k as math.BigNumber)

);

// t=(-1)\^{}k * (2x)\^{}(n-2k)

t = this.math.multiply(m, this.math.pow(x2, n2k as math.BigNumber));

t = this.math.divide(t, den); // t = t/(k! * (n-2k)!)

// sum = sum+t

sum = this.math.add(sum, t);

k++;

\}

// sum = n! * sum

sum = this.math.multiply(alphaFactorial, sum);

return sum.toString();

\}

calculate64(params: FunctionParamsForCalculationWithBigNumbers): string
\{

const \{ alpha, x \} = this.stringToBigNumber(params);

const sqrt2 = this.math.sqrt(BIG\_NUMBER\_CONSTANTS.TWO);

const xPhy = this.math.divide(x, sqrt2); // x/sqrt(2)

// yPhisicist = Hn(x/sqrt(t))

const yPhy = this.hermitePhysicist.calculate64(\{

...initializeParams64(),

alphaBig: params.alphaBig,

xBig: xPhy.toString(),

epsBig: params.epsBig,

\});

// 2\^{}(-n/2) * Hn(x/sqrt(2))

const factor = this.math.pow(

BIG\_NUMBER\_CONSTANTS.TWO,

this.math.divide(

this.math.unaryMinus(alpha),

BIG\_NUMBER\_CONSTANTS.TWO

) as math.BigNumber

);

const solution = this.math.multiply(this.math.bignumber(yPhy), factor);

return solution.toString();

\}

\begin{itemize}
\item
  \protect\phantomsection\label{_Toc48}{}Чебишевљеви полиноми
\end{itemize}

За Чебишевљеве полиноме прве и друге вресте имплементација је врло
једноставна будући да су њихове дефиниције дате кроз тригонометријске
функције:

\(T_{n}(x) =\)\(\cos(n\arccos x),|x| \leq 1\),

\(U_{n}(x) =\)\(\frac{\sin((n + 1)x)}{\sin x}\).

calculate64(params: FunctionParamsForCalculationWithBigNumbers): string
\{

const \{ alpha, x \} = this.stringToBigNumber(params);

let result: MathType = this.math.acos(x);

result = this.math.multiply(result, alpha); // n * arccos(x)

result = this.math.cos(result as BigNumber); // cos(n * arccos(x))

return result.toString();

\}

calculate64(params: FunctionParamsForCalculationWithBigNumbers): string
\{

const \{ alphaBig, xBig \} = params;

const alpha = this.math.bignumber(alphaBig);

const arccosX = this.math.acos(this.math.bignumber(xBig));

let result: MathType = this.math.add(alpha, BIG\_NUMBER\_CONSTANTS.ONE);

result = this.math.multiply(result, arccosX);

result = this.math.sin(result as BigNumber);

result = this.math.divide(result, this.math.sin(arccosX));

return result.toString();

\}

\begin{itemize}
\item
  \protect\phantomsection\label{_Toc49}{}Јакобијеви полиноми
\end{itemize}

Израчунавање вредности Јакобијевих полинома може да се имплементира
користећи следећи облик:

\(P_{n}^{(\alpha,\beta)}\)\((x) =\)\(\frac{1}{2^{n}}\overset{n}{\sum_{k = 0}}\left( \begin{array}{r}
n + \alpha \\
k
\end{array} \right)\left( \begin{array}{r}
n + \beta \\
n - k
\end{array} \right){(x - 1)}^{n - k}{(1 + x)}^{k}\),
\(n \in {\mathbb{N}}\),

изведен у поглављу 10.2.

применом Лајбницовог правила на Родригову формулу.

Ако ради оптимизације изведемо однос сваког сабирка са претходним,
добићемо:

\[t_{k} = \left( \begin{array}{r}
n + \alpha \\
k
\end{array} \right)\left( \begin{array}{r}
n + \beta \\
n - k
\end{array} \right){(x - 1)}^{n - k}{(1 + x)}^{k}\]\[\Rightarrow t_{0} = (n + \alpha)(n + \beta){(x - 1)}^{n}\]

\[R_{k + 1} = \frac{\left( \begin{array}{r}
n + \alpha \\
k + 1
\end{array} \right)\left( \begin{array}{r}
n + \beta \\
n - k - 1
\end{array} \right){(x - 1)}^{n - k - 1}{(1 + x)}^{k + 1}}{\left( \begin{array}{r}
n + \alpha \\
k
\end{array} \right)\left( \begin{array}{r}
n + \beta \\
n - k
\end{array} \right){(x - 1)}^{n - k}{(1 + x)}^{k}} =\]

\[= \frac{\frac{(n + \alpha)!}{(n + \alpha - k - 1)! \cdot (k + 1)!} \cdot \frac{(n + \beta)!}{(n + \beta - n + k + 1)! \cdot (n - k - 1)!}}{\frac{(n + \alpha)!}{(n + \alpha - k)! \cdot k!} \cdot \frac{(n + \beta)!}{(n + \beta - n + k)! \cdot (n - k)!}}\]\[\cdot \frac{x + 1}{x - 1}\]\[=\]

\(= \frac{n + \alpha - k}{k + 1} \cdot \frac{n - k}{\beta - k} \cdot \frac{x + 1}{x - 1}\).

calculate64(params: FunctionParamsForCalculationWithBigNumbers): string
\{

const \{ alpha, x, a, b \} = this.stringToBigNumber(params);

const xPlus1 = this.math.add(x, BIG\_NUMBER\_CONSTANTS.ONE);

const xMinus1 = this.math.subtract(x, BIG\_NUMBER\_CONSTANTS.ONE);

const alphaPlusA = this.math.add(alpha, a);

const alphaPlusB = this.math.add(alpha, b);

let sum = BIG\_NUMBER\_CONSTANTS.ZERO;

let t1, t2, t, kBig: MathType;

for (

let k = 0;

Number(this.math.compare(this.math.bignumber(k), alpha)) \textless= 0;

k++

) \{

kBig = this.math.bignumber(k);

const nk = this.math.subtract(alpha, kBig); // n-k

const bca = binomialCoefficient64(alphaPlusA, kBig); // n+a over k

const bcb = binomialCoefficient64(alphaPlusB, nk); // n+b over n-k

const factorNK = this.math.pow(xMinus1, nk); // (x-1)\^{}(n-k)

const factorK = this.math.pow(xPlus1, kBig); // (x+1)\^{}k

t1 = this.math.multiply(bca, bcb);

t2 = this.math.multiply(factorNK, factorK);

t = this.math.multiply(t1, t2);

sum = this.math.add(sum, t) as BigNumber;

\}

const res = this.math.divide(

sum,

this.math.pow(BIG\_NUMBER\_CONSTANTS.TWO, alpha)

);

return res.toString();

\}

\}

Петља која итерира од 0 до \emph{n} израчунава сабирке и сумира их:
израчунавају се биномни коефицијенти \(\left( \begin{array}{r}
n + \alpha \\
k
\end{array} \right)\) и \(\left( \begin{array}{r}
n + \beta \\
n - k
\end{array} \right)\) коришћењем помоћне функције
\emph{binomialCoefficient64,} множе се одговарајућим степенима \(x - 1\)
и \(1 + x\), који су пре уласка у петљу иницијализовани као константе и
све то се множи. Као константе су такође пре петље израчунати и
\(n + \alpha\) и \(n + \beta\). На крају се цела сума дели с \(2^{n}\).

\begin{itemize}
\item
  \protect\phantomsection\label{_Toc50}{}Ограничења апликације
\end{itemize}

Приликом коришћења апликације \emph{Калкулатор неких специјалних
функција} потребно је имати у виду и одређена ограничења.

\begin{itemize}
\item
  Нумеричка прецизност -- апликација подржава рад са бројевима до 64
  значајне цифре, али није гарантовано да ће све цифре бити поуздане
  услед ограничења апроксимационих метода. На примеру гама функције смо
  видели да се коришћењем најбоље нумеричке апроксимације постиже до 50
  децимала иза зареза, док се за бета функцију, која се израчунава
  коришћењем вредности гама функције, прецизност додатно умањује услед
  акумулирања грешке.
\item
  Конвергенција приликом израчунавања вредности функција -- брзина и
  стабилност нумеричке апроксимације зависе од својстава одабране
  функције и вредности унетих параметара. За функције које споро
  конвергирају, време израчунавања расте, а код великих вредности
  параметара конвергенција може бити још спорија, што утиче на
  ефикасност апликације. Пример функција код којих се то најбоље
  примећује јесу Беселове функције или Лагерови полиноми -- са повећањем
  аргумената, поготову реда, повећава се и број потребних корака за
  апроксимацију.
\item
  Ограничења улаза -- улазна поља имају максималну дужину дефинисану у
  оквиру апликације, што логично спречава унос превеликих или
  неподржаних вредности.
\item
  Ресурси -- као што је могло да се види у алгоритму за конверзију
  инфиксног израза у постфиксни и алгоритму за евалуацију вредности
  постфиксног израза, оба раде са стеком, тако да се као природно
  ограничење намеће величина стека. Величина стека је такође ограничење
  и за нумеричке методе, будући да многе од њих користе рекурзивни
  приступ. За комплексније изразе и високе редове полинома потребно је
  алоцирати довољно меморијских ресурса.
\end{itemize}

\begin{itemize}
\item
  \protect\phantomsection\label{_Toc51}{}Поређење са постојећим
  програмским пакетима исте намене
\end{itemize}

Увидом у постојеће калкулаторе специјалних функција, долази се до
неколико веб апликација које су по доста критеријума другачије од
апликације која је предмет овог рада. У овом поглављу ће бити дато
поређење тих апликација са предметном апликацијом.

Један од таквих алата је \emph{WolframAlpha
(\ul{\mbox{\url{https://www.wolframalpha.com/}})}.} Овај софтвер постоји
већ дуги низ година и временом је проширен на далеко већи број области
од само математике, па тако данас има и одељке: Науке и технологија,
Друштво и култура, Свакодневни живот. У том смислу, можда поређење са
овом апликацијом није најадекватније, будући да је \emph{WolframAlpha}
огромна апликација с великим бројем примена, те у контексту овог рада
говоримо само о израчунавању специјалних функција које представља само
делић овог програмског пакета. У погледу специјалних функција, до
калкулатора нпр. Беселове функције може се доћи са почетне странице
једноставном претрагом кључних речи. Тада апликација води на страницу
која нуди унос променљиве, реда (параметра) функције и падајући мени за
избор врсте Беселове функије, као на наредној слици
14.1.\includegraphics[width=6.30683in,height=5.74204in,alt={pasted-movie.png}]{media/image22.png}

Слика 14.1: Калкулатор Беселових функција на алату \emph{WolframAlpha.}

Осим Беселове функције прве врсте, овај алат нуди могућност израчунавања
вредности и Беселових функција друге и треће врсте, сферне функције,
итд. Такође, одмах испод резултата су приказане и различите
репрезентације ове функције.

Оно што се може узети као мана је то, што да би се уопште дошло до овог
калкулатора, мора да се изврши претрага (ако корисник не зна да уопште
постоји ова опција можда је и неће претражити већ ће ручно унети
функцију), а када је израчуната вредност, прикаже се резултат са само
шест децимала; да би се добило више децимала мора да се кликне на дугме.
Такође, у овом резултату нема графичког приказа, тако да уколико би
корисник желео то, поново мора да предузима одређене акције. Осим што
корисник мора да обави већи број акција да би дошао до жељеног
резултата, што може бити фрустрирајуће и одузима више времена, може се
рећи да ово и није претерано интуитивно. Наравно, то је последица
чињенице да специјалне функције уопште и нису превасходни циљ тог алата,
већ само једна од могућности које он пружа.

С друге стране, када се претраже све друге специјалне функције које нуди
\emph{Калкулатор специјалних функција}, резултат претраге неће бити
калкулатор специјалне функције већ репетиторијумске информације о самој
функцији - различите дефиниције, графици за првих неколико вредности
реда, интеграл и извод функције, тако да ако би корисник имао потребу да
израчуна вредности неке од њих, то би могао да уради само мануелним
уносом функције.

Када се говори о \emph{WolframAlpha,} неизбежно је споменути још један
производ исте компаније (\emph{Wolfram Research) ---} комерцијални
програмски пакет \emph{Mathematica,} који је у употреби још од 1988.
године {[}54{]}. \emph{Mathematica} је веома популаран алат широке
намене који је практично зачетник модерног симболичког израчунавања.
Важност овог алата огледа се и у чињеници да је то први алат који садржи
толики опсег функционалности: нумеричка, алгебарска, графичка
израчунавања, и друге, који су се до тада могли наћи само као
индивидуални пакети намењени за неку специфичну примену. Кроз протекле
три и по деценије, \emph{Wolfram Mathematica} постаје стандардни алат за
истраживања и производњу у инжењерству, представља основу истраживања у
бројним научним радовима и решава одређене теоријске проблеме и питања,
добија важну улогу и у развоју софтвера и информатици, а користи се и за
финансијско моделовање, анализу и планирање {[}43{]}.

Такође комерцијалнo симболичко и нумеричко окружење којe опстаје и дуже
од пакета \emph{Mathematica} је \emph{Maple,} развијан на Вотерлу
универзитету у Канади почетком осамдесетих година прошлог века, с циљем
да се направи алат који би могао да ради на јефтиним корисничким
рачунарима, будући да су за постојеће програме предуслов били хардверски
веома захтевни рачунари. Крајем осамдесетих година алат је
комерцијализован. Власник је \emph{Maplesoft™}. Он пружа подршку за
нумеричке прорачуне, визуелизацију, симболичка израчунавања и омогућава
декларисану произвољну прецизност. Корисник се може користити
стандардном математичком нотацијом, али и програмским језиком {[}28{]}.

Интересантно је када говоримо о овим софтверским алатима осврнути се и
на рад \emph{``Actual computer algebra systems and their future
development'',} који се управо бави анализом софтвера исте намене,
актуелних до краја двадесетог века, те анализира њихове карактеристике,
проблеме и прогнозира њихов даљи развој {[}21{]}.

Још један веома популаран алат, који се појављује као први резултат
претраге, јесте \emph{Symbolab
(\ul{https://www.symbolab.com/solver/functions-calculator}).} Слично
претходним описаним алатима, и \emph{Symbolab} је алат који задовољава
различите потребе корисника из области математике, физике, хемије,
економије, итд. Битан део ове апликације јесу управо функције. Међутим,
тај алат не нуди опције за конкретно специјалне функције, већ корисник
мора мануално да унесе функцију чије израчунавање му је потребно.

Овај алат као резултат рада даје неке основне особине функције и график,
али одређене битне могућности алата су доступне само у комерцијалној
верзији алата. Такође, он неће заправо израчунати неку задату вредност,
већ у бесплатној верзији даје само особине унете функције и график, иако
се овај део апликације зове
''\href{https://www.symbolab.com/solver/functions-line-calculator}{\emph{Functions
\& Line Calculator}}'', а вредност се може израчунати у потпуно другом
делу апликације.

На слици испод је калкулатор који нуди
\emph{Symbolab.}\includegraphics[width=6.26223in,height=3.51429in,alt={pasted-movie.png}]{media/image23.png}

Слика 14.2: Калкулатор функција апликације \emph{Symbolab.}

Апликација која је по намени најсличнија апликацији из овог рада је
\emph{Special Functions Calculator
(\ul{http://www.mhtl.uwaterloo.ca/old/onlinetools/func\_calc/func\_calc.html})}.
Ова апликација је врло интуитивна и једноставна, али нуди свега 4
функције: Гаусову функцију грешке, бета и гама функције и Беселову
функцију. Даје информацију за до 6 значајних цифара и резултате за прва
два реда Беселове функције прве и друге врсте и модификоване Беселове
функције друге врсте. Ово је уједно и једина функционалност коју ова
апликација нуди - не постоје графици функција, особине, итд. Апликација
је, у ствари, некомерцијални део програмског пакета
\emph{Maple.}\includegraphics[width=2.78976in,height=2.16063in,alt={pasted-movie.png}]{media/image24.png}

Слика 14.3: Рад апликације \emph{Special Functions Calculator.}

Веб апликација \emph{123calculus}
(\emph{\ul{https://www.123calculus.com/en}}) покрива сличне области као
\emph{Symbolab,} али у контексту специјалних функција има функционалност
сличну главној функционалности исхода овог рада, тј. за одабрану врсту
специјалне функције и унете параметре даје резултат њене вредности.
Међутим, она нуди потпуно други избор функција у односу на апликацију
која је предмет овог рада - Гама функцију, функцију грешке и различите
елиптичке интеграле.

После детаљног увида у предметни софтвер, може се закључити да остатак
бесплатних веб апликација које врше израчунавања специјалних функција
углавном се своде на калкулаторе појединачних функција (најчешће су то
калкулатори Беселових функција) и углавном им је то и једина
функционалност.

\protect\phantomsection\label{_Toc52}{}Закључак

Специјалне функције своју примену налазе у решавању многобројних
проблема у физици и математици, као и у техничким наукама. Самим тим,
јавља се и потреба за ефикасним, а ипак што једноставнијим начином за
израчунавањем вредности ових функција. Одатле и долази модерна идеја за
интуитивну апликацију добрих перформанси која би омогућила кориснику да
у што мање корака апроксимира вредност функције за очекивану задату
тачност.

Овим радом обухваћене су неке често коришћене специјалне функције:
Беселове функције (прве врсте), Лежандров и Лагеров полином, Ермитов
(физички и пробабилистички) полином, Јакобијеви полиноми, те Чебишевљеви
полиноми (прве и друге врсте).

У првом делу рада свака функција теоријски је обрађена почевши од
диференцијалне једначине чије је она партикулатно решење, дефинисана је
функција на свом домену ортогоналности, дате су њене различите
репрезентације, основне особине, рекурзивне везе, те бројне корисне
релације, као и релације између фунција различитих врста.

Потом је дат преглед апликације креиране у склопу овог рада, која врши
израчунавање одабране специјалне функције за задате вредности. Осим ове
функционалности, при калкулацији се добија и графички приказ функције, а
апликација нуди и репетиторијумски преглед дефиниција и особина ових
функција. Унос параметара функције, осим нумеричке вредности може бити и
израз, што је реализовано кроз алгоритме за конверзију инфиксног израза
у постфиксни и евалуацију постфиксног израза који су детаљно описани у
раду. Такође су дати и делови кода за евалуацију вредности сваке
специјалне функције обрађене у овом раду, што значи да је ова апликација
\emph{open source}.

На крају, изнесено је поређење са постојећим веб апликацијама које нуде
могућност израчунавања специјалних функција. Неколико популарних
апликација које нуде огроман спектар различитих израчунавања обухватају
и специјалне функције, али у мало другачијем формату - кроз више корака,
мање интуитиван интерфејс, итд. Постоје неки мање комерцијални алати
који нуде ову функционалност на сличан начин као предметна апликација,
али кроз минималистички интерфејс. Они обухватају другачији скуп
функција или чак само једну функцију, а дају и мању прецизност, не
поседују графички приказ, и слично. Као закључак произилази да је
предметна апликација јединствена са тако широким скупом функција и
функционалностима, да има веома интуитиван интерфејс и да корисник
једноставно долази до циља у односу на друге алате споменуте у раду.

Коначно, постоје неке не толико широко коришћене специјалне функције
чија би имплементација у ову апликацију била природно проширење, а што
би био правац проширења неких нових њених будућих верзија.

\protect\phantomsection\label{_Toc53}{}ЛИТЕРАТУРА

{[}1{]} Aboites, V., \emph{Laguerre polynomials and linear
algebra.}~Memorias Sociedad Matemática Mexicana~52 (2017); pp. 3-13.

{[}2{]} Abramowitz, M., Stegun I. A. (eds)~\emph{Handbook of
mathematical functions with formulas, graphs, and mathematical tables}.
Vol. 55. US Government printing office, 1948, p. 222, 333, 361.

{[}3{]} Aggarwal, S., Chauhan, R., Sharma, N., ``Mohand transform of
Bessel\RL{'}s functions".~\emph{International Journal of Research in
Advent Technology},~6(11), 2018, pp. 3034-3038.

{[}4{]} Andrews, G. E., Askey, R., Roy, R., \emph{Special functions}.
Cambridge: Cambridge university press, 1999.

{[}5{]} Arfken, G. B., Weber, H. J., Harris, F. E.,~\emph{Mathematical
methods for physicists: a comprehensive guide. Seventh Edition}.
Academic press, Elsevier, 2013, pp. 718-719.

{[}6{]} Artin, E.,~\emph{The gamma function}. Dover Publications, 2015.

{[}7{]} Asfeha, T., \emph{Legendre Differential Equation and its
applicаtion,} Department of mathematics, Addis Ababa University\emph{,}
2018\emph{.} available at
\href{http://213.55.95.56/handle/123456789/19254}{\ul{http://213.55.95.56/handle/123456789/19254}}

{[}8{]} Banahan, M., Brady, D., Doran, M., \emph{The C book}. GBdirect,
1991.

{[}9{]} Barkatou, M., ``Symbolic Methods for Solving Systems of Linear
Ordinary Differential Equations''\emph{, International Symposium on
Symbolic and Algebraic Computations (ISSAC)}, Munich, Germany (2010),
pp. 7-8.

{[}10{]} Bell, W. W.,~\emph{Special functions for scientists and
engineers}. Courier Corporation, 2004.

{[}11{]} Chen, C. P., \emph{``}A more accurate approximation for the
gamma function\emph{''.~Journal of Number Theory~164} (2016), pp.
417-428.

{[}12{]} Cuchta, T., Pavelites, M., Tinney, R., \emph{The Chebyshev
Difference Equation.}~Mathematics,~

8(1) (2020), p.74.

{[}13{]} Sanderson, C., Curtin, R., \emph{Armadillo: a template-based
C++ library for linear algebra.} Journal of Open Source Software 1 (2)
(2016), p. 26.

{[}14{]} David, C. W., \emph{Frobenius Solution for
Legendre\textquotesingle s Equation, Rodrigue\textquotesingle s Formula
and Normalization}, Chemistry Education Materials. (2009), p. 68.

{[}15{]} García, E. M., \emph{Bessel functions and equations of
mathematical physics}. Universidad del Pais Vasco, Facultad de Ciencia y
Tecnología, 2016.

{[}16{]} Erić, J., \emph{Neke Specijalne funkcije - diplomski rad}.
Sveučilište J. J. Strossmayera u Osijeku, Odjel za matematiku. Dostupno
na:
\href{https://www.mathos.unios.hr/~mdjumic/uploads/diplomski/ERI05.pdf}{\emph{\ul{https://www.mathos.unios.hr/\textasciitilde mdjumic/uploads/diplomski/ERI05.pdf}}},
2011.

{[}17{]} de Fériet, J. K., \emph{Fonctions de la physique mathematique}.
Centre national de la recherche scientifique, 1957, pp. 42-43.

{[}18{]} Godfrey, P., \emph{Lanczos Implementation of the Gamma
Function}, available at
\emph{http://www.numericana.com/answer/info/godfrey.htm (See also
http://my.fit.edu/}∼\emph{gabdo/gamma.txt})

{[}19{]} Hasan, S. Q., Abed, D. R., ``Fractional Operational Matrices
for Solving Multi-Fractional Nonlinear Differential Equations with Mixed
Boundary Conditions''.\emph{~Al-Mustansiriyah Journal of Science},~27
(2016), pp. 27-28.

{[}20{]} Imani, A., Aminataei, A., Imani, A., ``Collocation method via
Jacobi polynomials for solving nonlinear ordinary differential
equations''\emph{.}~\emph{International Journal of Mathematics and
Mathematical Sciences},~(2011), Article ID 673085, pp. 1-11.

{[}21{]} Iričanin, B., ``Actual computer algebra systems and their
future development'', in: \emph{6th World MultiConference on Systemics,
Cybernetics and Informatics, Vol V, Proceedings - Computer Sci I},
Orlando, U.S.A., 2002, vol. V, pp. 207-212.

{[}22{]} Korenev, B. G., \emph{Bessel functions and their applications}.
CRC Press, 2002.

{[}23{]} Kreh, M., \emph{Bessel functions.}~Lecture Notes, Penn
State-Göttingen Summer School on Number Theory~82 (2012), pp. 161-162.

{[}24{]} Lazić, S., Iričanin, B., "Calculator of some special
functions''\emph{,} in: \emph{The Book of Abstracts, XV Serbian
Mathematical Congress}, Belgrade, Serbia, 2024, p. 124.

{[}25{]} Lazić, S., Iričanin, B., ``Calculator of Some Special
Mathematical Functions''\emph{, Serbian Journal of Electrical
Engineering}, Vol 21 (2024), No 3, pp 345-357.

{[}26{]} Lee, K., Hubbard, S., \emph{Data structures and algorithms with
python}, ~Berlin/Heidelberg, Germany, Springer, 2015, pp 123-124.

{[}27{]} Lutz, M.,~\emph{Learning Python™: Powerful object-oriented
programming}. O\textquotesingle Reilly Media, Inc, 2013.

{[}28{]}~~~~ Maplesoft, a division of Waterloo Maple Inc.,
\emph{Maple~2024}. Waterloo, Ontario.

{[}29{]} Miličić, P. M., Ušćumlić, M. P.,~\emph{Zbirka zadataka iz više
matematike II}. VI izdanje. Naučna knjiga,1989, p. 375.

{[}30{]} Milosavljević, D., \emph{Rešavanje svojstvenog problema
Hamiltonijana: Numerov-Kulijev metod.} Prirodno matematički fakultet,
Univerzitet u Nišu, 2016.

{[}31{]} Milovanović, G., \emph{Numerička analiza i teorija
aproksimacija--Uvod u numeričke procese i rešavanje jednacina}.~Zavod za
udžbenike, Beograd, 2014.

{[}32{]} Mitrinović D., Janić R., \emph{Uvod u specijalne funkcije}.
Građevinska knjiga; 1986.

{[}33{]} Mortici, C., ``A continued fraction approximation of the gamma
function''. \emph{Journal of Mathematical Analysis and Applications},
402(2), (2013). pp. 405-410.

{[}34{]} Mortici, C., ``Gamma function estimates via completely
monotonicity arguments''\emph{.} \emph{Carpathian Journal of
Mathematics,} 28(1) (2012), pp. 93-102.

{[}35{]} Oldham, K. B., Myland, J., Spanier, J., \emph{An atlas of
functions: with equator, the atlas function calculator. Second edition}.
New York, NY, USA: Springer, 2009, pp. 56, 189-192, 217-221, 538,
567-570, 615.

{[}36{]} Olver, F. W. J. (ed.)~\emph{NIST handbook of mathematical
functions hardback and CD-ROM}. Cambridge university press, 2010, pp.
224.

{[}37{]} Patarroyo, K.Y., \emph{A digression on Hermite
polynomials.}~arXiv preprint arXiv:1901.01648, 2019.

{[}38{]} Penjić, S., \emph{Ortogonalni, Hermiteovi i Jacobijevi
polinomi}. Univerzitet u Sarajevu, Odsjek za matematiku, 2012, p. 29.

{[}39{]} Rivlin, T. J,.~\emph{Chebyshev polynomials}. Courier Dover
Publications, 2020.

{[}40{]} Schildt, H.,~\emph{Java: the complete reference}. McGraw-Hill
Education Group, 2014.

{[}41{]} Sebah, P., Gourdon, X., ``Introduction to the Gamma function''.
\emph{American Journal of Scientific Research}, Vol 2 (2002), pp. 2-18.

{[}42{]} Solter, N. A., Kleper, S. J.,~\emph{Professional C++}. John
Wiley \& Sons., 2005.

{[}43{]} Stanimirović, P. S., Milovanović, G. V., \emph{Programski paket
MATHEMATICA i primene}. Elektronski fakultet, Niš, 2002.

{[}44{]} Škrobar D.,~\emph{Specijalne funkcije. Diss.} Josip Juraj
Strossmayer University of Osijek. Department of Mathematics. Chair of
Pure Mathematics. Algebra and Calculus Research Group, 2016.

{[}45{]} Tomašević, M., \emph{Algoritmi i strukture podataka}. Akademska
misao,~Beograd. 2004, pp. 71-81.

{[}46{]} Temme, N. M., \emph{Special functions: An introduction to the
classical functions of mathematical physics}. John Wiley \& Son, 1996.

{[}47{]} Volmut, V.,~\emph{Ortogonalni polinomi}. Diss. Josip Juraj
Strossmayer University of Osijek. Department of Mathematics. Chair of
Pure Mathematics. Algebra and Calculus Research Group, 2016.

{[}48{]} Zwillinger,
\emph{D.,~\href{http://www.amazon.com/exec/obidos/ASIN/0127843965/ref=nosim/ericstreasuretro}{Handbook
of Differential Equations, 3rd ed.}~}Boston\emph{,} MA: Academic Press
1997, p.~120.

{[}49{]} Watson, G. N. \emph{A treatise on the theory of Bessel
functions}~(Vol. 3). The University Press. Cambridge, 1922., pp. 21-22,
74-75.

{[}50{]} Wallet, \emph{Properties of Bessel functions,} Libre Text
Mathematics, 2022. Web Resource.
\emph{\ul{https://math.libretexts.org/}}

{[}51{]} Weisstein E. W.~\emph{Beta Function},
From:~\href{https://mathworld.wolfram.com/}{\emph{MathWorld}}-\/-A
Wolfram Web
Resource.~\href{https://mathworld.wolfram.com/BetaFunction.html}{\emph{\ul{https://mathworld.wolfram.com/BetaFunction.html}}}

{[}52{]} Weisstein E. W.~\emph{Jacobi Polynomial},
From:~\href{https://mathworld.wolfram.com/}{\emph{MathWorld}}-\/-A
Wolfram Web
Resource.~\href{https://mathworld.wolfram.com/JacobiPolynomial.html}{\emph{\ul{https://mathworld.wolfram.com/JacobiPolynomial.html}}}

{[}53{]} Weisstein E. W.~\emph{Orthogonal Polynomials},
From:~\href{https://mathworld.wolfram.com/}{\emph{MathWorld}}-\/-A
Wolfram Web
Resource.~\href{https://mathworld.wolfram.com/OrthogonalPolynomials.html}{\emph{\ul{https://mathworld.wolfram.com/OrthogonalPolynomials.html}}}

{[}54{]} Wolfram Research, Inc., Mathematica, Version 14.1, Champaign,
IL (2024).

{[}55{]} GSL - GNU Scientific Library,
\emph{\href{https://www.gnu.org/software/gsl/}{\ul{https://www.gnu.org/software/gsl/}},}
24.01.2024.

\protect\phantomsection\label{_Toc54}{}Списак слика

Слика 2.3.1: График Беселових функција прве врсте за целобројни ред,
\(n = 0(1)4.\) 9

Слика 2.5.1: Графици Нојманових функција за различите вредности реда 7

Слика 3.1.1: График гама функције 5

Слика 4.1.1: Графички приказ бета функције 31

Слика 6.2.1. Првих пет Лежандрових полинома 43

Слика 7.2.1: Графици првих неколико Лагерових полинома 47

Слика 8.2.1. Ермитови физички полиноми (прва четири) 51

Слика 8.2.2. Ермитови пробабилистички полиноми (првих четири) 52

Слика 9.2.1: Графици првих пет Чебишевљевих полинома прве врсте (лево)
56

Слика 9.3.1: Првих пет Чебишевљевих полинома друге врсте 59

Слика 10.2.1: Првих неколико Јакобијевих полинома за различите вредности
параметарa 64

Слика 13.1.1: Почетна страна апликације 72

Слика 13.1.2: Изглед форме за унос података 73

Слика 13.1.3: Пример некорректног уноса 73

Слика 13.1.4: Унос поља помоћу калкулатора 74

Слика 13.1.5: Опште информације о функцији 75

Слика 13.1.6: Избор језика у апликацији 75

Слика 13.1.7: Остале странице у апликацији 76

Слика 12.3.1. Приказ резултата рада класе која апроксимирa вредност
Беселове функције 84

Слика 14.1: Калкулатор Беселових функција на алату \emph{WolframAlpha}
95

Слика 14.2: Калкулатор функција апликације \emph{Symbolab} 97

Слика 14.3: Рад апликације \emph{Special Functions Calculator} 97

\protect\phantomsection\label{_Toc55}{}Списак табела

Табела 3.1: Вредности гама функције за неке вредности аргумента 24

Табела 4.1: Вредности бета функције за неке парове аргумената 31

Табела 9.3.1: Првих пет Чебишевљевих полинома друге врсте 58

Табела 11.1: Преглед најважнијих особина Беселових функције прве врсте
66

Табела 11.2: Преглед најважнијих особина Лежандрових полинома 66

Табела 11.3: Преглед најважнијих особина Лагерових полинома 67

Табела 11.4: Преглед најважнијих особина Ермитовог физичког полинома 67

Табела 11.5: Преглед најважнијих особина Ермитовог пробабилистичког
полинома 68

Табела 11.6: Преглед најважнијих особина Чебишевљевих полинома прве
врсте 68

Табела 11.7: Преглед најважнијих особина Чебишевљевих полинома друге
врсте 69

Табела 11.8: Преглед најважнијих особина Јакобијевих полинома 69

Табела 13.2.1: Приоритет оператора коришћен у алгоритму за конверзију
инфиксног

израза у постфиксни 78

Табела 13.3.1. Парови параметара који дају најбоље апроксимације 86

\protect\phantomsection\label{_Toc56}{}БИОГРАФИЈА

Сара Лазић је рођена 28. 06. 1997. године у Ваљеву. Завршила је основну
школу „Владика Николај Велимировић`` у Ваљеву као носилац Вукове дипломе
и признања Ђак генерације, а Ваљевску гимназију, одељење ученика
обдарених за математику, као носилац Вукове дипломе. Током школовања
освојила је више првих награда на окружним, као и похвала на државним
такмичењима из математике, астрономије, хемије и биологије.
Електротехнички факултет у Београду уписала је 2016. године. Дипломирала
је на одсеку за Рачунарску технику и информатику 2021. године с
просечном оценом 8,11. Дипломски рад на тему „Имплементација мрежне
телеметрије'' одбранила је у септембру 2021. године са оценом 10.

Студије другог степена академских студија на Електротехничком факултету
у Београду, на Модулу за софтверско инжењерство, започела је у октобру
2021. године. Све испите положила је са просечном оценом 10, а пријавила
је тему мастер рада „Калкулатор неких специјалних функција''. Коаутор је
стручног рада „Јава апликација за нумеричку интеграцију'' презентованог
на једанаестом симпозијуму „Математика и примене'' новембра 2021.
године. Први аутор је рада „Калкулатор неких специјалних функција``
презентованог на Српском математичком конгресу (СМАК) 2024, као и
истоименог рада објављеног у часопису \emph{Serbian Journal of
Electrical Engineering}, Vol 21 (2024), No 3, pp. 345--357 {[}24{]}.

Од 2021. године запослена је у софтверској компанији која развија
платформу за анализу и коришћење професионалних мрежа контаката, где
тренутно обавља функцију техничког вође тима.
